%% MoCHII — gas physics, atomic data, algorithms, and validation.
%% Author: Kwang-Il Seon
%% Written 2026-07-12 for my_memo.cls.
\documentclass[english,a4paper]{my_memo}
\synctex=1
\usepackage{amsmath}
\usepackage{amssymb}
\usepackage{microtype}
\usepackage{booktabs}
\usepackage{longtable}

\newcommand{\code}[1]{\texttt{#1}}
\newcommand{\cmc}{\mathrm{cm^{3}\,s^{-1}}}
\newcommand{\xHI}{x_{\mathrm{HI}}}

\begin{document}

\title{MoCHII: Physics, Atomic Data, Algorithms, and Validation}
\author{Kwang-Il Seon}
\date{\docmoddate}
\maketitle

\begin{abstract}
MoCHII (MOnte Carlo for H\,\textsc{ii} regions) adds a photoionization gas
layer to the validated octree-AMR dust radiative-transfer engine of MoCafe
v2.00.  This memo documents the
ionizing frequency band and its zero-variance transport, the complete
atomic-data inventory with sources and the analytic fitting functions
evaluated at run time, the iteration and equilibrium algorithms
(ionization, thermal balance, trace-metal cascade, explicit diffuse
field), the $n$-level emissivity solver, the ported nebular
continuum and Storey \& Hummer recombination lines, and the
quantitative validation gates: rate
integrals vs.\ analytic attenuation (bias $+0.004\%$); the Str\"omgren
sphere vs.\ an exact 1D reference ($+0.56\%$ in ionized volume; AMR
$\equiv$ uniform); the thermal H/He nebula (median $0.14\%$ in $T_e$);
the MOCASSIN Lexington benchmarks ($T_e = 8169$~K at HII40 and 6925~K at
HII20 with the density-suppressed cooling default, matching a current
\textsc{Cloudy} c25.00 run to $0.50\%$ and $0.22\%$; the zero-density limit
gives 7894 and 6376~K, and both modern codes sit above the 2000-era
published ranges); and the line-flux comparison (H$\beta$ within a few
percent of the published values, the optical collisional lines within
the published spread).
Also documented: the solution-driven I-front re-refinement
(re-refined vs.\ native high-resolution octree $+0.008\%$ in ionized
volume with $4.5\times$ fewer cells) and an Illustris-TNG
post-processing demonstration; the completed dust coupling
(ionizing-band dust absorption with the dusty Str\"omgren reduction
0.724 reproduced to $+0.34\%$; EUV/FUV dust \emph{scattering} with
forced-first interaction sampling, validated by bracketing; live
ionization-dependent dust with a PAH-specific survival factor; and the
equilibrium dust temperature + IR spectrum consuming the transported
grain heating, energy closure 1.000); and the incremental metal
additions (argon, magnesium, iron, and silicon/chlorine/calcium) as
pure data operations.
\end{abstract}

\tableofcontents

%==========================================================================
\section{Scope and design rules}
MoCHII copies validated modules from MoCafe v2.00 (octree, event-driven
traversal, MPI shared memory, HDF5/FITS I/O) and adds gas microphysics
under three design rules:
(1) \emph{atomic data are fitted offline, evaluated online} --- the Python
pipeline under \code{tools/fitting/} is the auditable source of every
coefficient in \code{data/atomic/}, and Fortran only evaluates closed
forms; (2) \emph{adding an ion is a data operation}, realized by the
species registry (Section~\ref{sec:cascade}); (3) \emph{validation by
quantitative gates}, each capability compared against an analytic or
external reference.

The capabilities documented here: the ionizing band + rate integrals;
ionization equilibrium + opacity feedback (case B, fixed $T_e$);
thermal balance + trace-metal frame; the explicit diffuse field
(case A); solution-driven I-front re-refinement + snapshot
post-processing; the dust coupling (absorption competition,
scattering, live PAH split, grain heating/temperature, and the SEDust
stochastic-emission coupling); the incremental ions (argon,
magnesium, iron, silicon, chlorine, calcium); plus the output-time
diagnostic layer: $n$-level emissivities, Storey \& Hummer
recombination lines, and the nebular continuum.

%==========================================================================
\section{The ionizing band and transport}
\subsection{Continuous energy and diagnostic bins}
The ionizing band covers photon energies
$[E_{\min}, E_{\max}] = [13.598, 100]$~eV.  By default
(\code{par\%ion\_energy\_mode = 'continuous'}), a packet energy is sampled
from the normalized source-energy CDF, rather than from a frequency bin.
For a Planck source the CDF is built adaptively from the Planck energy
spectrum; for a tabulated source it is the cumulative integral of the
input $L_E$.  H\,\textsc{i}, He\,\textsc{i}, He\,\textsc{ii}, and active
metal thresholds are retained as knots in the adaptive Planck CDF.  With
$u\in[0,1)$, the draw is
\begin{equation}
 F_E(E) = \frac{\int_{E_{\min}}^E L_{E'}\,{\rm d}E'}
                {\int_{E_{\min}}^{E_{\max}} L_{E'}\,{\rm d}E'},
 \qquad E_{\rm pkt}=F_E^{-1}(u), \qquad
 L_{\rm pkt}=\frac{L_{\rm band}}{N_{\rm pkt}} .
 \label{eq:band-source-cdf}
\end{equation}
Here $L_{\rm band}$ is the luminosity integrated over the whole ionizing
band, not the luminosity of one energy interval.  Each packet is therefore
an equal-luminosity sample of the source spectrum.

The packet retains $E_{\rm pkt}$ throughout transport.  Gas opacity,
photoionization cross sections, and the excess energy used for heating are
evaluated at this exact energy.  The $n_\nu$ bins
(\code{par\%nnu\_ion}; gates use 32) receive only a diagnostic $J_E$ tally:
the sampled energy is located by a binary search of the bin-edge array and
assigned an output index.  Consequently the number or placement of those
bins cannot alter the continuous-mode physical solution; it changes only
the spectral resolution of saved tallies and rate diagnostics.  This is
important at thresholds: a photon exactly on a stored edge is assigned to
the bin above it, but its interaction physics still uses its own energy.

The validated continuous-energy path currently accepts one internal point
source with a Planck or tabulated spectrum.  It rejects unsupported
combinations (external sources, FUV extension, ionizing dust, secondary
ionization, metastable-He, and peel-off options) rather than falling back
silently to grouped energy physics.

\paragraph{Grouped mode.}  With \code{par\%ion\_energy\_mode = 'grouped'},
packets instead sample a luminosity bin and use its representative energy;
this is an explicit compatibility/reproduction mode, not the default
physical-energy treatment.
In this explicit compatibility mode, \code{par\%ion\_align\_edges} (default
on) pins the ionizing-band bin edges onto the ionization thresholds --- H\,\textsc{i}
(13.6~eV), He\,\textsc{i} (24.6~eV), He\,\textsc{ii} (54.4~eV), plus the
active metal photoionization thresholds gathered from the registry ---
so that no bin straddles a threshold.  The sub-segments between
consecutive thresholds are log-uniform, and bins are distributed among
the segments by log width (largest-remainder, floor 1).  When the
thresholds outnumber the bins the lowest-priority (metal) edges are
dropped gracefully --- the H/He edges are always kept --- so the option
is safe at any \code{nnu\_ion}.  Setting it \code{.false.} uses a
single log grid.  The
alignment removes a real discretization error: a coarse 32-bin log grid
straddles the 24.6~eV He\,\textsc{i} edge and delivers $\sim$18\% too
many He-ionizing photons against a true 40~kK blackbody
($Q_{\rm HeI}/Q_{\rm HI} = 0.127$ versus the exact 0.108), over-ionizing
He (and, downstream, the metals); aligning the edges restores the exact
ratio at fixed bin count (Section~\ref{sec:hii}).

\paragraph{FUV extension.}  With \code{par\%add\_fuv} the band grows a
second, lower segment: \code{par\%nnu\_fuv} log bins on
$[E_{\rm FUV}, E_{\min}]$ (\code{par\%efuv\_min}, default 6~eV)
\emph{below} the unchanged ionizing bins, so widening the band does not
coarsen the ionizing resolution (the earlier single-segment caveat).
\code{par\%luminosity} keeps its meaning as the ionizing
$[E_{\min}, E_{\max}]$ luminosity; the FUV part of the same source
spectrum rides on top ($L_{\rm FUV}/L_{\rm ion} = 1.097$ for a
4$\times$10$^4$~K Planck source at 6~eV), so the ionizing photon budget
$Q_{\rm H}$ is preserved without manual band-ratio rescaling and packets
carry $L_{\rm pkt} = (L_{\rm ion} + L_{\rm FUV})/N$.  H and He cross
sections vanish below their thresholds, so the FUV bins interact only
with dust and with the low-threshold metal stages (next subsection) ---
they penetrate the neutral gas beyond the I-front and photoionize
Mg\,\textsc{i} (7.65~eV), C\,\textsc{i} (11.26~eV), S\,\textsc{i}
(10.36~eV), and Fe\,\textsc{i} (7.90~eV) there, producing the
PDR-side Mg\,\textsc{ii} and C\,\textsc{ii} that the
ionizing-only band cannot.

\subsection{Source components: multiple point sources and an external
field}
\label{sec:sources}
The band is fed by any number of \emph{source components}:
\code{par\%nsource} internal point sources (0--16; positions
\code{par\%src\_x/y/z}, ionizing luminosities \code{par\%src\_lum} ---
all set, or none for an equal split of \code{par\%luminosity}, which
becomes the sum) plus, independently, one isotropic external radiation
field, on when \code{par\%ext\_intensity} $> 0$.  Each packet draws a
component in proportion to its band-total luminosity, is emitted from
it, and draws its frequency bin from that component's own spectrum CDF;
every packet carries $L_{\rm pkt} = \sum_c L_c/N$.  A single-component
run takes the legacy code path (the RNG stream is unchanged, so earlier
results reproduce).  The spectrum of internal source $i$ resolves as
\code{par\%src\_spectrum\_file} column $i$ $>$ \code{src\_tstar}$(i)$
Planck $>$ the global \code{ion\_spectrum} file $>$ the global
\code{tstar}; \code{src\_spectrum\_file} is one multi-column text file
--- column 1 the photon energy $E$ [eV] ascending, then one $L_E$
column for each source in arbitrary units, each independently
normalized so source $i$'s ionizing segment carries
\code{src\_lum}$(i)$.

The external field is an isotropic band-integrated mean intensity $J$
[erg\,s$^{-1}$\,cm$^{-2}$\,sr$^{-1}$] entering through an illuminating
surface (\code{par\%ext\_geometry}): \code{'rec'}, the box faces
(face-area-weighted --- correct for a non-cubic box), or \code{'sph'},
the bounding sphere of radius \code{par\%rmax} about the box center.
The inward flux through the surface is the hemisphere integral of
$J\cos\theta$, so the entering band power is
\begin{equation}
  F_{\rm in} = \int_{2\pi} J\cos\theta\,d\Omega = \pi J,
  \qquad
  L = \pi J A_{\rm surf}.
  \label{eq:extlum}
\end{equation}
Entry directions are cosine-weighted (Lambert) into the inward
hemisphere, $\mu = \sqrt{\xi}$, which reproduces an isotropic interior
field with energy density $u = 4\pi J/c$: an optically thin box tallies
an interior mean intensity equal to $J$ (the gate of
Section~\ref{sec:extval}).  Its spectrum resolves as
\code{ext\_spectrum} (file) $>$ \code{ext\_tstar} Planck $>$
the global spectrum.

\paragraph{Physical-unit spectra.}  Every file slot is read in the
column units set by \code{par\%spectrum\_type} (default \code{'shape'},
arbitrary units renormalized to the scale --- the legacy behavior).  The
physical types \code{'per\_ev'}, \code{'per\_hz'}, \code{'per\_ang'},
\code{'per\_um'} give the first column as $E$ [eV], $\nu$ [Hz], or
$\lambda$ [\AA{} or $\mu$m] and the value columns as a luminosity
density $L_X$ (internal sources) or an intensity density $J_X$ (external
field) in the matching spectral unit; wavelength tables are converted to
energy and sorted.  The \code{'shape'} value columns are read as the
per-eV density $L_E$ on the $E$~[eV] grid; a shape tabulated in
wavelength or frequency must use the matching \code{per\_*} type, whose
conversion applies the $\mathrm{d}\lambda/\mathrm{d}E$ (or
$\mathrm{d}\nu/\mathrm{d}E$) Jacobian, so a raw re-grid onto $E$ under
\code{'shape'} would be wrong by $E^{-2}$.  A physical spectrum is
\emph{absolute}: its bin
luminosities follow the tabulated integral directly.  If the scale
(\code{luminosity} / \code{src\_lum}$(i)$ / \code{ext\_intensity}) is
unset it is derived from the file and reported; if set, the spectrum is
rescaled to it and acts as a shape.  With \code{add\_fuv} the same table
fills the FUV bins with no separate normalization; a spectrum with no
luminosity inside the band aborts.

\paragraph{ISRF presets.}  \code{par\%ext\_spectrum} may instead name one
of three analytic interstellar radiation fields: \code{'draine'}
(Draine 1978), \code{'habing'} (the Draine shape at Habing 1968
normalization, $/1.71$), and \code{'mathis'} (Mathis, Mezger \&
Panagia 1983), the shape of the SEDust radiation-field module.  They
are FUV-only (5--13.6~eV, exactly zero in the ionizing bins), so they
require \code{add\_fuv}, and \code{par\%ext\_scale} multiplies them.  The
Draine field is the energy-weighted photon polynomial
$F_{\rm ph}(E) = 1.658{\times}10^{6}\,E^{2} - 2.152{\times}10^{5}\,E^{3}
+ 6.919{\times}10^{3}\,E^{4}$ with $J_E = e\,F_{\rm ph}$
($e = 1.602{\times}10^{-12}$~erg\,eV$^{-1}$; $F_{\rm ph} = E\,N(E)$ is
energy-weighted --- a common pitfall).  The presets carry the canonical
FUV energy densities $u_{\rm FUV} = 8.94$, $5.23$,
$6.01{\times}10^{-14}$~erg\,cm$^{-3}$ (Draine / Habing / Mathis;
$G_0 = 1.69$, $0.99$, $1.14$ in Habing units).  When \code{ext\_intensity}
is set the preset or physical file is rescaled to that band-integrated
$J$ (a note is printed; \code{ext\_scale} is then ignored).

\subsection{Opacity and the zero-variance edge walk}
Each leaf carries the band absorption coefficient per code length
\begin{equation}
  \kappa_b \;=\; n_{\rm H}\left[ \xHI\,\sigma_{\rm HI}(E_b)
   + y_{\rm He}\left( x_{\rm HeI}\sigma_{\rm HeI} +
     x_{\rm HeII}\sigma_{\rm HeII}\right)(E_b)\right] d_{\rm cm}
   \;(+\;\kappa_{\rm dust}),
\end{equation}
stored as the full $(n_\nu, n_{\rm leaf})$ array \code{kap\_ion} --- the
grey $s_{\rm ext}$ rescaling of the dust SED band cannot apply because
the gas opacity varies with both leaf and bin.
With \code{par\%ion\_metal\_abs} (default on when \code{use\_metals})
the registry metals join the absorption,
$\Delta\kappa_b = n_{\rm H}\sum_{\rm el} A_{\rm el} \sum_i
f_i\,\sigma_{{\rm VFKY96},i}(E_b)\,d_{\rm cm}$, with the stage
fractions $f_i$ from the same product chain as the cooling, iterated by
the opacity feedback exactly like $\xHI$.  Above 13.6~eV this is a
$\sim$10$^{-22}$~cm$^2$ per H perturbation next to H/He and dust; in the
FUV bins it is the only \emph{gas} opacity.  The absorbed share returns
to the metals through their own $\Gamma$ integrals, so the bookkeeping
stays consistent (metal photoheating is not added to the thermal
balance --- negligible at trace abundances).

Where the band has no scattering and the diffuse field is emitted
explicitly, every packet's contribution to the mean-intensity tally is
computed \emph{analytically} along its ray (the forced-first-flight
estimator of MoCafe, made exact):
\begin{equation}
  \Delta J_b({\rm leaf}) \;=\; w\,L_{\rm pkt}\,
  \frac{e^{-\tau_{\rm in}} - e^{-\tau_{\rm out}}}{\kappa_b},
  \label{eq:edgewalk}
\end{equation}
accumulated leaf by leaf until the domain edge
(\code{raytrace\_ion\_to\_edge\_amr}).  No interaction sampling is
needed; the only Monte Carlo variance is in direction and bin sampling.
The rate-integral gate verifies this estimator to a mean bias of
$+0.004\%$.

\subsection{Rate integrals}
With $V$ the leaf volume and $J_b$ the reduced tally, the code forms for
each absorber $i$
\begin{equation}
  \Gamma_i = \sum_b \frac{[4\pi J\,\Delta\nu]_b\,\sigma_i(E_b)}{h\nu_b},
  \qquad
  \mathcal{H}_i = \sum_b [4\pi J\,\Delta\nu]_b\,\sigma_i(E_b)
  \left(1 - \frac{E_{{\rm th},i}}{E_b}\right),
  \label{eq:rates}
\end{equation}
the photoionization rate [s$^{-1}$] and photoheating [erg\,s$^{-1}$]
for each particle of species $i$.

%==========================================================================
\section{Atomic data}
\label{sec:atomic}
Every coefficient file in \code{data/atomic/} carries a provenance
header (source, version, fit range, maximum fit error).  The generators
live in \code{tools/fitting/}.

\subsection{Photoionization cross sections}
The one-electron ions H\,\textsc{i} ($Z{=}1$) and He\,\textsc{ii}
($Z{=}2$) use the \emph{exact} ground-state hydrogenic cross section
(Osterbrock \& Ferland 2006, Eq.~2.4; \code{sigma\_hydrogenic}, the same
form as EXHALE):
\begin{equation}
  \varepsilon = \sqrt{E/E_{\rm th}-1},\qquad
  \sigma(E) = \frac{\sigma_0}{Z^2}\left(\frac{E_{\rm th}}{E}\right)^{4}
  \frac{\exp\!\big(4 - 4\arctan\varepsilon/\varepsilon\big)}
       {1 - e^{-2\pi/\varepsilon}},
  \label{eq:sighyd}
\end{equation}
with $\sigma_0 = 6.30{\times}10^{-18}\,\mathrm{cm^2}$ the threshold value
and $E_{\rm th}$ the true ionization potential (13.598 and 54.416~eV),
tied to the physical edge (not $13.6\,Z^2$) so it stays aligned with the
band grid.  He\,\textsc{i} (two electrons) and every metal ion use the
full VFKY96 single-shell analytic form (Verner, Ferland, Korista \&
Yakovlev 1996, ApJ 465, 487, Eq.~1):
\begin{equation}
  x = \frac{E}{E_0} - y_0,\quad z = \sqrt{x^2 + y_1^2},\quad
  Q = 5.5 - \tfrac{1}{2}P,\quad
  \sigma(E) = \sigma_0\left[(x-1)^2 + y_w^2\right]
  z^{-Q}\left(1 + \sqrt{z/y_a}\right)^{-P},
  \label{eq:vfky}
\end{equation}
zero below threshold, $\sigma_0$ in Mb.  Parameters were parsed directly
from the authors' \code{phfit2.f} \code{PH2} (outer-shell) table and
verified against the EXHALE implementation; H/He values are compiled in
\code{photo\_xsec.f90}, the registry ions in
\code{element\_<el>.txt}.  For the elements \emph{absent} from the
VFKY96 outer-shell table (P, Cl, K --- Opacity-Project gaps; the
MOCASSIN \code{ph2.dat} read loop skips the same $Z$), the
Verner \& Yakovlev (1995) outer-shell subshell fit is used instead
(\code{sigma\_vy95}; \code{PHOTO2} lines):
$\sigma = \sigma_0[(y-1)^2 + y_w^2]\,y^{-(5.5+l-P/2)}
(1+\sqrt{y/y_a})^{-P}$ with $y = E/E_0$ and $l$ the subshell orbital
quantum number --- the same VFKY96/VY95 hybrid as MOCASSIN
\code{phFitEl}.  Chlorine is the first user of this path.  Among H/He,
only the He\,\textsc{i} ground state uses VFKY96:

\begin{center}\small
\begin{tabular}{lrrrrrrrr}
\toprule
ion & $E_{\rm th}$[eV] & $E_0$ & $\sigma_0$[Mb] & $y_a$ & $P$ & $y_w$ & $y_0$ & $y_1$ \\
\midrule
He\,\textsc{i}  & 24.587 & 13.61  & 949.2 & 1.469 & 3.188 & 2.039 & 0.4434 & 2.136 \\
\bottomrule
\end{tabular}
\end{center}
plus C\,\textsc{i}--\textsc{iii}, N\,\textsc{i}--\textsc{ii},
O\,\textsc{i}--\textsc{iii}, Ne\,\textsc{i}--\textsc{iii},
S\,\textsc{i}--\textsc{iv}, Ar\,\textsc{i}--\textsc{iv} (20 rows, in
the element files).  The rate-integral gate cross-checks the curves
against EXHALE's independent implementations.  The exact H\,\textsc{i}
and He\,\textsc{ii} cross sections give the threshold values
$\sigma_0/Z^2 = 6.30$ and $1.575$~Mb; the VFKY96 fits they replace for
H/He agreed with the exact curve to $0.73\%$ (H\,\textsc{i}) and
$0.43\%$ (He\,\textsc{ii}), and $\sigma_{\rm HeI} = 7.43$~Mb (VFKY96).

\subsection{H/He recombination}
\label{sec:hherec}
The recombination coefficients in the ionization balance are chosen by
\code{par\%recomb\_model}.  The default \code{'badnell\_mao'} takes the
total (case A) radiative recombination $\alpha_A$ from the same Badnell
(2023) RR fit used for the metals [Eq.~(\ref{eq:badnellrr})], obtains the
ground-level ($n{=}1$) direct-capture rate $\alpha_1$ from the Milne
relation applied to the \emph{same} photoionization cross sections the
transport absorbs with [Eq.~(\ref{eq:milne-alpha1})], and forms case B by
removing the ground term, $\alpha_B = \alpha_A - \alpha_1$.  At run time
$\alpha_1$ is evaluated from the seven-coefficient functional form of
Mao \& Kaastra (2016, A\&A 587, A84) ($T_{\rm eV} = k_B T$; $a_0$ in
$10^{-10}\,\cmc$),
\begin{equation}
  \alpha_1(T) = a_0\,10^{-10}\,
  T_{\rm eV}^{-b_0 - c_0\ln T_{\rm eV}}\,
  \frac{1 + a_2\,T_{\rm eV}^{-b_2}}{1 + a_1\,T_{\rm eV}^{-b_1}},
  \label{eq:mao}
\end{equation}
but the coefficients are fitted to that Milne integral instead of being
taken from Mao \& Kaastra; \S\ref{sec:diffuse} lists them with their fit
residuals and gives the measurement that forced the replacement.  The
namelist value \code{'badnell\_mao'} and the routine name \code{rr\_mao}
in \code{recomb\_mod} therefore record the functional form and the
history, not the origin of the numbers.  Because $\alpha_1$ and the
ground-continuum emissivity of \S\ref{sec:diffuse} are now one integral
over one cross section, the case A/B split and the diffuse packet count
cannot drift apart.  He\,\textsc{ii}$\to$He\,\textsc{i} adds
the Badnell three-term dielectronic sum to $\alpha_A$ (negligible below
$\sim5{\times}10^4$~K); H\,\textsc{ii} and He\,\textsc{iii} are bare
nuclei with no DR.  The coefficients at $10^4$~K
($10^{-13}\,\cmc$), with the independently level-resolved Cloudy c25.00
values in parentheses, are
\begin{center}\small
\begin{tabular}{lrrr}
\toprule
reaction & $\alpha_A$ & $\alpha_1$ & $\alpha_B$ \\
\midrule
H\,\textsc{ii}$\to$H\,\textsc{i}     & 4.19 (4.16) & 1.579 (1.584) & 2.61 \\
He\,\textsc{ii}$\to$He\,\textsc{i}   & 4.37 (4.36) & 1.609 (1.603) & 2.76 \\
He\,\textsc{iii}$\to$He\,\textsc{ii} & 21.9 (21.8) & 6.512 (6.515) & 15.4 \\
\bottomrule
\end{tabular}
\end{center}
Setting \code{par\%recomb\_model = 'hui\_gnedin'} restores the earlier
Hui \& Gnedin (1997, MNRAS 292, 27) case-A/B fits with
$\lambda = 2T_{\rm TR}/T$ ($T_{\rm TR} = 157807$, $285335$, $631515$~K
for H\,\textsc{i}, He\,\textsc{i}, He\,\textsc{ii}):
\begin{align}
  \alpha_A^{\rm HII} &= 1.269{\times}10^{-13}\,
    \lambda^{1.503}\big/\big[1+(\lambda/0.522)^{0.470}\big]^{1.923}, &
  \alpha_B^{\rm HII} &= 2.753{\times}10^{-14}\,
    \lambda^{1.500}\big/\big[1+(\lambda/2.740)^{0.407}\big]^{2.242},
    \nonumber\\
  \alpha_A^{\rm HeII} &= 3.000{\times}10^{-14}\,\lambda^{0.654}, &
  \alpha_B^{\rm HeII} &= 1.260{\times}10^{-14}\,\lambda^{0.750},
\end{align}
He\,\textsc{iii} hydrogenic ($Z{=}2$).  These are retained so the
recorded gates reproduce; they agree with the default to $<0.6\%$ for
H\,\textsc{ii} and He\,\textsc{iii}, but give
He\,\textsc{ii}$\to$He\,\textsc{i} $5\%$ lower ($\alpha_B = 2.62$ vs
$2.76{\times}10^{-13}\,\cmc$).
Recombination \emph{cooling} uses the companion fits
\begin{equation}
  \beta_B^{\rm HII} = 3.435{\times}10^{-30}\,T\,\lambda^{1.970}
  \big/\big[1+(\lambda/2.250)^{0.376}\big]^{3.720}
  \;\mathrm{erg\,cm^3\,s^{-1}},
\end{equation}
(case A analog with $1.778{\times}10^{-29}$, $0.541$, $1.965$, $0.502$,
$2.697$), He\,\textsc{iii} scaled as for $\alpha$, and He\,\textsc{ii} as
$k_B T\,\alpha_{A,B}^{\rm HeII}$ (their recommendation; He\,\textsc{ii}
dielectronic recombination is negligible at nebular temperatures).

\subsection{H/He microphysics: comparison with other codes}
\label{sec:codecmp}
The H/He cross sections and recombination can be placed next to three
reference codes.  For the photoionization cross section, MoCHII and
EXHALE evaluate the exact hydrogenic formula [Eq.~(\ref{eq:sighyd})] for
H\,\textsc{i}/He\,\textsc{ii}; MOCASSIN fits them with VFKY96 (agreeing
with the exact curve to $<1\%$), and Cloudy resolves them from its
iso-electronic H-like/He-like model atoms.  For the recombination that
closes the ionization balance (Table~\ref{tab:codecmp}), MoCHII's default
takes $\alpha_A$ from the modern Strathclyde (Badnell) data and splits off
$\alpha_1$ with the Milne relation, which is how Cloudy resolves its
H-like and He-like recombination level by level; MOCASSIN uses the
older Verner \& Ferland (1996) total; EXHALE, a planetary-atmosphere
code, adopted the Badnell RR $+$ Mao \& Kaastra $\alpha_1$ construction as
its default on 2026-07-17, before the $\alpha_1$ correction of
\S\ref{sec:diffuse} (Hui \& Gnedin remains a legacy option).  MOCASSIN
and Cloudy always work in
case A and transport the diffuse ground-recombination field explicitly,
whereas MoCHII defaults to case B (on-the-spot) with an optional explicit
diffuse field, and EXHALE is case B only.

\begin{table}[htbp]
\centering\small
\caption{H/He recombination in the ionization balance: source and
$\alpha$ at $10^4$~K ($10^{-13}\,\cmc$; case A / case B where both
apply).}
\label{tab:codecmp}
\begin{tabular}{llccc}
\toprule
code & source & H\,\textsc{ii}$\to$\textsc{i}
     & He\,\textsc{ii}$\to$\textsc{i} & He\,\textsc{iii}$\to$\textsc{ii} \\
\midrule
MoCHII (default)          & Badnell RR $+$ Milne $\alpha_1$            & 4.19 / 2.61 & 4.37 / 2.76 & 21.9 / 15.4 \\
MoCHII (\code{hui\_gnedin})& Hui \& Gnedin 1997                        & 4.30 / 2.59 & 4.23 / 2.62 & 22.3 / 15.5 \\
EXHALE (default)           & Badnell RR $+$ Mao\,\&\,Kaastra $\alpha_1$ & -- / 2.61   & -- / 2.81   & -- / 15.4   \\
EXHALE (\code{legacy\_hhe\_rates}) & Hui \& Gnedin 1997 (case B)         & -- / 2.59   & -- / 2.62   & -- / 15.5   \\
MOCASSIN                   & Verner \& Ferland 1996 (case A)           & 4.19        & 4.84        & 21.9        \\
Cloudy c25.00              & level-resolved Milne, own $\sigma_{\rm pi}$ (case A) & 4.16 & 4.36 & 21.8 \\
\bottomrule
\end{tabular}
\end{table}
The H\,\textsc{ii} and He\,\textsc{iii} rates agree across all four codes
to within a few percent.  The one substantive difference is
He\,\textsc{ii}$\to$He\,\textsc{i}: MoCHII's Badnell total ($4.37$) and
Cloudy's explicit sum over levels ($4.36$) agree to $0.3\%$, MOCASSIN's
Verner \& Ferland value ($4.84$) is the high outlier, and the earlier
Hui \& Gnedin ($4.23$) was low.  Moving MoCHII onto the modern Badnell
value therefore raises its case-B He recombination by $\sim5\%$, in the
direction that reduces the He over-ionization seen relative to converged
MOCASSIN in the HII40 benchmark (\S\ref{sec:validation} discusses that
residual).  Because that $11\%$ spread in a single coefficient is the
largest code-to-code difference anywhere in the H/He data, \S\ref{sec:henetwork}
takes it apart.

\subsection{The helium ionization network, code by code}
\label{sec:henetwork}
Neutral helium is the one species in the H/He inventory on which the
photoionization codes disagree at more than the percent level.  Reading
each code's own data rather than its published description gives, for
$\alpha_A$(He\,\textsc{i}) at $10^4$~K,
\begin{center}\small
\begin{tabular}{@{}llrr@{}}
\toprule
code & how $\alpha_A$(He\,\textsc{i}) is obtained
     & $10^{-13}\,\cmc$ & vs.\ Cloudy \\
\midrule
Cloudy c25.00 & level-by-level Milne from its own $\sigma_{\rm pi}$,
                total cached & 4.3583 & --- \\
MoCHII        & Badnell RR total $+$ DR & 4.3719 & $+0.3\%$ \\
Wood \code{ionize} / CMacIonize
              & sum of the four BSS99 \emph{effective} coefficients
              & 4.2630 & $-2.2\%$ \\
MOCASSIN      & Verner \& Ferland 1996 total fit & 4.8394 & $+11.0\%$ \\
\bottomrule
\end{tabular}
\end{center}
The MOCASSIN entry is decoded from \code{data/radrec.dat}, replicating the
read order of \code{update\_mod.f90} lines 1707--1735; the decode is
validated on the two channels whose answer is known independently, giving
H\,\textsc{i} $4.1923$ against Cloudy's $4.1642{\times}10^{-13}$ ($+0.7\%$)
and He\,\textsc{ii} $2.1880$ against $2.1817{\times}10^{-12}$ ($+0.3\%$).

\paragraph{The spread is not the extent of the cascade sum.}  A natural
first reading of an $11\%$ spread is that the codes stop the recombination
cascade at different levels.  None of them does.  Three treatments are in
play and all three are complete: Cloudy sums explicitly over resolved
levels; MoCHII and MOCASSIN use totals over all $n$ (Badnell, VF96); Wood
and CMacIonize use \emph{effective} coefficients, in which the cascades
from higher levels are already folded into the $n = 2$ terms.  MoCHII's
Badnell total agrees with Cloudy's explicit level sum to $0.3\%$, which is
what a truncation on either side would break.  (The 20-level list in
Cloudy's \code{source/sanity\_check.cpp} is a unit test of threshold cross
sections, not the extent of its recombination sum.)

The control that settles it is the electron count:
\begin{center}\small
\begin{tabular}{@{}lcc@{}}
\toprule
species & electrons & code-to-code spread in $\alpha_A$ \\
\midrule
H\,\textsc{i}   & 1 (hydrogenic) & $0.7\%$ \\
He\,\textsc{ii} & 1 (hydrogenic) & $0.3\%$ \\
He\,\textsc{i}  & \textbf{2}     & $\mathbf{11}$--$\mathbf{13\%}$ \\
\bottomrule
\end{tabular}
\end{center}
The one-electron species have an exact analytic solution, so there is
nothing for two calculations to disagree about, and they do not; the
two-electron one has no closed form and carries screening and the
singlet/triplet structure, and it spreads by $11$--$13\%$.  The disagreement
is therefore in the atomic data for a two-electron atom, not in how far
anybody sums.

\paragraph{Three of the four cluster.}  Within He\,\textsc{i} the values are
not scattered but split three-to-one: Cloudy's level sum $4.3583$, Badnell
$4.3719$ and BSS99 $4.2630{\times}10^{-13}\,\cmc$ all lie within $2.5\%$ of
each other, and Verner \& Ferland (1996) alone sits at
$4.8394{\times}10^{-13}$.  MOCASSIN applies its \code{radRecFit} uniformly
with no special case for H or He
(\code{source/update\_mod.f90:1636--1639}, whose own comment names
``H-like, He-like, Li-like, Na-like --- Verner \& Ferland, 1996, ApJS, 103,
467''), and adds dielectronic recombination only for
$Z_{\rm elem} \ge 3$, so its He\,\textsc{i} rate is the VF96 radiative
value alone.  Offered as reasoning rather than as a verified claim: VF96's
He-like fits lean on hydrogenic scaling for the excited levels, and $Z = 2$
is where that approximation is worst, since neutral helium is exactly where
screening and the singlet/triplet splitting matter most.

\paragraph{A caution for the MOCASSIN regression suite.}  MOCASSIN's
11-case suite is used here as a three-dimensional validation reference.
Every helium comparison against it inherits a known $+11\%$ in
$\alpha_A$(He\,\textsc{i}), and that difference has been traced to VF96
by the first-principles Cloudy determination above, so it should be
subtracted before a disagreement in a helium ionization front or a helium
line intensity is attributed to transport or geometry.  The same
comparison across the metal registry, where no such independent
determination is available, is in \S\ref{sec:recaudit}.

\paragraph{Excited He\,\textsc{i} levels are not a hidden channel.}  What
Cloudy's resolved He\,\textsc{i} model atom has and an effective treatment
does not is a population in the $n = 2$ terms that can itself be ionized.
MoCHII cannot photoionize those levels at all: the band floor is
\code{par\%eion\_min} $= 13.598$~eV while every excited He\,\textsc{i}
level has an ionization potential below it ($2^3$S 4.767, $2^1$S 3.971,
$2^3$P 3.623, $2^1$P 3.369~eV).  The size of what is missed is bounded by
the metastable population, since $2^1$S decays by two-photon emission at
$51.3\,{\rm s^{-1}}$ and is left $10^7$ times less populated than $2^3$S:
with $n(2^3{\rm S})/n({\rm He^0}) = 6.2{\times}10^{-5}$ at $r = 3$~pc,
$Q({>}4.767~{\rm eV})/Q({>}24.587~{\rm eV}) = 33.8$ for a 40-kK blackbody
and Cloudy's own $2^3$S threshold cross section of
$5.48{\times}10^{-18}\,{\rm cm^2}$, the extra ionization is
$0.2$--$0.6\%$, and collisional ionization out of $2^3$S adds a comparable
${\sim}0.6\%$.

Taken together with the front sensitivity measured in
\S\ref{sec:diffuse} ($0.80$~pc per unit fractional change in the
He\,\textsc{i} recombination rate), the $11$--$13\%$ code spread is worth
$0.09$~pc --- half the HII40 front residual against Cloudy --- but MoCHII
itself lies within $0.3\%$ of Cloudy, so that lever is not available to
MoCHII: it measures how much of a helium front comparison between two
arbitrary codes can be recombination data alone.

\subsection{Collisional ionization (Voronov 1997)}
\begin{equation}
  k(T) = A\,\frac{1 + P\sqrt{U}}{X + U}\,U^{K} e^{-U},
  \qquad U = \frac{\Delta E}{kT},
\end{equation}
with parameters $(\Delta E\,[\mathrm{eV}], P, A, X, K)$ parsed from the
author's \code{cfit.f} table:
\begin{center}\small
\begin{tabular}{lrrrrr}
\toprule
ion & $\Delta E$ & $P$ & $A\,[\cmc]$ & $X$ & $K$ \\
\midrule
H\,\textsc{i}   & 13.6 & 0 & $2.91{\times}10^{-8}$ & 0.232 & 0.39 \\
He\,\textsc{i}  & 24.6 & 0 & $1.75{\times}10^{-8}$ & 0.180 & 0.35 \\
He\,\textsc{ii} & 54.4 & 1 & $2.05{\times}10^{-9}$ & 0.265 & 0.25 \\
\bottomrule
\end{tabular}
\end{center}
plus the 16 registry rows in the element files.  An optional
\code{par\%ci\_model = 'dere\_hybrid'} multiplies each Voronov rate by
the constant Dere (2007)/Voronov ratio per (element, stage) taken from
Cloudy c23.01's \code{DereRatio} table; the default \code{'voronov'}
leaves the rate unchanged.  In photoionized gas at $\sim$$10^4$~K the
metal collisional ionization is Boltzmann-suppressed far below
photoionization, so the hybrid shifts the ionization structure only in
shielded, low-ionization zones where collisions are the sole
metal-ionization channel.

\subsection{Metal recombination: Badnell RR + DR}
For the registry cascade, total recombination of the recombining ion is
Badnell (2006, ApJS 167, 334) radiative recombination
\begin{equation}
  \alpha_{\rm RR}(T) = A\Big[\sqrt{T/T_0}\,
  \big(1+\sqrt{T/T_0}\big)^{1-b}\big(1+\sqrt{T/T_1}\big)^{1+b}\Big]^{-1},
  \qquad b = B + C e^{-T_2/T},
  \label{eq:badnellrr}
\end{equation}
plus the Badnell-group dielectronic fits
$\alpha_{\rm DR}(T) = T^{-3/2}\sum_i c_i e^{-E_i/T}$.
Coefficients are read from the Badnell/Strathclyde AMDPP TAMOC tables
(the \code{clist\_K} files, 2023 release, copied into
\code{data/atomic/badnell\_\{rr,dr\}.dat} and verified md5-identical to
the copies Cloudy c23.01 redistributes), using the $M{=}1$ ground-level
fit of each recombining ion.  Most ions are numerically identical to the
CHIANTI v11.0.2 \code{.rrparams}/\code{.drparams} (which \emph{are}
Badnell data), the notable exception being Si\,\textsc{ii}:
CHIANTI~11 still carries the Abdel-Naby et al.\ (2012) DR, whose
$\sim$$10^4$-K resonance is stronger, so the 2023 refit lowers the
total $\mathrm{Si}^+ + e \to \mathrm{Si}^0$ recombination to $0.48\times$
the CHIANTI value at $10^4$~K.  Ions Badnell does not tabulate (the low
Ar/Fe/Cl/Ca stages, confirmed absent at every $M$) are taken from the
CHIANTI files instead.
Where CHIANTI carries no Badnell DR (Ar\,\textsc{ii}, Ar\,\textsc{iii},
Cl\,\textsc{ii}, Ca\,\textsc{ii}, Ca\,\textsc{iv}, Ca\,\textsc{v}), the
\code{.drparams} type-2 entries are the Shull \&
Van Steenberg (1982) form
$\alpha_{\rm DR} = A\,T^{-3/2} e^{-T_0/T}(1 + B\,e^{-T_1/T})$, written
as \code{DR2} lines and supported by the same reader (negligible at
nebular temperatures; kept for completeness).  One further CHIANTI RR
type appears with iron and the higher calcium stages: the
Fe\,\textsc{ii}--\textsc{iv}, Ca\,\textsc{iv} and Ca\,\textsc{v}
\code{.rrparams} are type-3 power laws, written as \code{RR2} lines and
evaluated as $\alpha_{\rm RR} = A\,(T/10^4\,{\rm K})^{-\eta}$.
Section~\ref{sec:recaudit} audits all of this entry by entry.

\subsection{Recombination data across the registry: an audit}
\label{sec:recaudit}
The $\alpha_1$ source-mixing defect of \S\ref{sec:diffuse} was found in the
H/He channels, and it prompted a sweep of every recombination coefficient the
code carries --- 32 registry entries over 11 elements, plus the hardcoded H/He
channels.  Three questions were put to them: is each coefficient the one its
provenance header claims, is the functional form the best available for that
ion, and how far does the answer depend on which data set is used.

\paragraph{Transcription: no errors.}  Of the 32 registry entries, 22 take
their coefficients from the Badnell tables, and all 22 reproduce
\code{data/atomic/badnell\_\{rr,dr\}.dat} at $10^4$~K to a ratio of $1.0000$.
The key is fixed by the registry convention: \code{TRANSITION} $i$ of element
$Z$ is the recombination of the ion of charge $i$, so the Badnell row is
$(Z,\ N = Z - i,\ M = 1)$.  MoCHII's hardcoded H\,\textsc{i} coefficients
($8.318{\times}10^{-11}$, $0.7472$, $2.965$, $7.001{\times}10^{5}$) are exactly
the $Z = 1$, $N = 0$ row of the same table.  There is no transcription error
anywhere in the registry, and no coefficient traceable to the wrong ion.

\paragraph{The older forms cannot be migrated, and should not be.}  Badnell's
tabulation covers $N \le 15$ (through the Al-like sequence) plus $N = 18$
(Ar-like) for $Z \ge 20$.  Every use of \code{RR2} or \code{DR2} in the
registry falls at a $(Z, N)$ outside that coverage.  Where Badnell has nothing,
\code{tools/fitting/make\_element\_data.py} takes CHIANTI v11.0.2 and writes
\emph{exactly the functional form CHIANTI declares}: \code{rrparams} types 1
and 2 are the Badnell form and become \code{RR}, type 3 is the power law and
becomes \code{RR2}; \code{drparams} type 1 becomes \code{DR}, and type 2, the
Shull \& Van Steenberg (1982) form, becomes \code{DR2}.  Checked against the
CHIANTI file of the recombining ion, entry by entry:
\begin{center}\small
\begin{tabular}{@{}llccc@{}}
\toprule
registry entry & recombining ion & $N$ & Badnell & CHIANTI
  \code{rr}/\code{dr} type $\to$ form \\
\midrule
Ar 1        & Ar\,\textsc{ii}    & 17 & absent  & 2 / 2 $\to$ \code{RR} $+$ \code{DR2} \\
Ar 2        & Ar\,\textsc{iii}   & 16 & absent  & 1 / 2 $\to$ \code{RR} $+$ \code{DR2} \\
Ar 3        & Ar\,\textsc{iv}    & 15 & covered & 2 / 1 $\to$ \code{RR} $+$ \code{DR}  \\
Cl 1        & Cl\,\textsc{ii}    & 16 & absent  & 1 / 2 $\to$ \code{RR} $+$ \code{DR2} \\
Ca 1        & Ca\,\textsc{ii}    & 19 & absent  & 1 / 2 $\to$ \code{RR} $+$ \code{DR2} \\
Ca 2        & Ca\,\textsc{iii}   & 18 & covered & 2 / 1 $\to$ \code{RR} $+$ \code{DR}  \\
Ca 3        & Ca\,\textsc{iv}    & 17 & absent  & 3 / 2 $\to$ \code{RR2} $+$ \code{DR2} \\
Ca 4        & Ca\,\textsc{v}     & 16 & absent  & 3 / 2 $\to$ \code{RR2} $+$ \code{DR2} \\
Fe 1, 2, 3  & Fe\,\textsc{ii}--\textsc{iv} & 25, 24, 23 & absent
  & 3 / 1 $\to$ \code{RR2} $+$ \code{DR} \\
\bottomrule
\end{tabular}
\end{center}
Ar\,\textsc{iv} and Ca\,\textsc{iii} are the controls: Badnell covers $N = 15$,
and $N = 18$ for $Z \ge 20$, and the registry duly uses \code{RR} $+$
\code{DR} at both.  So the Shull \& Van Steenberg form is not a stale choice
made in preference to Badnell --- it is what CHIANTI itself still ships for
those ions.  There is nothing to migrate, and the registry already uses the
best source available at every element and stage.  A suggestion recorded
earlier, to move the \code{RR2} and \code{DR2} entries onto Badnell, is
therefore withdrawn on both counts: Badnell has no row for them, and CHIANTI's
current data for them is in those forms.

\paragraph{Code to code: variance between data sets, not an error bar on
MoCHII.}  What transcription cannot settle is which data set is right.
Comparing radiative recombination only --- dielectronic recombination is
excluded on both sides, because MOCASSIN adds it separately and keeping it on
one side alone would compare unlike quantities --- MoCHII's Badnell/CHIANTI
values against the Verner \& Ferland (1996) fits MOCASSIN uses, at $10^4$~K:
\begin{center}\small
\begin{tabular}{@{}lr@{}}
\toprule
recombination forming & MoCHII / VF96 \\
\midrule
C\,\textsc{ii}, C\,\textsc{iii}, N\,\textsc{ii}, Ne\,\textsc{ii}
  & $0.999$--$1.016$ \\
N\,\textsc{i}                   & 0.971 \\
Si\,\textsc{iv} (Na-like)       & 0.932 \\
Ne\,\textsc{i}                  & 0.898 \\
C\,\textsc{i}                   & 0.861 \\
O\,\textsc{i}                   & 0.837 \\
Mg\,\textsc{ii} (Na-like)       & 0.791 \\
\bottomrule
\end{tabular}
\end{center}
Eleven ions: seven within $10\%$, ten within $20\%$, median $0.971$.  The
spread is concentrated in the rates that form \emph{neutral atoms}, which is
where He\,\textsc{i} sits as well, at $11\%$ (\S\ref{sec:henetwork}).

What that table is has to be stated precisely.  It is the variance between two
data sets, \emph{not} a measured error in MoCHII, and for the metals there is
no third arbiter available here: Cloudy ships \code{data/badnell\_rr.dat} and
\code{badnell\_dr.dat} md5-identical to MoCHII's copies and evaluates them with
\code{RR\_Badnell\_rate\_coef} and \code{DR\_Badnell\_rate\_coef} in
\code{source/ion\_recomb.cpp}, so Cloudy and MoCHII agree on metal
recombination by construction and Cloudy cannot adjudicate.  With only two
determinations in play, which one is right for C\,\textsc{i}, O\,\textsc{i},
Ne\,\textsc{i} or Mg\,\textsc{ii} is undetermined.

He\,\textsc{i} is the one exception, and there MoCHII is demonstrably right.
Cloudy does \emph{not} use Badnell for the He-like and H-like sequences; it
derives them level by level from its own photoionization cross sections through
the Milne relation, which is an independent first-principles determination.
Against it, Badnell (MoCHII) is $+0.3\%$ and VF96 (MOCASSIN) is $+11.0\%$
(\S\ref{sec:henetwork}), so for He\,\textsc{i} the spread is attributable to
VF96.  One may reason by analogy that the metal neutrals behave the same way
--- Badnell (2006) is the later and more careful calculation, VF96 (1996) leans
on simplified treatments of many-electron systems, and He\,\textsc{i} is the
case where that was actually demonstrated --- but that is inference, not
measurement.  Settling it needs an independent determination for each ion:
R-matrix values for C\,\textsc{i}, O\,\textsc{i} and Ne\,\textsc{i}
recombination, or CHIANTI totals that do not derive from Badnell.

An earlier phrasing, ``a $10$--$20\%$ uncertainty in MoCHII's
neutral-forming recombination rates'', is withdrawn as imprecise: it reads a
spread between two data sets as an error bar on one of them.  The two accurate
statements are that MoCHII's He\,\textsc{i} rate agrees with an independent
first-principles calculation to $0.3\%$ while the VF96 set MOCASSIN uses is
$11\%$ high (measured), and that for the metal neutrals Badnell and VF96 differ
by $10$--$16\%$ with no third source in this configuration able to say which is
right (a code-to-code spread).

\paragraph{What the audit establishes, and what it does not.}  It found no
defect in MoCHII's recombination data: no transcription error, no form that a
better source could replace, and no coefficient attached to the wrong ion.
That is not the same statement as ``the data are accurate''.  For H and He an
independent determination exists and MoCHII matches it to $0.3$--$0.7\%$; for
the metals no arbiter is available in this configuration, so the verdict there
is the absence of a defect found, not accuracy established.

\subsection{Charge exchange with hydrogen}
\label{sec:cx}
First-stage reactions use the verified transcription of Huang et al.\
(2023, ApJ 951, 123) Table~4 (sources Stancil et al.\ 1998/1999, Lin et
al.\ 2005, Zhao et al.\ 2005, Kingdon \& Ferland 1996), e.g.\ for oxygen
\begin{align}
  k({\rm O} + {\rm H}^+) &= 2.08{\times}10^{-9}\,T_4^{0.405}
   + 1.11{\times}10^{-11}\,T_4^{-0.458},\nonumber\\
  k({\rm O}^+ + {\rm H}) &= \big(1.26{\times}10^{-9}\,T_4^{0.517}
   + 4.25{\times}10^{-10}\,T_4^{0.00669}\big)\,e^{-227/T}.
\end{align}
Higher stages ($\mathrm{X}^{i+} + \mathrm{H}^0 \to
\mathrm{X}^{(i-1)+}$, recombination direction) use the Kingdon \&
Ferland (1996) fits $k = a\,T_4^{b}\,(1 + c\,e^{d\,T_4})$, transcribed
from the MOCASSIN \code{update\_mod.f90} \code{chex} table and evaluated
at $T$ clamped into the fit validity ($\sim$5\,000--50\,000~K).
\emph{Stage mapping (corrected 2026-07-13).}  MOCASSIN's
\code{chex}$(Z,\mathrm{ion})$ is the reaction
$\mathrm{X}^{\mathrm{ion}+} + \mathrm{H}^0 \to
\mathrm{X}^{(\mathrm{ion}-1)+}$, and its inline comment labels the
\emph{product} (lower) ion.  The recombining ion of transition $i$ is
$\mathrm{X}^{i+}$, so the entry is \code{chex}$(Z,i)$.  An earlier
version read the product-labeled comment as the reactant and used
\code{chex}$(Z,i{+}1)$, placing every higher-stage rate one ionization
stage too low; the corrected assignment (verified against MOCASSIN's own
ratio loop and CMacIonize) is below.

\begin{center}\small
\begin{tabular}{lrrrr}
\toprule
reaction & $a$ [$10^{-9}\cmc$] & $b$ & $c$ & $d$ \\
\midrule
C$^{2+}$ + H  & $1.67{\times}10^{-4}$ & 2.79 & 304.72 & $-4.07$ \\
C$^{3+}$ + H  & 3.25 & 0.21 & 0.19 & $-3.29$ \\
N$^{2+}$ + H  & 0.305 & 0.60 & 2.65 & $-0.93$ \\
O$^{2+}$ + H  & 1.04 & 0.27 & 2.02 & $-5.92$ \\
Ne$^{2+}$ + H & $10^{-5}$ & 0 & 0 & 0 \\
S$^{2+}$ + H  & $10^{-5}$ & 0 & 0 & 0 \\
S$^{3+}$ + H  & 2.29 & 0.0402 & 1.59 & $-6.06$ \\
Ar$^{2+}$ + H & $10^{-5}$ & 0 & 0 & 0 \\
Ar$^{3+}$ + H & 4.57 & 0.27 & $-0.18$ & $-1.57$ \\
Si$^{2+}$ + H & 1.23 & 0.24 & 3.17 & $4.18{\times}10^{-3}$ \\
Si$^{3+}$ + H & 0.49 & $-0.0874$ & $-0.36$ & $-0.79$ \\
Si$^{4+}$ + H & 7.58 & 0.37 & 1.06 & $-4.09$ \\
Mg$^{2+}$ + H & $8.58{\times}10^{-5}$ & $2.49{\times}10^{-3}$ & 0.0293 & $-4.33$ \\
Fe$^{2+}$ + H & 1.26 & 0.0772 & $-0.41$ & $-7.31$ \\
Fe$^{3+}$ + H & 3.42 & 0.51 & $-2.06$ & $-8.99$ \\
\bottomrule
\end{tabular}
\end{center}
Mg and Fe were already correct: their Huang Table~4 reactant labels
happen to land on \code{chex}$(Z,i)$.  Silicon's first stage uses the
Huang A9/A10 fits ($\mathrm{Si} + \mathrm{H}^+$ and
$\mathrm{Si}^+ + \mathrm{H}^0$, the latter
$2.371{\times}10^{-12}\cmc$ at $10^4$~K); Cl and Ca have no KF96/Huang
transcription.  The Mg and Fe first-stage reactions follow Huang
Table~4: Mg + H$^+$ is the KF96 form with an added reverse barrier
$e^{-6.91/T_4}$ (endothermic); Fe + H$^+$ is a fast constant
$4.0{\times}10^{-9}\cmc$ with reverse $a\,(T/300)^{b}\,e^{-6.61/T_4}$.
The correction raises the HII40 doubly/triply-ionized fractions
(Section~\ref{sec:hii}); the earlier good agreement of some metal
fractions with the underconverged MOCASSIN baseline was partly
fortuitous.

\subsection{Cooling fits}
\label{sec:cool}
Collisional line cooling for the thermal loop uses the EXHALE
multi-exponential form fitted to direct CHIANTI v11.0.2 evaluations
(\code{fit\_cooling.py}; reader \code{chianti\_cooling.py} parses
\code{.elvlc}/\code{.scups} with Burgess \& Tully 1992 descaling):
\begin{equation}
  \Lambda_{\rm ion}(T) = T^{-1/2}\sum_{i=1}^{N} A_i\,e^{-T_i/T}
  \qquad [\mathrm{erg\,cm^3\,s^{-1}}\ \text{per } (n_e n_{\rm ion})],
  \label{eq:cool}
\end{equation}
evaluated in the optically thin $n_e\rightarrow0$ limit of statistical
equilibrium: every ion sits in its lowest level, excitation happens out
of that level only, and the cascade radiates the full excitation energy,
so $\Lambda_{\rm ion}$ is the excitation power out of the ground level
and equals what the $n$-level solve (Section~\ref{sec:nlev}) returns as
$n_e\rightarrow0$.  Transition energies are the observed \code{.elvlc}
level differences in both the Boltzmann factor and the radiated photon
(the \code{.scups} \code{de} is the Burgess--Tully scaling energy, and
detailed balance ties the excitation rate to the level gap that appears
in the equilibrium matrix).  Fit range $10^3$--$10^5$~K;
Levenberg--Marquardt on $\log\Lambda$ with the $T_i$ ladder derived from
the excitation temperatures $\Delta E_{1u}/k$ of each ion.

\begin{center}\small
\begin{tabular}{lcc@{\qquad}lcc}
\toprule
ion & $N$ & max err & ion & $N$ & max err \\
\midrule
H\,\textsc{i}   & 4 & 0.52\% & O\,\textsc{i}    & 6 & 0.20\% \\
C\,\textsc{i}   & 6 & 0.27\% & O\,\textsc{ii}   & 4 & 0.27\% \\
C\,\textsc{ii}  & 5 & 0.35\% & O\,\textsc{iii}  & 5 & 0.32\% \\
C\,\textsc{iii} & 6 & 0.50\%$^{a}$ & S\,\textsc{ii} & 3 & 2.23\% \\
N\,\textsc{i}   & 5 & 0.58\% & S\,\textsc{iii}  & 5 & 0.23\% \\
N\,\textsc{ii}  & 5 & 0.66\% & S\,\textsc{iv}   & 3 & 0.61\% \\
N\,\textsc{iii} & 5 & 0.70\% & Ne\,\textsc{ii}  & 5 & 0.40\% \\
 & & & Ne\,\textsc{iii} & 5 & 0.34\% \\
 & & & Ar\,\textsc{iii} & 2 & 0.22\%$^{b}$ \\
 & & & Ar\,\textsc{iv}  & 4 & 0.27\% \\
\bottomrule
\end{tabular}\\[2pt]
{\footnotesize $^a$in the relevant range ($\Lambda > 10^{-3}\max$);
5.5\% at the negligible low-$T$ tail.
$^b$at $10^4$~K; 3.9\% at the $10^5$-K endpoint.}
\end{center}
By construction Eq.~(\ref{eq:cool}) carries no collisional de-excitation,
so it is exact for lines whose critical density is far above the local
$n_e$ and increasingly too large for those below it --- the
low-$n_{\rm crit}$ lines [C\,\textsc{ii}] 158\,$\mu$m,
[O\,\textsc{iii}] 52/88\,$\mu$m and the [O\,\textsc{ii}]/[S\,\textsc{ii}]
optical doublets, whose electron critical densities are
$10^2$--$10^{3.5}\,{\rm cm^{-3}}$, so at the $n_e\sim10^2\,{\rm cm^{-3}}$
of an H\,\textsc{ii} region they overcool the line budget by $\sim$12\%
(measured against a \textsc{Cloudy} c25.00 cooling decomposition on a
fixed state).  The thermal loop therefore treats
Eq.~(\ref{eq:cool}) as the $n_e\rightarrow0$ \emph{carrier} and
suppresses it to the local density.  A $(T, n_e)$ table, built once at
setup (\code{nlevel\_cooling\_mod}) by solving the same $n$-level atom
over a grid ($\log_{10}T = 2$--$6$, $\log_{10}n_e = -6$--$10$), stores
for each ion the suppression ratio
\begin{equation}
  S_{\rm sup}(T,n_e) = \frac{\Lambda_{\rm SE}(T,n_e)}
       {\Lambda_{\rm SE}(T,n_e\!\to\!0)}\in(0,1],
  \qquad
  \Lambda_{\rm SE}(T,n_e) = \frac{1}{n_e}\sum_k n_u^{(k)}A_k\,\Delta E_k,
\end{equation}
and the low-density reference $\Lambda_{\rm SE}(T,n_e\!\to\!0)$; the
runtime cooling coefficient combines the two as
\begin{equation}
  \Lambda(T,n_e) = \Lambda_{\rm fit}(T)\bigl[\,f_{\rm cov}(T)\,
    S_{\rm sup}(T,n_e) + (1 - f_{\rm cov}(T))\,\bigr],
  \qquad
  f_{\rm cov} = \min\!\Bigl(1,\ \tfrac{\Lambda_{\rm SE}(T,n_e\to0)}
       {\Lambda_{\rm fit}(T)}\Bigr).
  \label{eq:coolcombine}
\end{equation}
Here $f_{\rm cov}$ is the fraction of the low-density cooling the (often
truncated) $n$-level model resolves; the remainder $1-f_{\rm cov}$ ---
the high-excitation resonance lines, of far higher critical density ---
is left in its low-density limit, so a truncated model cannot bias the
total.  As $n_e\rightarrow0$, $S_{\rm sup}\rightarrow1$ and
$\Lambda\rightarrow\Lambda_{\rm fit}$ \emph{exactly}, and an ion with no
$n$-level model would keep the plain fit --- there is none left, since
C\,\textsc{i} and N\,\textsc{i}, the last two gaps, now carry models of
their own (their far-infrared fine-structure lines have
$n_{\rm crit} \sim 2$--$10\,{\rm cm^{-3}}$, so the suppression is severe:
$S_{\rm sup} = 0.43$ at $10^4$~K and $n_e = 10^4\,{\rm cm^{-3}}$).
This density-suppressed cooling
(\code{par\%cooling\_model = local\_ne}, the default) is built from the
very solve that produces the diagnostic lines (Section~\ref{sec:nlev}),
so the thermal balance and the emergent line list are one physics; the
pure zero-density fit is retained as the \code{low\_density} option,
exact in the diffuse gas the limit is written for.  The
density-dependent emissivities themselves come from the $n$-level data.

\subsection{$n$-level (diagnostic) data}
\label{sec:nlev}
For output-time diagnostics, \code{fit\_nlevel.py} exports for
each ion the lowest $n$ levels (energies, statistical weights),
the radiative $A$-values (\code{.wgfa}), and the effective collision
strengths $\Upsilon(T)$ as low-order Chebyshev fits in Burgess--Tully
scaled space: with $e_t = kT/\Delta E$,
\begin{equation}
  s = \begin{cases} 1 - \ln C/\ln(e_t + C) & \text{types 1, 4}\\
  e_t/(e_t + C) & \text{types 2, 3, 5, 6}\end{cases},
  \qquad
  u = \frac{2(s - s_{\rm lo})}{s_{\rm hi} - s_{\rm lo}} - 1,
  \qquad y = \sum_k c_k T_k(u),
\end{equation}
descaled by type ($y\ln(e_t{+}e)$; $y$; $y/(e_t{+}1)$; $y\ln(e_t{+}C)$;
$y/e_t$; $10^y$), with a log flag storing $\log_{10} y$ for weak peaked
intercombination transitions.  The $n$-level export now covers the 24
line-emitting ions of the eleven-element registry
(Table~\ref{tab:metallines_phys}); the worst transition fit error is
1.7\% (O\,\textsc{iii} $^1$S$_0$--$^5$S$_2$, $\Upsilon\sim10^{-4}$),
typically $\lesssim 0.5\%$.  Ar\,\textsc{ii} has no level/collision data in
CHIANTI v11 (only RR/DR), so that stage is tracked without cooling or
lines ([Ar\,\textsc{ii}] 6.98\,$\mu$m omitted).  Verification is layered: the fitted files
reproduce direct CHIANTI spline evaluation of the classic ratios to
0.05--0.5\%; the PyNeb cross-check gives [O\,\textsc{iii}] 5007/4363
within 1.2\% and [N\,\textsc{ii}] 6584/5755 within 0.3\% (the density
diagnostics differ by 6--10\% at $10^4\,{\rm cm^{-3}}$ from PyNeb's
independent collision data --- database spread); and the Fortran solver
(Section~\ref{sec:nlevel}) matches the Python reference implementation
to four significant digits on 26 lines.

Table~\ref{tab:metallines_phys} lists the principal collisional
diagnostic lines each metal block produces (the complete transition list
for each ion is written to \code{<base>\_lines.txt}); only active-registry
ions with an \code{nlevel\_<el>\_<stage>} file get a block.  The
H\,\textsc{i}/He\,\textsc{ii}/He\,\textsc{i} recombination-line set is
documented separately in the recombination-lines subsection below.

\begingroup\small
\begin{longtable}{@{}l p{10.8cm}@{}}
\caption{Metal collisional diagnostic lines from the $n$-level solve
  (block \code{emis\_<el>\_<stage>} in \code{<base>\_emis}; principal lines,
  \AA\ unless $\mu$m noted).  C, N, O, Ne, S are active by default; Ar, Mg,
  Fe, Si, Cl, Ca are activated by \code{par\%abund\_<El>}.}
  \label{tab:metallines_phys}\\
\toprule
ion (block) & principal lines \\
\midrule
\endfirsthead
\multicolumn{2}{@{}l}{\small\itshape Table~\ref{tab:metallines_phys} (continued)}\\
\toprule
ion (block) & principal lines \\
\midrule
\endhead
\bottomrule
\endfoot
\code{c\_2} [C\,\textsc{ii}]   & 2325 (UV multiplet), 158\,$\mu$m \\
\code{c\_3} C\,\textsc{iii}]   & 977, 1907, 1909, 178\,$\mu$m \\
\code{n\_2} [N\,\textsc{ii}]   & 3064/3071, 5756, 6550, 6585, 122\,$\mu$m, 205\,$\mu$m \\
\code{n\_3} [N\,\textsc{iii}]  & 1750 (UV multiplet), 57\,$\mu$m \\
\code{o\_1} [O\,\textsc{i}]    & 5579, 6302, 6366, 63\,$\mu$m, 145\,$\mu$m \\
\code{o\_2} [O\,\textsc{ii}]   & 2471, 3727, 3730, 7321/7331 \\
\code{o\_3} [O\,\textsc{iii}]  & 1661/1666, 2322/2332, 4364, 4960, 5008, 52\,$\mu$m, 88\,$\mu$m \\
\code{ne\_2} [Ne\,\textsc{ii}] & 12.8\,$\mu$m \\
\code{ne\_3} [Ne\,\textsc{iii}]& 3343, 3870, 3969, 15.55\,$\mu$m, 36\,$\mu$m \\
\code{s\_2} [S\,\textsc{ii}]   & 4070/4078, 6718, 6733, 1.03\,$\mu$m, 315\,$\mu$m \\
\code{s\_3} [S\,\textsc{iii}]  & 3722, 6314, 8831, 9070, 9532, 12\,$\mu$m, 18.7\,$\mu$m, 33.5\,$\mu$m \\
\midrule
\multicolumn{2}{@{}l}{\itshape Optional --- activate with \code{par\%abund\_<El>}}\\
\midrule
\code{ar\_3} [Ar\,\textsc{iii}]& 3006/3110, 5193, 7138, 7753, 8.99\,$\mu$m, 21.8\,$\mu$m \\
\code{ar\_4} [Ar\,\textsc{iv}] & 2854/2869, 4713, 4741, 7170 (multiplet), 56\,$\mu$m, 78\,$\mu$m \\
\code{mg\_2} Mg\,\textsc{ii}   & 2796, 2804 \\
\code{fe\_2} [Fe\,\textsc{ii}] & 1.26/1.32\,$\mu$m, 5.34\,$\mu$m, 17.9\,$\mu$m, 26\,$\mu$m (40 lines total) \\
\code{fe\_3} [Fe\,\textsc{iii}]& 4659, 4703, 4881, 5011, 22.9\,$\mu$m (38 lines; 16/34-level with \code{fe\_levels\_full}) \\
\code{si\_2} [Si\,\textsc{ii}] & 2335 (multiplet), 34.8\,$\mu$m \\
\code{si\_3} Si\,\textsc{iii}] & 1206, 1883, 1892, 38\,$\mu$m \\
\code{si\_4} Si\,\textsc{iv}   & 1394, 1403 \\
\code{cl\_2} [Cl\,\textsc{ii}] & 3587/3679, 6164, 8581, 9126, 14.4\,$\mu$m, 33\,$\mu$m \\
\code{cl\_3} [Cl\,\textsc{iii}]& 3344/3354, 5519, 5539, 8480 (multiplet) \\
\code{cl\_4} [Cl\,\textsc{iv}] & 3120/3205, 5325, 7263, 7533, 11.8\,$\mu$m, 20.3\,$\mu$m \\
\code{ca\_2} Ca\,\textsc{ii}   & 3935, 3970, 7293/7326, 8500/8544/8665 \\
\code{ca\_5} [Ca\,\textsc{v}]  & 3999, 5311, 6088, 3.05\,$\mu$m, 4.16\,$\mu$m, 11.5\,$\mu$m \\
\end{longtable}
\endgroup
\noindent\small Some registry ions are tracked in the ionization balance but
carry no line block because CHIANTI v11 has no level/collision data for them:
Ar\,\textsc{ii}, Cl\,\textsc{v}, Ca\,\textsc{iii}, Ca\,\textsc{iv}.\normalsize

\subsection{Free-free Gaunt factors and continuum data}
\label{sec:contdata}
\code{gaunt.f90} is the verbatim MOCASSIN copy of the Hummer (1988)
two-dimensional Chebyshev fit of the Karzas \& Latter (1961) free-free
Gaunt factor ($\leq 0.7\%$; interface
\code{getGauntFF(z, log10Te, xlf, g)}).  The thermal loop uses the
thermally averaged $\bar g(T, Z) = \int g_{\rm ff}(u)\,e^{-u}\,du$
tabulated at setup ($\bar g(10^4\,{\rm K}, Z{=}1) = 1.2626$), now over
net charge $Z=1$--$4$ so the trace-metal free-free
(\code{par\%metal\_freefree}) has its Gaunt factor.  The van
Hoof et al.\ (2014) table is now available as a drop-in
(\code{par\%gaunt\_vh14}; \code{data/gauntff\_vh14.dat}, an
$81 \times 146$ grid in $\log_{10}\gamma^2 = \log_{10}(Z^2\,{\rm Ry}/kT)$
and $\log_{10} u = \log_{10}(h\nu/kT)$, selected at both call sites ---
the cooling average and the continuum $\gamma_{\rm ff}$).  The two
compilations agree: switching moves $T_e$ by a median
$5{\times}10^{-5}$, at most $0.5\%$, consistent with the sub-percent
Hummer accuracy van Hoof et al.\ report.  Hummer remains the default
so the recorded gates reproduce.

The nebular continuum tables \code{data/gamma\{HI,HeI,HeII\}.dat}
(MOCASSIN inheritance; 21 temperatures $\times$ paired-edge frequency
points) were \emph{numerically verified} (2026-07-12) against a
hydrogenic Kramers Milne sum: they contain free-bound \emph{plus}
free-free combined coefficients (at $\nu = 0.0069$~Ryd, $10^4$~K the
table value 9.47 decomposes into ff $\simeq 8.9$ + fb $\simeq 0.6$ in
$10^{-40}$ units; near the Balmer edge the residual is $\sim$10\%, the
size of the neglected bound-free Gaunt factor).  Consequently the port
never adds free-free inside the table range.  The Ercolano \& Storey
(2006) item is resolved as \emph{provenance}: the MOCASSIN 2.02 tables
ship in the ES06 compact paired-edge format, so no swap is needed ---
the verification above is the record.  Two-photon continua
follow MOCASSIN: Nussbaumer \& Schmutz (1984) spectral shape with
Pengelly (1964) $\alpha_{\rm eff}(2s)$ and the Osterbrock
$q(2^2{\rm S}{-}2^2{\rm P})$ collisional suppression for H\,\textsc{i};
the $Z$-scaled form with Storey \& Hummer $\alpha_{\rm eff}$ for
He\,\textsc{ii}; Almog \& Netzer (1989) shape
(\code{data/HeI2phot.dat}) with Benjamin, Skillman \& Smits (1999)
$\alpha_{\rm eff}$ for He\,\textsc{i}.

\subsection{H\,\textsc{i} and He\,\textsc{ii} recombination lines
(Storey \& Hummer 1995)}
\code{make\_sh95\_lines.py} parses the SH95 case-B data (\code{e1bx.d})
with the reference reader indexing
$k = (2n_{\rm top} - n_{ll} - n_l + 1)(n_l - n_{ll})/2 + n_{\rm top} -
n_u + 1$ and exports H$\alpha$, H$\beta$, H$\gamma$, H$\delta$,
P$\alpha$, P$\beta$, Br$\gamma$ on the full $(T, n_e)$ grid
(11 $\times$ 13) to \code{data/atomic/sh95\_hi\_caseB.txt}; Fortran
interpolates bilinearly in $(\log T, \log n_e)$.  Checks:
$4\pi j_{{\rm H}\beta}/(n_e n_p) = 1.235{\times}10^{-25}$ at
$(10^4\,{\rm K}, 10^2\,{\rm cm^{-3}})$ and
H$\alpha$/H$\beta = 2.864$ (canonical values).

He\,\textsc{ii} ($Z = 2$) comes from the \emph{original} CDS file
\code{e2b.d}, whose format differs from the repackaged \code{e1bx.d}
(six-token block headers; index
$k = (n_{\rm cut}-n_u)(n_{\rm cut}+n_u-1)/2 + n_l$ with
$n_{\rm cut} = 25$, from the SH95 \code{intrat.f} reader): six lines
(1640, 4687, 3204, 10126, 5413, 4543~\AA) on a $12 \times 13$ grid in
\code{data/atomic/sh95\_heii\_caseB.txt}, weighted in the line table by
$n_e n_{\rm HeIII}$.  Checks:
$4\pi j_{4687}/(n_e n_{\rm HeIII}) = 1.495{\times}10^{-24}$ at
$(10^4\,{\rm K}, 10^4\,{\rm cm^{-3}})$
$= \alpha^{\rm eff}_{4687}(3.57{\times}10^{-13})\,h\nu$, and
$1640/4687 = 6.56$ there.

He\,\textsc{i} lines use the same interpolator (a third
table in \code{sh95\_mod}) but a different source: the Porter et al.\
(2012, 2013) case-B emissivities, parsed by \code{make\_hei\_lines.py}
from Cloudy c23.01's \code{he1\_case\_b.dat} (copied into
\code{data/atomic}) onto a $21 \times 14$ $(T, n_e)$ grid.  Unlike the
H\,\textsc{i} and He\,\textsc{ii} tables, these are case-B
collisional-radiative TOTALS, not pure recombination cascades: the
tabulated $4\pi j/(n_e n_{\rm He^+})$ is density-dependent (the 10833~\AA\
coefficient rises 6.6x from $n_e = 10$ to $10^4$~cm$^{-3}$ at $10^4$~K)
because it already includes collisional excitation out of the 2$^3$S
metastable.  Eight H\,\textsc{ii}-region diagnostics are exported
(3890, 4027, 4473, 5877, 6680, 7067, 7283, 10833~\AA, vacuum), weighted in
the line table by $n_e n_{\rm HeII}$ (He$^+$ recombining).  Emergent
ratios are physical: $5877/4473 = 2.88$, $6680/4473 = 0.82$,
$7067/4473 = 0.49$ at $\sim$7~kK.

The H\,\textsc{i} and He\,\textsc{ii} tables follow \code{par\%case\_ab}:
case-A tables (\code{sh95\_\{hi,heii\}\_caseA.txt}, parsed from the
original CDS files \code{e1a.d}/\code{e2a.d} onto the $10\times9$/$12\times9$
grids case A tabulates) are loaded when the ionization runs case A, so the
line emissivities match the balance; case B (H$\alpha$/H$\beta = 2.94$) is
the default.  The He\,\textsc{i} Porter set is case B only.  An optional
collisional H\,\textsc{i} Balmer component (\code{par\%h\_coll\_effects},
off by default) solves a 25-level H\,\textsc{i} atom (from CHIANTI
\code{h\_1} through the $n$-level pipeline) with collisional excitation from
the ground level, groups the $l$-resolved emission into
H$\alpha$/H$\beta$/H$\gamma$, weights it by $n_{\rm HI}$, and writes it as
separate \code{hc} rows --- a first-cut estimate (optically-thin cascade)
for the front and hot-cell bias, kept apart from the recombination lines.

\subsection{He\,\textsc{i} 2$^3$S metastable diagnostic}
\label{sec:heimeta}
The Porter case-B 10833~\AA\ table above is a collisional-radiative
\emph{total} emissivity, correct in the dense, optically-thin,
recombination-dominated H\,\textsc{ii} regime it was built for.  It does
\emph{not} give the 2$^3$S population $n_3$, its line-of-sight column, or
the 10833~\AA\ \emph{absorption} opacity --- the observables of exoplanet
transit, chromospheric, and PDR 10833 work, where the line is seen against
a background.  \code{hei\_metastable\_mod} (\code{par\%hei\_metastable})
adds these as a separate diagnostic; it does not add or modify any emission
row, and in particular \emph{nothing is summed onto} the Porter 10833 line
(which already carries the metastable-collisional term --- summing an
explicit collisional 10833 would double-count).

\paragraph{The 2$^3$S balance (Oklopcic \& Hirata 2018).}  OH18 track the
neutral singlet ground 1$^1$S, the metastable triplet 2$^3$S, and the ion.
Because $n_3/n_{\rm He}\sim10^{-7}$ the balance is linear in $n_3$ --- a
closed-form quotient in each leaf, no iteration:
\begin{equation}
n_3 = \frac{\alpha_3\,n_e\,n_{\rm HeII} + q_{13}\,n_e\,n_{\rm HeI}}
           {A_{31} + \Phi_3 + n_e\,(q_{31a}+q_{31b}) + n_{\rm HI}\,Q_{31}},
\end{equation}
with source terms recombination into the triplet system
($\alpha_3\,n_e\,n_{\rm HeII}$) and e-impact excitation
1$^1$S\,$\to$\,2$^3$S ($q_{13}\,n_e\,n_{\rm HeI}$), and sinks the radiative
decay $A_{31}$ (2$^3$S\,$\to$\,1$^1$S), photoionization $\Phi_3$,
spin-changing e-impact transfer into the singlets
($q_{31a} = $ 2$^3$S\,$\to$\,2$^1$S, $q_{31b} = $ 2$^3$S\,$\to$\,2$^1$P ---
\emph{not} 2$^3$S\,$\to$\,1$^1$S), and Penning/associative ionization with
neutral hydrogen ($n_{\rm HI}\,Q_{31}$).  Two limits anchor the gates: the
low-density radiative limit $n_3/n_{\rm He^+} = \alpha_3 n_e/A_{31}$ and the
high-density collisional plateau $\alpha_3/(q_{31a}+q_{31b})$, crossing at
$n_{\rm crit} = A_{31}/(q_{31a}+q_{31b})$.

The rate coefficients are closed-form fits transcribed from the EXHALE
triplet implementation (offline into the coefficient file, Fortran only
evaluates): $\alpha_3(T) = 2.10{\times}10^{-13}(T/10^4)^{-0.778}$ (OH18);
$q_{13}, q_{31a}, q_{31b}$ from Bray et al.\ (2000) collision strengths via
OH18; $A_{31} = 1.272{\times}10^{-4}$~s$^{-1}$ (the CHIANTI \code{he\_1.wgfa}
$A(1\!-\!2) = 1.73{\times}10^{-4}$ is $\sim$36\% high and is \emph{not} used);
$Q_{31}$ from Taylor et al.\ (2025) Table~2 (Garcia Mu\~noz 2025 recorded as
the alternative, 1.2--2.2$\times$ lower --- the documented uncertainty of the
neutral/PDR sink).  A consistency check: at $10^4$~K,
$\alpha_3 = 0.998\times0.75\,\alpha_B({\rm He\,II})$, matching the
\code{HEI\_F3S} $=0.75$ triplet branching the diffuse field already uses.

\paragraph{The 2$^3$S photoionization rate.}
$\Phi_3 = \sum_\nu 4\pi J_\nu\,\sigma_3(\nu)\,\Delta\nu/(h\nu)$ over the
\emph{full} band (the reduced tally \code{jt\_ion}, the same integral pattern
as $\Gamma_i$).  The threshold is 4.78~eV, but $\sigma_3 \sim
1$--$2{\times}10^{-18}$~cm$^2$ remains significant far above 13.6~eV
(Cooper minimum near 34.7~eV, bump near 45.6~eV), so the ionizing band alone
gives $\Phi_3 > 0$ wherever a field exists; \code{add\_fuv} with
\code{efuv\_min}$=4.78$ adds the soft-UV part $[4.78, 13.6)$~eV (a rank-0 note
fires when the band floor sits above 4.78~eV).  $\sigma_3(E)$ is two VFKY96
wings joined by a log-linear bridge, a $\sim$4\% fit to Norcross (1971)\,+\,TOPbase.

\paragraph{10833~\AA\ absorption (Stage 2).}  From $n_3$ the line-center
opacity of each fine-structure component is
\begin{equation}
\kappa_{0,i} = \frac{\sqrt{\pi}\,e^2}{m_e c}\,f_i\,\frac{\lambda_i}{b}\,n_3,
\qquad b = \sqrt{2 k T_e/m_{\rm He} + v_{\rm turb}^2},
\end{equation}
the classical integrated cross section $\pi e^2/(m_e c) = 0.02654$~cm$^2$\,Hz
divided by $\sqrt{\pi}\,b$ for the Gaussian peak.  The three
2$^3$S\,$\to$\,2$^3$P$_J$ components (vacuum 10832.06 / 10833.22 /
10833.31~\AA; $f = 0.0599 / 0.1797 / 0.2996$, ratio 1:3:5; $A =
1.0216{\times}10^7$~s$^{-1}$ each) take the NIST ASD values (Drake 2007),
cross-checked against p-winds; the CHIANTI \code{he\_1.wgfa} $A \simeq
1.15{\times}10^7$ ($\sim$13\% high) is \emph{not} used.  The $J'=1,2$ pair sits
0.089~\AA\ apart, far inside the $\sim$0.4~\AA\ thermal Doppler width at
$10^4$~K, so they blend into the observed ``doublet'' (their line-center
opacities are summed); $J'=0$ at 10832.06~\AA\ is resolved.  The Python
\code{HeiMetaData} reader deposits $n_3$ and $\kappa_0$ along the line of
sight (the flux-conserving deposit of the emissivity maps) into column
[cm$^{-2}$] and $\tau_0 = \int\!\kappa_0\,{\rm d}l$ maps, and warns when
line-center $\tau_0 > 1$.

The module runs once at output time (no hot-loop cost, no feedback --- $n_3$
is a $\sim$10$^{-7}$ trace of He).  Caveats recorded at the code site:
(i)~\emph{local equilibrium} --- the advection term $v\,{\rm d}f_3/{\rm d}r$
is absent, valid where the 2$^3$S destruction time is short against any flow
time (dense H\,\textsc{ii}), not a faithful outflow benchmark; (ii)~10833 can
be optically thick in dense metastable columns, where the optically-thin
$\tau_0$ diagnostic breaks down (line transfer is out of scope).

%==========================================================================
\section{Algorithms}
\subsection{The iteration loop}
\begin{center}
\begin{minipage}{0.86\textwidth}
\begin{tabbing}
xxx\=xxx\=xxx\=\kill
read grid $\rightarrow$ persist gas leaf state $\rightarrow$ fill \code{kap\_ion}\\
\textbf{do} iter = 1, \code{gas\_niter}\\
\> zero the band tally $J_b$\\
\> sample and transport stellar packets (analytic edge walk, Eq.~\ref{eq:edgewalk})\\
\> [diffuse] rebuild the diffuse emission table from the current state;\\
\> \> sample and transport $L_{\rm dif}/L_{\rm pkt}$ diffuse LyC packets\\
\> ALLREDUCE $J_b$; rate integrals (Eq.~\ref{eq:rates}); metal $\Gamma$'s\\
\> solve each leaf: ionization only or coupled $(x, T_e)$\\
\> refill \code{kap\_ion} from the new state (the stellar tally is\\
\> \> recomputed from zero each iteration --- no reuse)\\
\> converged by \code{par\%conv\_crit}: cell-max or volume-integrated
   tests (below)\\
\textbf{end do}\\
outputs: rates + state maps; line luminosities; nebular continuum
\end{tabbing}
\end{minipage}
\end{center}
Every rank runs the identical cell solves (inputs are the reduced
tallies), so all ranks agree on convergence; only the node master writes
the shared state arrays.  Under-relaxation on $x$
(\code{par\%ion\_relax}) is available for sharp I-fronts.  Practical
lessons recorded: a fully neutral start advances the front only
$\sim$one cell per iteration (the photons die in the first neutral
cell), while an ionized start converges globally; and
$\max|\Delta x_{\rm HII}|$ stalls at the front-cell Monte-Carlo noise
floor ($\sim$2--4${\times}10^{-2}$ at $4{\times}10^7$ packets), so
tolerances below that are unreachable by that criterion --- the
diffuse field, interestingly, smooths the front and converges to
$3{\times}10^{-3}$.

\subsection{Stellar packet launch}
\label{sec:stellar-launch}
For each direct packet, MoCHII first samples its energy with
Eq.~(\ref{eq:band-source-cdf}); it then samples an isotropic point-source
direction with $\mu=2u_\mu-1$ and $\phi=2\pi u_\phi$.  The source position,
the exact energy, the direction, and the common packet luminosity are placed
in the packet before the octree edge walk begins.  Thus energy sampling is
not conditional on the diagnostic binning.  At every traversed cell the
path estimator uses the packet's exact-energy opacity and cross sections;
the energy-bin index is accumulated only alongside this estimator to form
the reported spectrum.  Rebuilding the direct tally from zero on every
nonlinear iteration means that the changing opacity is always applied to a
new sample of the source spectrum rather than to a reweighted old tally.

For ordinary pseudo-random launching the inverse-CDF variate and the two
direction variates are independent Mersenne-Twister draws.  The same
probability law is used by the Sobol option described next.  The packet
weight is not divided by the probability density: equal packet luminosity
combined with inverse-CDF sampling is the importance sampling rule that
recovers the source luminosity spectrum in expectation.

\paragraph{Convergence measures.}  Two pairs are computed every
iteration and \code{par\%conv\_crit} selects which gates the loop:
the cell-max tests ($\max|\Delta x_{\rm HII}| <$ \code{gas\_tol},
$\max|\Delta T_e|/T_e <$ \code{gas\_tol\_te}) and the
volume-integrated tests
\begin{equation}
  \frac{\sum_{\rm leaf} V\,|\Delta x_{\rm HII}|}
       {\sum_{\rm leaf} V\,x_{\rm HII}} < \code{gas\_tol},
  \qquad
  \frac{\sum_{\rm leaf} V\,n_e\,|\Delta T_e|}
       {\sum_{\rm leaf} V\,n_e\,T_e} < \code{gas\_tol\_te}.
\end{equation}
The volume measures fix both stall modes of the cell-max tests: front
jitter occupies little volume (the front shell is a few thousand cells
against $\sim$10$^5$ ionized ones, so its noise enters the $L_1$ norm
suppressed by that ratio), and the $n_e$ weight removes the vacuum and
no-heating cells whose \code{te\_min}-pinned oscillation dominates
$\max|\Delta T_e|$ in FUV/PDR configurations.  \code{'vol'} is the
recommended production setting; \code{'cell'} (the historical
behavior) remains the default so the recorded gates reproduce.
Measured on the FUV dusty Str\"omgren case at $8{\times}10^6$ packets
(where the cell-max tests sat at 0.1 and 6.4 indefinitely): the volume
measures decay smoothly to well-defined Monte Carlo floors of
$2.7{\times}10^{-3}$ ($x$) and $6.6{\times}10^{-3}$ ($T_e$) by
iteration $\sim$27, and running 30 further iterations leaves the
solution unchanged ($V_{\rm ion}$ to 0.014\%) --- the floor IS
convergence.  Set the tolerances just above the floor for the packet
count in use (here $5{\times}10^{-3}$ and $10^{-2}$ would stop the
loop at iteration $\sim$23); the floors scale as
$1/\sqrt{N_{\rm packets}}$.
With \code{par\%require\_convergence} (default off) a non-converged gas
iteration prints a warning and, when the switch is set, aborts before
writing output; either way the rates-file header records the convergence
status in the keywords \code{CONVERGD}, \code{NITERDON}, \code{FINALDX},
and \code{FINALDTE}.

\subsection{Quasi-random photon launching (optional)}
With \code{par\%launch\_sequence = 'sobol'} the packet launch draws its
variables --- source component where relevant, energy, direction, and any
entry-surface coordinates --- from an Owen-scrambled Sobol net
(\code{src/qmc\_mod.f90}) instead of the Mersenne Twister, while every
post-launch transport event stays pseudo-random.  The Sobol point is
addressed by the \emph{global} photon number, so the launch set is
independent of the MPI task count, and \code{par\%qmc\_seed} is held fixed
across the nonlinear iterations (common random numbers).  Because the band
has no scattering, the direct field is an analytic function of the launch
variables (Eq.~\ref{eq:edgewalk}), so a low-discrepancy launch set reduces
its noise directly: on the fixed-state analytic-attenuation gate the
cell-wise $\Gamma_{\rm HI}$ error scales as $\sim 1/N$ (Monte Carlo
$1/\sqrt{N}$), with effective photon gains of $29\times$ at $N = 2^{17}$ and
$222\times$ at $2^{20}$, at unchanged wall time and no significant bias.  The
method, the fixed dimension layouts, and all gates are documented in the
companion memo \code{MoCHII\_QMC}.  In the validated single point-source
continuous mode the three Sobol coordinates are
$(u_E,u_\mu,u_\phi)$: $u_E$ is inverted through the source CDF and the
remaining coordinates set the isotropic direction.  QMC therefore reduces
the discrepancy of the energy samples as well as of the angular samples;
it does not replace continuous energy sampling by bin selection.  Diffuse
packets use a separate scrambled stream and the nine coordinates listed in
\S\ref{sec:diffuse}.

\subsection{Ionization equilibrium}
\label{sec:ioneq}
Per leaf at temperature $T$ (\code{solve\_ion\_cell}): with
$R_i = \Gamma_i + n_e C_i(T)$ the ionization rates and $\alpha_i$ the
case-A/B recombination rates,
$\xHI = \alpha n_e/(R_{\rm HI} + \alpha n_e)$ and the He stage chain
$n_{i+1}/n_i = R_i/(\alpha_{i+1} n_e)$, closed by
$n_e = n_{\rm H}[x_{\rm HII} + y_{\rm He}(x_{\rm HeII} +
2x_{\rm HeIII})]$ and solved by a damped fixed point on $n_e$ (mixing
$\tfrac12$, converged at $10^{-10}$, $\leq 200$ iterations).

\subsection{Thermal balance}
\label{sec:thermal}
When \code{par\%solve\_te} is set, \code{gas\_thermal\_update} replaces
the fixed-$T$ \code{gas\_equilibrium\_update} (\S\ref{sec:ioneq}) and
updates each leaf's temperature once per iteration, from that iteration's
radiation field, by finding the equilibrium
$\mathcal{H}(T) = \mathcal{C}(T)$ --- the root of the net rate
\begin{equation}
  \mathcal{N}(T) = \mathcal{H}(T) - \mathcal{C}(T),\qquad
  \mathcal{H} = n_{\rm HI}\mathcal{H}_{\rm HI} + y_{\rm He} n_{\rm H}
  (x_{\rm HeI}\mathcal{H}_{\rm HeI} + x_{\rm HeII}\mathcal{H}_{\rm HeII}),
  \label{eq:netrate}
\end{equation}
with the photoheating rates $\mathcal{H}_i$ (per atom) from
Eq.~(\ref{eq:rates}).  The ionization state is \emph{re-solved}
(\code{solve\_ion\_cell}, \S\ref{sec:ioneq}) at every trial $T$: both
$\mathcal{H}$ and $\mathcal{C}$ depend on the fractions
$(x_{\rm HI}, x_{\rm HeI}, x_{\rm HeII}, n_e)$, which themselves depend on
$T$.  The ionization response is mild, so $\mathcal{N}(T)$ is monotone
through its root and bisection on $\log T_e$ in
[\code{te\_min}, \code{te\_max}] is safe: each step takes the geometric
midpoint $\sqrt{T_{\rm lo}T_{\rm hi}}$ (up to 60 steps, stopping at
$T_{\rm hi}/T_{\rm lo}-1 < 10^{-5}$), and if cooling already wins at
\code{te\_min} ($\mathcal{N}\le 0$) the leaf is pinned to \code{te\_min}
while if heating still wins at \code{te\_max} it is pinned to
\code{te\_max}.

The cooling $\mathcal{C}(T)$ (\code{cooling\_mod}) assembles recombination
(Hui--Gnedin $\beta$), free-free
($1.42554{\times}10^{-27}\sqrt{T}\,Z^2\,\bar g\,n_e n_{\rm ion}$, summed
over H\,\textsc{ii}, He\,\textsc{ii} ($Z=1$) and He\,\textsc{iii} ($Z=2$)
with $Z$ the \emph{net} ion charge; under \code{par\%metal\_freefree} (on
by default) the trace metals are added with $Z_{\rm ion}=$ (stage $-1$),
not the nuclear $Z$, using $\bar g(T, Z_{\rm ion})$ tabulated for
$Z=1$--$4$ --- singly-ionized metals dominate this small
$\sim$$10^{-4}$-of-the-H/He term),
collisional ionization ($k_{\rm ci}\times$ ionization potential),
H\,\textsc{i} line cooling, and the trace-metal sum
$\sum n_e\,n_{\rm ion}\Lambda_{\rm ion}(T,n_e)$ with the cascade fractions
(\S\ref{sec:cascade}) re-evaluated at each trial $T$.  The line-cooling
coefficient $\Lambda_{\rm ion}(T,n_e)$ is the density-suppressed value of
Eq.~(\ref{eq:coolcombine}) by default (\code{par\%cooling\_model =
local\_ne}), the plain fit $\Lambda_{\rm ion}(T)$ under
\code{low\_density}.  PDR runs add three
optional terms to $\mathcal{N}$: metal photoheating
(\code{par\%metal\_heat}); grain photoelectric heating with grain
recombination cooling (Bakes \& Tielens 1994, \code{par\%grain\_pe},
driven by the local FUV field $G_0$, both weighted by $x_{\rm HI}$ and
evaluated at $\min(T, 10^4\,{\rm K})$ so the fits stay in range --- two
guards the PDR gate needed against a 12--35~kK runaway); and the
H-impact [C\,\textsc{ii}] 158\,$\mu$m / [O\,\textsc{i}] 63\,$\mu$m
fine-structure cooling that balances the photoelectric heating in the
mostly neutral gas.

Under \code{par\%use\_sec\_ion} (Shull \& van Steenberg 1985) a
photoelectron of kinetic energy $E_0 = h\nu - E_{{\rm th},s}$ from absorber
$s$ (H\,\textsc{i}, He\,\textsc{i}, He\,\textsc{ii}) with $E_0 > 40$~eV
deposits only $f_{\rm heat}(x)$ of its excess as heat (below the threshold
it thermalizes fully), so the photoheating integrand is weighted by
$f_{\rm heat}(x)$; the un-thermalized balance drives
$f_{\rm ion,HI}(x)\,E_0/E_{\rm th,HI}$ secondary H\,\textsc{i} and
$f_{\rm ion,HeI}(x)\,E_0/E_{\rm th,HeI}$ secondary He\,\textsc{i}
ionizations that add to the corresponding photoionization rates in
\code{solve\_ion\_cell}.  Here
$x = (n_{\rm HII}+n_{\rm HeII}+n_{\rm HeIII})/(n_{\rm H}+n_{\rm He})$ is the
ionized fraction of the H+He nuclei and the partitions are
$f_{\rm heat} = 0.9971[1-(1-x^{0.2663})^{1.3163}]$,
$f_{\rm ion,HI} = 0.3908(1-x^{0.4092})^{1.7592}$,
$f_{\rm ion,HeI} = 0.0554(1-x^{0.4614})^{1.6660}$.  The effect vanishes in
fully ionized gas ($f_{\rm ion}\to 0$ as $x\to 1$, so the HII-region gates
are unchanged) and matters only in partially neutral gas under a hard
field (fronts, PDR, X-ray/AGN); metals are neglected as absorbers (trace).
Default is off.

The solved $T_e$ is written directly (only the ionization fractions carry
the under-relaxation \code{par\%ion\_relax}), $n_e$ is rebuilt from the
new state --- including the trace-metal electrons under
\code{par\%metal\_ne} --- and, as in the ionization-only path, only the
node-local rank~0 runs the bisection (the reduced band tally is identical
on every node master) and writes the shared state, with the convergence
metrics broadcast within the node.  The loop is gated by the $T_e$
convergence measures defined above (\code{gas\_tol\_te}, cell-max or
volume-integrated per \code{par\%conv\_crit}).

\subsection{The trace-metal cascade}
\label{sec:cascade}
The registry (\code{species\_mod}) activates an element when its
abundance is positive.  In the trace approximation (no feedback on
$n_e$ or opacity), stage fractions follow the product chain
\begin{equation}
  \frac{n_{i+1}}{n_i} \;=\;
  \frac{\Gamma_i + n_e C_i + n_{\rm HII}\,k^{\rm CX}_{{\rm ion},i}}
       {n_e\,[\alpha_{\rm RR} + \alpha_{\rm DR}]_{i+1}
        + n_{\rm HI}\,k^{\rm CX}_{{\rm rec},i+1}},
\end{equation}
with metal $\Gamma_i$ computed once per iteration from the same band
tally.  Adding an ion is one entry in the data generators plus its
cooling and $n$-level files.

\subsection{Metal opacity in the He-ionizing window}
\label{sec:metalop}
``Trace'' names the \emph{data} treatment, not a switched-off feedback.  The
founding plan lists metals feeding back on $n_e$ and on the opacity as a
non-goal, but the production inputs run with \code{par\%metal\_ne},
\code{par\%ion\_metal\_abs} and \code{par\%metal\_heat} all true, so both
feedbacks are active: \code{species\_opacity\_add} sums over \emph{all}
active stages with their VFKY96 cross sections and the self-consistent stage
fractions, and the sum is iterated together with the opacity it changes.
The cross sections themselves are not in question.  For the eight ions that
carry the He-ionizing window --- C\,\textsc{i}--\textsc{iii},
N\,\textsc{ii}, O\,\textsc{ii}, Ne\,\textsc{i}, S\,\textsc{ii},
S\,\textsc{iii} --- MoCHII reproduces Verner's own \code{phfit2.f} to
$10^{-6}$ relative, the single-precision round-off of \code{phfit2} itself.

What is genuinely absent is metal recombination continua as a diffuse
\emph{source}: the diffuse field has exactly four channels
(\code{dif\_ch(4, nleaf)}) and all four are H/He.  Summing the metal
ground-state recombination continua whose photons lie above 24.587~eV
(C\,\textsc{iv}, N\,\textsc{iii}, N\,\textsc{iv}, O\,\textsc{iii},
O\,\textsc{iv}, Ne\,\textsc{iii}, Ne\,\textsc{iv}, S\,\textsc{iii},
S\,\textsc{iv}) gives $0.29\%$ of the He\,\textsc{i} ground-continuum
channel, about $0.001$~pc of front position.  Metals as secondary-ionization
absorbers are absent as well, but \code{par\%use\_sec\_ion} is off in the
benchmarks.

\paragraph{The measured lever.}  Scaling the summed metal cross section on
the converged HII40 model (68 ranks, same Sobol seed, $0.01$-pc shell
diagnostic) gives
\begin{center}\small
\begin{tabular}{@{}lrrr@{}}
\toprule
variant & $r({\rm He^0}{=}0.5)$ [pc] & $\Delta r$ & pc per unit fraction \\
\midrule
production                 & 4.11970 & ---      & ---   \\
metal $\sigma \times 0.70$ & 4.04764 & $-0.07206$ & 0.240 \\
metal $\sigma \times 1.30$ & 4.15063 & $+0.03093$ & 0.103 \\
metal absorption off       & 3.50470 & $-0.61500$ & 0.615 \\
metals off entirely        & 1.93335 & $-2.18635$ & (also removes all metal
  cooling; $T_e$ 8162 $\to$ 18681~K) \\
\bottomrule
\end{tabular}
\end{center}
The response is strongly asymmetric and saturating: removing metal opacity
costs $0.615$~pc per unit while adding it gains only $0.103$~pc per unit.
Extrapolating the upward branch, closing the $0.178$~pc Cloudy residual of
\S\ref{sec:diffuse} would take $+180\%$ metal opacity ($\times 2.8$), and the
saturation means still more.  That is not available, and it retires the
metal opacity as a single-parameter explanation.

The prediction recorded before these runs was ``metal opacity $\times 1.30$
moves $r({\rm He^0}{=}0.5)$ outward, $\times 0.70$ inward, of order
$0.1$--$0.2$~pc''.  The \emph{sign} was right both ways; the
\emph{magnitude} was wrong by factors of $1.5$--$6$, because the prediction
extrapolated the on/off bound of $0.615$~pc as if the response were
symmetric.  An on/off bound measures what removal costs and says little
about what addition could gain.  The companion statement that ``a ${\sim}30\%$
error in the metal opacity would be worth ${\sim}0.18$~pc'' is accordingly
withdrawn: $30\%$ is worth $0.031$--$0.072$~pc.

\paragraph{Mechanism: spectral hardening, not cooling.}  The metal stages
with thresholds inside $21$--$40$~eV (Ne\,\textsc{i} 21.56,
S\,\textsc{ii} 23.33, C\,\textsc{ii} 24.38, N\,\textsc{ii} 29.60,
S\,\textsc{iii} 34.83, O\,\textsc{ii} 35.12~eV) preferentially remove the
\emph{softest} He-ionizing photons, and
$\sigma({\rm He\,\textsc{i}})/\sigma({\rm H\,\textsc{i}})$ rises steeply
with energy --- $6.1$ at 25~eV, $9.0$ at 35, $13.9$ at 55, $19.9$ at
95~eV.  The hardened surviving field is therefore taken up by He$^0$ rather
than H$^0$ and the He$^+$ zone extends; remove the metal opacity and the
field stays soft, H$^0$ takes a larger share, and the zone collapses.  The
alternative readings are excluded by the same runs: between production and
\code{ion\_metal\_abs = .false.}\ the hydrogen structure and the mean
temperature barely move ($r({\rm H^0}{=}0.1) = 4.61633$ vs.\ $4.61551$~pc;
$T_e$ 8162 vs.\ 8183~K), so the $0.615$~pc is neither a thermal effect nor a
change in the hydrogen solution.

\paragraph{What Cloudy has here that MoCHII does not.}  Cloudy's base is the
same Verner \code{phfit}.  Its
\code{opacity\_\allowbreak createall.cpp} then adds
channels MoCHII has no equivalent for: N$^+$ excited-level photoionization
(Henry 1970), O\,\textsc{i} valence replaced wholesale by \code{ofit} when
PHFIT96 is in use, O\,\textsc{i} $^1$S, O$^{2+}$ $^1$D and $^1$S,
photoionization into excited states of O$^+$, and a Mg$^+$ excited level,
several of them fitted to Opacity Project data.  So Cloudy photoionizes
metals out of excited levels and MoCHII only out of the ground level.  As
opacity these are small: the excited-level populations are
$10^{-5}$--$10^{-4}$ of the ground, so the O$^{2+}$ $^1$D opacity is of order
$10^{-4}$ of the O$^+$ ground opacity.

An independent Opacity Project comparison is possible for C, N and O
through the TOPbase tables carried by \code{sirocco}.  Photon-weighted with
a 40-kK blackbody over $24.587$--$54.4$~eV, OP/VFKY96 is $0.946$
(C\,\textsc{ii}), $0.978$ (N\,\textsc{ii}) and $0.863$ (O\,\textsc{ii}) ---
OP is $2$--$14\%$ \emph{lower}.  With the measured upward derivative that is
$-0.005$ to $-0.03$~pc, the wrong direction for the residual.  Ne and S are
absent from that table and were not checked.

\paragraph{A justification that had to be retracted in the code.}  Two
comments --- at the \code{ion\_metal\_abs} declaration in \code{define.f90}
and above \code{species\_opacity\_add} --- used to justify the metal opacity
as ``negligible next to H/He above 13.6~eV''.  The measurement above
contradicts that, and both now carry the measured derivatives instead: the
opacity is small in \emph{magnitude} but not in \emph{effect}, because it
acts on the spectral shape in the band where He$^0$ and H$^0$ compete.  An
on/off statement is the form such a justification has to take here, since
the response is asymmetric and a single number in either direction
misstates it.

\subsection{Explicit diffuse LyC field and packet sampling}
\label{sec:diffuse}
Case B on the spot is replaced by case A rates plus explicit emission of
ground-recombination packets.  For each leaf and each of the three ground
continua H\,\textsc{ii}$\rightarrow$H\,\textsc{i}($1s$),
He\,\textsc{ii}$\rightarrow$He\,\textsc{i}($1^1S$), and
He\,\textsc{iii}$\rightarrow$He\,\textsc{ii}($1s$), the emissivity uses
$\alpha_1(T)=\alpha_A(T)-\alpha_B(T)$.  At the start of every gas
iteration the code rebuilds the energy-luminosity table
\begin{equation}
 L_{{\rm ch},l} = n_{e,l}n_{{\rm rec},l}\alpha_{1,\rm ch}(T_l)
      \left<E\right>_{n,\rm ch}(T_l)\,V_l,
 \qquad L_{\rm dif}=\sum_{l,{\rm ch}} L_{{\rm ch},l},
 \label{eq:diffuse-luminosity}
\end{equation}
with $\left<E\right>_n$ the photon-number-weighted mean energy of the
emitted continuum.

\paragraph{Ground-continuum spectrum (\code{par\%diffuse\_energy\_model}).}
Detailed balance between photoionization and radiative recombination fixes
the free--bound spectrum.  Per unit photon energy the photon-number
emissivity of each ground continuum is
\begin{equation}
 n(E)\;\propto\;E^2\,\sigma_{\rm pi}(E)\,
   \exp\!\left[-\frac{E-E_{\rm th}}{kT_e}\right],
 \qquad E\ge E_{\rm th},
 \label{eq:milne-pdf}
\end{equation}
evaluated with the \emph{same} ground-state cross section the transport
absorbs with, so a packet is absorbed with the cross section it was emitted
with.  This is the default, \code{'milne'}.  Restoring the Milne prefactor
turns the same integral into the ground-state recombination coefficient,
\begin{equation}
 \alpha_1(T)=\frac{g_n}{2g_+}\,\frac{2}{\sqrt{\pi}}\,(kT)^{-3/2}
   \sqrt{\frac{2}{m_e}}\,\frac{1}{m_ec^2}
   \int_{E_{\rm th}}^{\infty}\!E^2\sigma_{\rm pi}(E)
   e^{-(E-E_{\rm th})/kT}\,dE ,
 \label{eq:milne-alpha1}
\end{equation}
so the spectrum and the emission rate are not independent assumptions.
Equation~(\ref{eq:milne-alpha1}) reproduces
$\alpha_{1s}({\rm H\,\textsc{i}},10^4~{\rm K})=1.579\times10^{-13}$
cm$^3$\,s$^{-1}$ against the tabulated $1.58\times10^{-13}$.  The code
reports that integral at start-up against the $\alpha_A-\alpha_B$ that sets
the packet count and rescales neither side.  Since the correction described
next, the largest disagreement over the three channels is $0.01\%$.  It was
$3.44\%$ before, all of it in He\,\textsc{i}, and that number is what
exposed the defect.

\paragraph{The $\alpha_1$ source-mixing defect.}  \code{recomb\_mod} took
$\alpha_A$ from the Badnell total, $\alpha_1$ from the \emph{level-resolved}
Mao \& Kaastra fit, and then defined $\alpha_B = \alpha_A - \alpha_1$.  That
subtraction crosses two unrelated determinations.  Nothing in the
construction forces the Mao $n{=}1$ value to be the ground term of the
Badnell total it is subtracted from, and whatever separates them lands
entirely in $\alpha_B$ --- and, through $\alpha_1$, in the diffuse packet
count as well.

The cross-section side is not the error.  MoCHII's \code{sigma\_HeI}
reproduces Verner's own \code{phfit2.f} to a ratio of $1.00000$ at every
energy over $24.6$--$105$~eV, and its threshold value, $7.436$~Mb, sits
$0.5\%$ above the measured $7.40\pm0.07$~Mb of Samson et al.\ (1994),
inside the experimental uncertainty.  Nor is the start-up comparison between
unlike quantities: because $\alpha_B = \alpha_A - \alpha_1$, the ``rates''
side of it \emph{is} $\alpha_1$(He\,\textsc{i}) itself, so the $3.44\%$
separated two determinations of one number.

\paragraph{Cloudy settles which side is right.}  Cloudy c25.00 derives its
H-like and He-like recombination level by level from its \emph{own}
photoionization cross sections through the exact Milne relation
(\code{source/iso\_radiative\_recomb.cpp}), caching the result in
\code{data/h\_iso\_recomb.dat} and \code{data/he\_iso\_recomb.dat}.  In those
files row 0 of each element is the ground level and the last row is the
total --- verified on hydrogen, where row 0 is $1.5840{\times}10^{-13}$ and
the last row $4.1642{\times}10^{-13}$ against the known
$\alpha_{1s} = 1.58{\times}10^{-13}$ and
$\alpha_A = 4.18{\times}10^{-13}\,\cmc$.  At $10^4$~K, in
$10^{-13}\,\cmc$:
\begin{center}\small
\begin{tabular}{@{}lrrrrr@{}}
\toprule
& \multicolumn{2}{c}{$\alpha_A$} & \multicolumn{3}{c}{$\alpha_1$} \\
\cmidrule(r){2-3}\cmidrule(l){4-6}
channel & MoCHII & Cloudy & Mao \& Kaastra & Milne integral & Cloudy \\
\midrule
H\,\textsc{i}   & 4.1933 & 4.1642 & 1.5873          & 1.5790 & 1.5840 \\
He\,\textsc{i}  & 4.3719 & 4.3583 & \textbf{1.5650} & 1.6086 & 1.6032 \\
He\,\textsc{ii} & 21.895 & 21.817 & 6.5366          & 6.5116 & 6.5145 \\
\bottomrule
\end{tabular}
\end{center}
Badnell's totals are accurate to $0.3$--$0.7\%$ on all three channels, and
the Mao $n{=}1$ values to $0.2$--$0.3\%$ for H\,\textsc{i} and
He\,\textsc{ii} --- but He\,\textsc{i} is $2.38\%$ low, and the subtraction
propagates that into $\alpha_B$(He\,\textsc{i}) as $+1.88\%$.  MoCHII's own
Milne integral matches Cloudy to $0.3\%$, $0.34\%$ and $0.04\%$.  The Milne
relation is exact, Cloudy applies it to a He$^0$ ground cross section that
agrees with Verner's to the $0.4\%$ implied by the $\alpha_1$ agreement, and
so \emph{the Milne route is the correct one and the level-resolved fit is
the low side}.

\paragraph{An earlier reading of the same numbers, withdrawn.}  This
document previously concluded the opposite --- that the Milne integral was
the outlier --- from the third determination available at the time,
Benjamin, Skillman \& Smits (1999) with J.~S.\ Mathis as carried by
K.~Wood's \code{ionize},
$\alpha_1({\rm He\,\textsc{i}}) = 1.54{\times}10^{-13}(T_4)^{-0.486}$
$\cmc$:
\begin{center}\small
\begin{tabular}{lrrrr}
\toprule
$T$ & BSS99 $+$ JSM & Mao \& Kaastra & Milne integral & Cloudy c25.00 \\
\midrule
$5000$~K       & 2.157 & 2.184 ($+1.3\%$) & 2.259 ($+4.8\%$) & --- \\
$10^4$~K       & 1.540 & 1.565 ($+1.6\%$) & 1.609 ($+4.5\%$) & 1.603 ($+4.1\%$) \\
$2{\times}10^4$~K & 1.100 & 1.128 ($+2.6\%$) & 1.151 ($+4.7\%$) & --- \\
\bottomrule
\end{tabular}
\end{center}
In $10^{-13}\,\cmc$, percentages against the BSS99 column.  The arithmetic
was right --- the Milne integral does sit $4.5$--$4.8\%$ above BSS99 and
$2.8\%$ above Mao while the two fits cluster within $2.6\%$ of each other
--- but the inference was a majority vote taken on the mutual closeness of
two fits, which is not evidence that either reproduces the cross sections in
use.  Cloudy, computing the same quantity level by level from its own cross
sections, lands within $0.34\%$ of the Milne integral.  The earlier
conclusion is withdrawn.

\paragraph{The correction.}  \code{tools/fitting/fit\_alpha1\_milne.py}
evaluates Eq.~(\ref{eq:milne-alpha1}) on the cross sections the transport
actually uses --- exact hydrogenic for H\,\textsc{i} and He\,\textsc{ii}
(\code{sigma\_hydrogenic}, $Z = 1$ and 2; note that \code{sigma\_HI} in
\code{src/photo\_xsec.f90} is the exact hydrogenic form, not VFKY96) and
VFKY96 for He\,\textsc{i} --- at 241 temperatures over $10^3$--$10^5$~K, and
refits the same seven coefficients of Eq.~(\ref{eq:mao}), so that no Fortran
structure changes:
\begin{center}\small
\begin{tabular}{@{}lp{112mm}@{}}
\toprule
function & $(a_0,\,b_0,\,c_0,\,a_1,\,b_1,\,a_2,\,b_2)$ \\
\midrule
\code{alpha1\_HII}   & $1.87007{\times}10^{-2}$, $1.29077$,
  $-2.17465{\times}10^{-3}$, $1.29189{\times}10^{1}$,
  $8.05060{\times}10^{-1}$, $8.44399{\times}10^{-2}$,
  $5.93940{\times}10^{-4}$ \\
\code{alpha1\_HeII}  & $3.26862{\times}10^{-3}$, $6.97361{\times}10^{-1}$,
  $3.14907{\times}10^{-4}$, $5.76341{\times}10^{1}$, $1.18534$,
  $2.58417{\times}10^{1}$, $9.90134{\times}10^{-1}$ \\
\code{alpha1\_HeIII} & $2.61525{\times}10^{-1}$, $1.39527$,
  $-3.22079{\times}10^{-4}$, $6.40708{\times}10^{1}$,
  $8.97582{\times}10^{-1}$, $5.01551{\times}10^{-1}$,
  $1.21533{\times}10^{-3}$ \\
\bottomrule
\end{tabular}
\end{center}
The largest $|{\rm fit}/{\rm Milne}-1|$ is $0.019\%$ (H\,\textsc{i}),
$0.001\%$ (He\,\textsc{i}) and $0.006\%$ (He\,\textsc{ii}) over
$10^3$--$10^5$~K, and $0.0075\%$, $0.0004\%$, $0.0021\%$ over the production
range $3{\times}10^3$--$3{\times}10^4$~K, so the functional form is not the
limiting approximation.  At $10^4$~K the fitted $\alpha_1$ is $1.5791$,
$1.6086$ and $6.5117{\times}10^{-13}\,\cmc$: within $0.35\%$ of Cloudy on
all three channels, with He\,\textsc{i} corrected from $-2.38\%$ to
$+0.34\%$.  The code's own cross-check confirms it, \code{milne\_setup}
reporting a largest $\alpha_1$ disagreement of $3.44\%$ before and $0.01\%$
after.

The gate is \code{tests/milne/check\_alpha1.f90}, wired into
\code{tests/milne/run\_check.sh}.  It re-implements the Milne quadrature
independently from \code{photo\_xsec}, checks the three channels at 81
temperatures over $10^3$--$10^5$~K to $10^{-3}$, checks
$\alpha_1 < \alpha_A$ everywhere (largest $\alpha_1/\alpha_A$ $0.560$,
$0.496$, $0.445$), and checks all three against the hardcoded Cloudy c25.00
values to $1\%$.  Reinstating the old Mao He\,\textsc{i} coefficients makes
it fail on both counts ($3.886{\times}10^{-2}$ from the quadrature, $2.383\%$
from Cloudy) and it passes again after reverting.

\paragraph{Effect on the benchmarks, graded against the prediction.}  What
was recorded before the runs: ``$\alpha_1$(He\,\textsc{i}) rises $2.4\%$, so
channel 2 emits more; $\alpha_B$(He\,\textsc{i}) falls $1.9\%$, so there is
less case-B recombination.  Both push the front outward, toward Cloudy, by
of order $0.015$--$0.02$~pc.''  Measured on the paired HII40/HII20
production runs (68 ranks, same Sobol seed, converged in 38 and 27
iterations as before):
\begin{center}\small
\begin{tabular}{@{}lrr@{}}
\toprule
 & HII40 & HII20 \\
\midrule
$\Delta r({\rm He^0}{=}0.5)$ & $+0.00811$~pc & $+0.00210$~pc \\
Cloudy residual & $-0.18561 \to -0.17750$~pc & --- \\
$\Delta L({\rm H}\beta)$ & $-0.3713\%$ & $-0.4159\%$ \\
mean $T_e(\xHI < 0.95)$ & $9053.4 \to 9051.7$~K & $8522.3 \to 8528.7$~K \\
\bottomrule
\end{tabular}
\end{center}
The \emph{sign} was right, outward and toward Cloudy on both benchmarks.
The \emph{magnitude} was about half the low end of the predicted range,
$+0.0081$ against $0.015$--$0.02$~pc, so the prediction was roughly a factor
two high.  \code{tests/continuous\_energy/check\_hii\_production.py} passes
on both, with energy closure $1.5{\times}10^{-16}$ and $0$.  The published
benchmark quantities move by less than the width of the published spread:
$\langle T[N_pN_e]\rangle$ from 8162 to 8169~K at HII40 and 6917 to 6925~K
at HII20, $T_{\rm inner}$ from 7535 to 7536 and 6936 to 6938~K, and $R_{\rm
out}$ from $8.89$ to $8.88{\times}10^{18}$~cm at HII20 with
$1.44{\times}10^{19}$~cm unchanged at HII40.

$L({\rm H}\beta)$ moves the \emph{wrong} way.  It falls $0.37$--$0.42\%$
($1.94 \to 1.93{\times}10^{37}$ and $4.70 \to 4.68{\times}10^{36}\,{\rm
erg\,s^{-1}}$), taking MoCHII from about $3\%$ below the published
$2.01$--$2.10{\times}10^{37}$ at HII40 to slightly further below.  The
likely route is channel 1: $\alpha_1$(H\,\textsc{i}) drops $0.52\%$ from the
Mao value to the Milne one, so the H\,\textsc{i} ground continuum recycles
marginally fewer H-ionizing photons.  This is not an argument against the
change --- the acceptance criterion is physical correctness, and an exact
relation was being violated --- but it is recorded rather than buried.

Two legacy options remain.  \code{'exponential'} is the historical
near-threshold approximation
\begin{equation}
 E_{\rm pkt}=E_{\rm th}+kT_e[-\ln(\max(u,\epsilon))],
 \label{eq:diffuse-energy}
\end{equation}
i.e.\ Eq.~(\ref{eq:milne-pdf}) with the $E^2\sigma_{\rm pi}$ weight dropped;
it is kept to reproduce earlier runs.  \code{'threshold'} emits at
$E_{\rm th}$ and is a sensitivity bracket, not a model.  The three
thresholds are 13.598, 24.587, and 54.416~eV, and every model is a sampled
photon energy, not a bin-center approximation.

\paragraph{Sampling weight.}  MoCHII emits equal-energy packets and fixes
their number from the channel's \emph{energy} luminosity.  The packet
energies must therefore follow the energy-luminosity density $E\,n(E)$,
not $n(E)$ itself.  With that pairing the represented photon rate is exactly
the physical one,
\begin{equation}
 N_{\rm pkt}L_{\rm pkt}\left<E^{-1}\right>_{En}
   = R\left<E\right>_n\frac{1}{\left<E\right>_n}=R ,
\end{equation}
whereas drawing from $n(E)$ while counting packets from the energy
luminosity inflates it by $\left<E\right>_n\left<E^{-1}\right>_n>1$
(1.002 at $8000$~K, 1.010 at $2\times10^4$~K).  The CDF is tabulated in
$x=(E-E_{\rm th})/kT$, interpolated in $\log T$ and inverted once, rather
than inverting two neighboring temperature rows and interpolating the
resulting energies.

\paragraph{Effect.}  Against the exponential approximation the converged
HII40 calculation moves $r({\rm He^0}=0.5)$ outward by $0.00063$~pc and
$r({\rm C^{2+}}=0.5)$ by $0.00052$~pc, closing $0.3\%$ of the residual
against Cloudy c25.  The two continua move in opposite directions because
each follows its own mean energy: the H\,\textsc{i} continuum softens
($14.460\rightarrow14.427$~eV at $10^4$~K) since the hydrogenic cross
section falls as $E^{-3}$, while the He\,\textsc{i} continuum hardens
($25.449\rightarrow25.460$~eV) because the Verner He\,\textsc{i} fit falls
more slowly near threshold.  The exact spectrum is therefore a correctness
fix, not the explanation of the Cloudy front residual.  The same holds for
the three He\,\textsc{i} cascade corrections described below: the two-photon
shape moves $r({\rm He^0}=0.5)$ by $-0.00030$~pc, the branch fractions by
$+0.00024$~pc, and the fate of the 584~\AA{} photon by $-0.00363$~pc.  All
four corrections to the emitted diffuse spectrum together move it
$-0.00306$~pc, i.e.\ $1.7\%$ of the $0.18$~pc residual and \emph{away} from
Cloudy rather than toward it.  Switching the whole fourth channel off moves
it $-0.103$~pc, so the channel's presence matters and its internal detail
does not.

\paragraph{Where the Cloudy front residual is not.}  Each of those four
corrections is independently correct --- three of them reproduce an external
reference implementation and the first an exact detailed-balance relation ---
and the acceptance criterion is physical correctness, not agreement with a
particular reference calculation, so a net displacement in the unhelpful
direction is not an argument against any of them.  What the sum does settle
is the diagnosis: \emph{the $0.18$~pc residual is not in what the diffuse
field emits.}  Four independent corrections, spanning the emitted spectrum's
shape, its branch weights and the fate of its strongest line, together
account for $1.7\%$ of the gap and with the wrong sign.

\paragraph{Bounding the remaining candidates.}  The quantity to be explained
is $r({\rm He^0}=0.5) = 4.12781$~pc against Cloudy c25.00's $4.30531$~pc,
i.e.\ $-0.17750$~pc after the $\alpha_1$ correction above.  The sensitivity
runs below were measured against the pre-correction production point,
$4.11970$~pc; the $0.008$~pc between the two baselines changes no lever, and
the raw front positions are quoted as measured.  The front's sensitivity to
the diffuse field is
calibrated two ways from the 584~\AA{} measurement above: differentially
(channel 4 lost $8.87\%$ of its energy, is $17.6\%$ of the diffuse field, and
moved the front $0.00363$~pc) and by switching that channel off entirely
($17.6\%$ of the field, $0.103$~pc), giving $0.233$ and $0.585$~pc per unit
fractional change in the total diffuse band luminosity.  Two further
sensitivity runs (HII40, 68 ranks, same Sobol stream, 38 iterations) scale
the He\,\textsc{i} atomic data by $3\%$: reducing
$\sigma_{\rm pi}$(He\,\textsc{i})
moves the front $+0.00799$~pc and reducing the He\,\textsc{i} recombination
rate moves it $+0.02410$~pc.  Every candidate then has a number attached to
it, and every number is too small:
\begin{center}\small
\begin{tabular}{@{}lll@{}}
\toprule
axis & bound measured & to reach $0.178$~pc \\
\midrule
what the diffuse field emits & four corrections, net $-0.00306$~pc
  & wrong sign \\
explicit transport vs.\ on-the-spot & 4.11970 to 4.18159~pc
  & the whole range is $0.062$~pc \\
total diffuse band luminosity & $0.233$--$0.585$~pc per unit fraction
  & $30$--$76\%$ \\
$\sigma_{\rm pi}$(He\,\textsc{i}) & $+0.00799$~pc per $-3\%$ ($0.27$~pc/unit)
  & ${\sim}70\%$ \\
$\alpha$(He\,\textsc{i}) & $+0.02410$~pc per $-3\%$ ($0.80$~pc/unit)
  & ${\sim}23\%$ \\
$\alpha_A$, $\alpha_1$, $\alpha_B$(He\,\textsc{i})
  & $0.3\%$, $0.34\%$, $1.9\%$ from Cloudy (\S\ref{sec:henetwork})
  & already there; code spread is $11\%$ \\
metal opacity & $0.103$~pc/unit upward, $0.615$ downward
  (\S\ref{sec:metalop}) & $+180\%$, and saturating \\
metal recombination continua & $0.29\%$ of channel 2 & ${\sim}0.001$~pc \\
He\,\textsc{i} excited-level ionization & $0.2$--$0.6\%$ extra, and below the
  band floor & structurally absent \\
dust & \code{dust\_model='none'} in the benchmark & excluded by construction \\
source photon budget & $Q$ ratios $1.00005$ and $1.00003$ & $14\%$ in
  $Q({\rm He^0})$ \\
Cloudy-side resolution & median zone width $0.0106$~pc at the front
  & $17\times$ finer already \\
\bottomrule
\end{tabular}
\end{center}
The transport entry is the decisive one, because it was the leading candidate:
switching the whole diffuse treatment between its extremes --- explicit
transport at $4.11970$~pc and the on-the-spot approximation at $4.18159$~pc
--- spans only $0.062$~pc, and Cloudy sits $0.124$~pc beyond the more
favorable end of that range.  No choice about how diffuse photons are
transported can reach it.  The photon budget is verified against the same
40-kK blackbody Cloudy was given: $Q({>}13.598~{\rm eV}) =
4.2656{\times}10^{49}$~s$^{-1}$ against $4.2653{\times}10^{49}$ and
$Q({>}24.587~{\rm eV}) = 4.6082{\times}10^{48}$ against
$4.6081{\times}10^{48}$, ratios $1.00005$ and $1.00003$.  Cloudy resolves the
front with 40 zones between 4.0 and 4.6~pc, median width $0.0106$~pc and
maximum $0.039$~pc, so the reference position is not a resolution artifact
either.

\paragraph{The He-ionizing photon budget.}  Where the sensitivity lives is
visible in the budget itself.  Integrated over the He$^+$ zone of the
converged HII40 model:
\begin{center}\small
\begin{tabular}{@{}lr@{}}
\toprule
HII40 & s$^{-1}$ \\
\midrule
stellar $Q({>}24.587$~eV) & $4.6082{\times}10^{48}$ \\
diffuse He\,\textsc{i} ground continuum, recycled & $1.6135{\times}10^{48}$ \\
total He-ionizing supply & $6.2217{\times}10^{48}$ \\
He$^0$ photoionizations $=$ case-A recombinations & $4.7008{\times}10^{48}$ \\
absorbed by H$^0$ or metals, or escaped
  & $1.5210{\times}10^{48} = 24.4\%$ of the supply \\
\bottomrule
\end{tabular}
\end{center}
Nearly a quarter of the He-ionizing supply is lost to absorbers other than
He$^0$, and halving that loss fraction alone would move the front
$0.211$~pc.  The H$^0$/He$^0$ competition is therefore the one place where a
few-percent difference has the leverage to matter, which is consistent with
the radial signature below and with the metal-opacity mechanism of
\S\ref{sec:metalop}.  The budget also closes geometrically: with the leaf
volume computed correctly the summed gas volume matches the geometric shell
$0.972$--$4.729$~pc to $0.25\%$.  (\code{LeafSize} in the output is the leaf
\emph{full} width, and the residual analysis script had used
$(2\,\mathrm{size})^3$, eight times the cell volume.  Because it enters only
as a relative weight within a shell, every front position reported from it is
unaffected and re-running after the fix reproduces all of them identically;
the expression was nevertheless wrong and is now $\mathrm{size}^3$.  The
table generator, the Python readers and the g1/g4 gates all used the
convention correctly and were never affected.)

\paragraph{The radial signature.}  What the residual does have is a shape.
The hydrogen structure agrees with Cloudy throughout the ionized zone; only
the helium structure drifts, and it drifts monotonically outward:
\begin{center}
\begin{tabular}{lrrrrrr}
\toprule
$r$ [pc] & 1.6 & 2.0 & 2.5 & 3.0 & 3.5 & 3.8 \\
\midrule
$x({\rm He^0})$, MoCHII/Cloudy & 1.007 & 1.015 & 1.031 & 1.059 & 1.109
  & 1.196 \\
$x({\rm H^0})$, MoCHII/Cloudy  & 0.981 & 0.984 & 0.989 & 1.000 & 1.005
  & 1.009 \\
\bottomrule
\end{tabular}
\end{center}
A normalization difference in a rate or a cross section would be
radius-independent; this one compounds with depth.  That is what a $0.7\%$
seed at 1.6~pc must do when He$^0$ carries almost all of the opacity above
24.6~eV: at 30~eV,
$\kappa({\rm He^0})/[\kappa({\rm He^0})+\kappa({\rm H^0})]$ rises from
$0.843$ at 1.6~pc to $0.956$ at 4.0~pc, so any excess He$^0$ shortens the
mean free path of the very photons that would have removed it.  The
signature therefore identifies a small helium-specific seed, amplified
geometrically, rather than a large error anywhere.

Closing $0.178$~pc is $4.1\%$ in radius and $13\%$ in volume, which asks for
$Q({\rm He^0})$ higher by ${\sim}14\%$ or $\alpha_B$(He\,\textsc{i}) lower by
${\sim}20\%$.  $Q({\rm He^0})$ is verified to $0.003\%$, and MoCHII's
$\alpha_A$, $\alpha_1$ and $\alpha_B$(He\,\textsc{i}) now sit within $0.3\%$,
$0.34\%$ and $1.9\%$ of Cloudy's own level-resolved values, with no
published determination anywhere near $20\%$ away.  \emph{The residual is
therefore unexplained}, and what has been achieved is the systematic closing
of the space of simple explanations: it is not in what the diffuse field
emits, not in how those photons are transported, not in the He\,\textsc{i}
photoionization cross section, not in the He\,\textsc{i} recombination rate,
not in the metal opacity, not in the metal recombination continua, not in
He\,\textsc{i} excited-level ionization, not in dust, not in the source
photon budget, and not in the reference calculation's resolution.

The two candidates named in earlier revisions have both been examined and
neither survives as a single parameter.  The helium ionization
\emph{network} has been compared code to code (\S\ref{sec:henetwork}): every
code sums the full cascade, MoCHII's $\alpha_A$(He\,\textsc{i}) agrees with
Cloudy's explicit level sum to $0.3\%$, and the He\,\textsc{i} excited levels
that a resolved model atom carries are worth $0.2$--$0.6\%$ of extra
ionization.  The metal opacity has been measured
(\S\ref{sec:metalop}): its upward derivative is $0.103$~pc per unit, so
closing the gap would need $+180\%$ against cross sections that reproduce
Verner's to $10^{-6}$ and an Opacity Project comparison that runs $2$--$14\%$
the other way.

What remains is the one hypothesis that no single-parameter experiment can
reproduce: several differences of a few percent each --- in the He\,\textsc{i}
recombination rate, in the shape of $\sigma_{\rm pi}$(He\,\textsc{i}), in the
metal opacity of the $21$--$40$~eV window, in Cloudy-side modeling choices
not yet identified --- acting together through the He$^0$ optical depth,
which compounds them with depth exactly as the radial signature shows.  Each
is individually below its own measured lever; their product over $24.4\%$ of
the He-ionizing supply lost to competing absorbers is not bounded by any of
the experiments above.  Testing it requires varying several ingredients
jointly, not one at a time, and that is the open item.

All diffuse packets carry the same luminosity as direct packets,
$L_{\rm pkt}=L_{\rm band}/N_{\rm pkt}$, so their number is
$N_{\rm dif}=\mathrm{nint}(L_{\rm dif}/L_{\rm pkt})$.  A packet first
selects a leaf from the CDF of $\sum_{\rm ch}L_{{\rm ch},l}$, then a channel
from that leaf's four-channel luminosity weights.  Its emission point is
uniform in the selected cubic leaf and its direction is isotropic.  It is
then transported with the same exact-energy edge walk and path estimator as
a stellar packet.  Hence diffuse luminosity is represented by packet count,
not by a different packet weight.  At HII40 the diffuse band luminosity
converges to $6.8{\times}10^{38}$~erg/s $=0.50\,L_\star$.

For output bookkeeping, the energy from Eq.~(\ref{eq:diffuse-energy}) is
mapped by binary search of \code{ion\_eedge} (\code{ion\_bin\_of}).  This
works for both threshold-aligned and logarithmic diagnostic grids; an energy
on an edge falls in the bin above it.  In continuous mode this assignment
does not modify $E_{\rm pkt}$ or its cross sections.  In grouped mode only,
the representative energy of that selected bin is used for transport.  The
binary search also corrects the former closed-form log-grid mapping, which
could place a 24.6-eV He\,\textsc{i} recombination photon below the
He\,\textsc{i} threshold on an aligned grid.

\paragraph{He\,\textsc{i} excited-level recombination
(\code{par\%hei\_diffuse}).}
A fourth channel adds the He\,\textsc{i} \emph{excited-level} recombination
radiation (previously those photons simply vanished).  Every case-B
He\,\textsc{ii} recombination reaches $1^1{\rm S}$ through exactly one
$n = 2$ term, so what the channel needs is not three absolute rates but
three \emph{branch fractions}; the absolute rate is already fixed by the
$\alpha_B({\rm He\,II})$ that sets the packet count.  MoCHII takes all three
from one internally consistent set of effective recombination coefficients,
\begin{equation}
\begin{aligned}
 \alpha_{\rm eff}(2^3{\rm S}) &= 2.10{\times}10^{-13}\,(T_4)^{-0.381}, \\
 \alpha_{\rm eff}(2^1{\rm S}) &= 2.06{\times}10^{-14}\,(T_4)^{-0.451}, \\
 \alpha_{\rm eff}(2^1{\rm P}) &= 4.17{\times}10^{-14}\,(T_4)^{-0.695},
\end{aligned}
 \label{eq:hei-branch}
\end{equation}
in $\cmc$ --- Benjamin, Skillman \& Smits (1999) Table~1 with J.~S.\ Mathis,
in the form carried by K.~Wood's \code{ionize} --- and normalizes by their
sum.  That all three come from one source matters more than any individual
value: a fraction assembled from fits with unrelated temperature dependences
is not a fraction of anything.  Normalizing by the sum is the source's own
choice, made for the reason it states, that the fitting formulae carry
different temperature dependences.  Outside $3{\times}10^3$--$2{\times}10^4$~K
the fits are clamped, not extrapolated.

The closure
\begin{equation}
 \alpha_{\rm eff}(2^3{\rm S}) + \alpha_{\rm eff}(2^1{\rm S})
   + \alpha_{\rm eff}(2^1{\rm P}) = \alpha_B({\rm He\,II})
 \label{eq:hei-closure}
\end{equation}
is physically true, but MoCHII does \emph{not} use it to derive a branch.
Its two sides are independent atomic data: the three fits of
Eq.~(\ref{eq:hei-branch}) account for $0.90\,\alpha_B$ at 8000~K and
$0.97\,\alpha_B$ at $10^4$~K, and the ratio is reported at start-up rather
than absorbed.  Forcing it onto $\alpha_B$ would dump the entire mismatch
between two data sets into whichever branch was solved for last.  The
resulting fractions, as the code prints them, are
\begin{center}
\begin{tabular}{lcccc}
\toprule
$T$ [K] & $2^3$S & $2^1$S & $2^1$P & $\sum\alpha_{\rm eff}/\alpha_B$ \\
\midrule
$5{\times}10^3$   & 0.74084 & 0.07629 & 0.18288 & 0.767 \\
$8{\times}10^3$   & 0.76183 & 0.07591 & 0.16226 & 0.896 \\
$10^4$            & 0.77121 & 0.07565 & 0.15314 & 0.970 \\
$1.5{\times}10^4$ & 0.78729 & 0.07507 & 0.13764 & 1.134 \\
\midrule
statistical & 0.75000 & 0.08333 & 0.16667 & --- \\
\bottomrule
\end{tabular}
\end{center}
against the AGN3 low-density statistical weights of the last row, which the
code used before.  The correction is a moderate temperature-dependent shift,
not a reordering.

Two provenance points belong with these numbers.  First,
$2.06{\times}10^{-14} + 4.17{\times}10^{-14} = 6.23{\times}10^{-14}$
exactly, so the $6.23{\times}10^{-14}(T_4)^{-0.827}$ that \code{nebcont\_mod}
took from MOCASSIN for the He\,\textsc{i} two-photon continuum
(Section~\ref{sec:twophoton}) is the \emph{singlet total}, not
$\alpha_{\rm eff}(2^1{\rm S})$.  An earlier version of this section read it
as the latter and closed $2^1$P against $\alpha_B$, obtaining
$0.746 : 0.224 : 0.031$ together with an argument that case-B trapping
inverts the singlet split; both the numbers and the argument are withdrawn.
Second, the $2^3$S coefficient here is deliberately \emph{not} shared with
the \code{ALPHA3} of \code{hei\_metastable\_mod}, which is
$2.10{\times}10^{-13}(T_4)^{-0.778}$ (Oklopcic \& Hirata 2018): the same
prefactor with a steeper exponent, $9\%$ higher at 8000~K.  That one is an
absolute production rate for the metastable population; this one is a member
of a set whose ratios are the branch fractions, and mixing sources would
break the internal consistency those ratios depend on.

The $2^1$S photons follow the tabulated Drake, Victor \& Dalgarno (1969)
distribution (\code{data/HeI2phot.dat}, the same table the nebular continuum
uses) rather than a flat draw.  That shape fixes the H-ionizing yield at
0.55644 photons and 8.9646~eV per decay, with a mean in-band energy of
16.110~eV, against the 0.56 and 17.109~eV of the flat surrogate.  The two
corrections are individually small and act in opposite directions.  Per
case-B He\,\textsc{ii} recombination the channel emits
\begin{center}
\begin{tabular}{lcc}
\toprule
 & H-ionizing energy [eV] & photons \\
\midrule
statistical weights, flat two-photon & 19.200 & 0.9634 \\
statistical weights, Drake shape     & 19.149 & 0.9633 \\
Eq.~(\ref{eq:hei-branch}), Drake shape, 8000~K & 19.223 & 0.9663 \\
\bottomrule
\end{tabular}
\end{center}
i.e.\ $-0.27\%$ in energy from the two-photon shape, then $+0.39\%$ in
energy and $+0.34\%$ in photons from the branch fractions, for a net
$+0.12\%$ and $+0.30\%$ against the statistical weights.  In transport (68
ranks, same Sobol stream, branch fractions against the Drake-only state) the
He\,\textsc{i} front moves by $+0.00024$~pc at HII40 and $+0.00006$~pc at
HII20, $r({\rm C^{2+}}=0.5)$ by $+0.00015$ and $0.00000$~pc, and
$L({\rm H}\beta)$ by $+0.0121\%$ and $+0.0027\%$; both benchmarks still pass
\code{tests/continuous\_energy/check\_hii\_production.py}.  Like the Milne
ground continua, this is a correctness fix and not the explanation of the
$0.18$~pc Cloudy residual.

\paragraph{The fate of the 584~\AA{} photon.}  The $2^1$P branch decays to
the ground state on He\,\textsc{i} Ly$\alpha$, 584~\AA{} $=21.22$~eV, which
is both an H-ionizing photon and a resonance photon in a medium optically
thick to it in He$^0$.  It therefore scatters many times before anything
irreversible happens, and two irreversible outcomes compete: photoionization
of a hydrogen atom (depositing $21.22-13.60 = 7.62$~eV of heat), or
collisional transfer $2^1{\rm P}\rightarrow 2^1{\rm S}$ in the last He atom
that absorbed it, after which the decay is the two-photon continuum instead.
Which one wins is set by the H$^0$ and He$^0$ opacities the trapped photon
sees, so it is a local, density- and ionization-dependent branching and not a
constant.  MoCHII resolves it with the escape-probability closure of Wood,
Mathis \& Ercolano (2004), Section~3.3,
\begin{equation}
 p_{584} = \left[1 + \frac{77\,x({\rm He^0})}
                          {\sqrt{T}\;x({\rm H^0})}\right]^{-1},
 \label{eq:hei-584}
\end{equation}
the probability that the 584~\AA{} photon ends up ionizing hydrogen
(\code{hei\_584\_ionizes\_H} in \code{hei\_cascade\_mod}).  Its two limits are
the correct ones: $x({\rm H^0})\rightarrow0$ gives $p_{584}\rightarrow0$,
nothing left to ionize and everything converted, while
$x({\rm He^0})\rightarrow0$ gives $p_{584}\rightarrow1$, nothing left to
convert.  K.~Wood's \code{ionize} and CMacIonize carry the same closure in
two algebraically different arrangements, and the gate
\code{tests/hei\_cascade/check\_hei\_584.f90} reproduces both to
$2.2{\times}10^{-16}$ over a grid in $(T, x_{\rm H^0}, x_{\rm He^0})$.  The
four codes in the reference trees sit at three different points on the
question:
\begin{center}
\begin{tabular}{llll}
\toprule
code & cascade coefficients & fate of the 584~\AA{} photon & transported? \\
\midrule
Wood \code{ionize} & Eq.~(\ref{eq:hei-branch}), sum-normalized &
  $p_{584}$ ionizes H; $1-p_{584}$ converts & no, on the spot \\
CMacIonize & the same set and normalization & the same split &
  no, on the spot \\
MOCASSIN & singlet total $6.23{\times}10^{-14}(T_4)^{-0.827}$ &
  all converted ($p_{584}=0$) & emissivity only \\
MoCHII & Eq.~(\ref{eq:hei-branch}), sum-normalized &
  Eq.~(\ref{eq:hei-584}) split & yes, full transport \\
\bottomrule
\end{tabular}
\end{center}
CMacIonize is also where MoCHII's flat-surrogate constant $0.56$ came from,
as $2\times28\%$ of the two-photon distribution above the Lyman limit;
Eq.~(\ref{eq:hei-branch}) and the Drake table replaced both of those
constants above.  MOCASSIN does not resolve the competition at all: its
comment states the assumption that all He\,\textsc{i} singlets end in
$2^1$S, and the $6.23{\times}10^{-14}$ singlet total is that
$p_{584}\rightarrow 0$ limit written as one coefficient.  MoCHII is now the
only one of the four that applies the branching and then \emph{transports}
what it produces: Wood and CMacIonize destroy the 584~\AA{} photon where it
was created, so the deposition of whatever the branch yields is assumed
local, while MoCHII carries the branch product through the grid with the
same edge walk as any other packet and lets the grid decide where it is
absorbed.

Because $p_{584}$ depends only on the emitting leaf's own $T$, $x({\rm H^0})$
and $x({\rm He^0})$, the split folds into the branch weights of
\code{hei\_branch\_energies} at build time and needs no random variate of its
own, so the nine-coordinate Sobol assignment below is unchanged:
\begin{equation}
\begin{aligned}
 e(2^3{\rm S}) &= f_{2^3{\rm S}}\times 19.82~{\rm eV}, \\
 e(2^1{\rm P}) &= f_{2^1{\rm P}}\,p_{584}\times 21.22~{\rm eV}, \\
 e(2^1{\rm S}) &= \left[f_{2^1{\rm S}}
                  + f_{2^1{\rm P}}(1-p_{584})\right]
                  \times 8.9646~{\rm eV}.
\end{aligned}
\end{equation}
A converted decay radiates the two-photon pair, $20.62$~eV, and not
$21.22$~eV; the difference of $0.60$~eV leaves as $2^1{\rm P}\rightarrow
2^1{\rm S}$ at 2.06~$\mu$m, which is non-ionizing and therefore correctly
absent from the band.  \code{gamma\_2q\_HeI} uses the same split for the
He\,\textsc{i} two-photon nebular continuum,
$\alpha_{2q} = \alpha_{\rm eff}(2^1{\rm S}) + [1-p_{584}]\,
\alpha_{\rm eff}(2^1{\rm P})$ (Section~\ref{sec:twophoton}), which closes an
internal inconsistency: \code{nebcont\_mod} had inherited MOCASSIN's singlet
total, the $p_{584}\rightarrow0$ limit, while \code{diffuse\_mod} transported
every $2^1$P decay as a 21.22-eV packet, the $p_{584}\rightarrow1$ limit, so
the two modules disagreed about the fate of the same decay.  The domain of
validity is recorded at the code site: this is an escape-probability closure
for a 584~\AA{} line assumed optically thick, static and dust free.  Dust
competes for the trapped photon and is not in the closure --- a small
correction at the Lexington dust content, not a small one at higher
dust-to-gas ratios.

\paragraph{Effect of the 584~\AA{} split.}  Weighted by the $2^1$P emission
rate on the converged grids, $\left<p_{584}\right> = 0.117$ at HII40 and
$0.0044$ at HII20: $88\%$ and $99.6\%$ of that branch converts to the
two-photon continuum, and the channel's output per case-B He\,\textsc{ii}
recombination falls from 19.22 to 17.47~eV at HII40 ($-9.1\%$) and from 19.24
to 17.12~eV at HII20 ($-11.0\%$), with the photon count down $6.6\%$ and
$7.9\%$.  In transport (68 ranks, same Sobol stream, against the branch-fraction
state):
\begin{center}
\begin{tabular}{lrr}
\toprule
 & HII40 & HII20 \\
\midrule
$\Delta r({\rm He^0}=0.5)$ [pc]  & $-0.00363$ & $-0.00022$ \\
$\Delta r({\rm C^{2+}}=0.5)$ [pc] & $-0.00332$ & $-0.00001$ \\
$\Delta r({\rm H^0}=0.1)$ [pc]   & $-0.01002$ & (density bounded) \\
$\Delta L({\rm H}\beta)$          & $-0.1997\%$ & $-0.0704\%$ \\
$\left<T_e\right>$ ($x_{\rm H^0}<0.95$) [K]
  & $9080.2\rightarrow9053.4$ & $8523.0\rightarrow8522.3$ \\
diffuse band luminosity           & $-1.56\%$ & $-0.22\%$ \\
\bottomrule
\end{tabular}
\end{center}
Both benchmarks still pass
\code{tests/continuous\_energy/check\_hii\_production.py}.  This is the
largest of the four corrections to the emitted diffuse spectrum --- six times
the Milne ground continua and twelve times the branch fractions --- and it
still moves the He\,\textsc{i} front only $2\%$ of the way to the Cloudy
residual, in the wrong direction.

Three predictions were recorded before these runs, and grading them is worth
keeping.  The sign was predicted correctly: losing channel-4 photons pulls
the fronts inward.  The magnitude, of order $0.01$~pc, was correct.  The
third, that HII20 would move more than HII40, was \emph{wrong} --- HII40
moved sixteen times further.  The prediction had been read off $p_{584}$
alone, which measures only what fraction of the $2^1$P branch converts;
front displacement instead scales with that branch's \emph{absolute}
contribution to the ionizing budget.  Inverting the measured luminosity drops
against the change per recombination puts channel 4 at $17.6\%$ of the HII40
diffuse field and $2.0\%$ of the HII20 one: a 20~kK spectrum makes few
photons above 24.6~eV, so He$^+$ is scarce and the channel fed by its
recombinations is small.  A larger converted fraction of a much smaller
channel moves the front less.  The lesson generalizes --- a local branching
probability predicts what happens per event, never how far a front moves;
that second factor is the channel's share of the budget, set by the hardness
of the source.

With \code{launch\_sequence='sobol'}, diffuse packets use a separate
Owen-scrambled Sobol stream so that their changing count does not disturb
the direct-source sequence.  Its fixed nine coordinates are: one for the
leaf CDF; three for the position within that leaf; two for $(\mu,\phi)$;
one for the channel CDF; one for the ground-continuum energy draw (or
the excited-He branch); and one for the two-photon energy,
drawn from the tabulated Drake shape, when that branch is selected.  Coordinates that are irrelevant to a selected channel
are still reserved, preserving the global packet-indexed layout across
iterations and MPI task counts.

\subsection{Solution-driven I-front re-refinement}
\label{sec:refine}
At iteration \code{refine\_iter} the octree is rebuilt from the current
solution: an old leaf is flagged as front when
$\epsilon < \xHI < 1{-}\epsilon$ or a face neighbor differs by more than
\code{refine\_dx}; the new tree carries level \code{refine\_lmax}
wherever a cell overlaps a front-flagged leaf and \code{refine\_lbase}
elsewhere.  The gas state and \code{rhokap} map over by position (each
new leaf samples the old leaf containing its center), and every
dependent array is recreated (the band opacity, tallies, metal
$\Gamma$ blocks, diffuse tables, rate arrays --- the last was the one
initially missed, an allocate-once array that faulted when
$n_{\rm leaf}$ grew).  All ranks rebuild identically from the shared
state, and the old grid's shared-memory windows are \emph{recycled}:
\code{destroy\_shared\_mem} (the collective counterpart of
\code{create\_shared\_mem}) frees the 18 old windows right after the
position sampling --- the old tree's last reader --- so repeated
re-refinement neither leaks nor holds old and new grids in memory
simultaneously.

\subsection{Dust in the ionizing band}
\label{sec:dust}
The three parts of the dust coupling, now complete:
\begin{itemize}
\item \textbf{Absorption competition.}  With \code{par\%ion\_add\_dust},
  each band bin adds
  $\kappa_{\rm d}(b,{\rm leaf}) = \rho\kappa_{\rm leaf}\,s_{\rm abs}(b)$
  where $s_{\rm abs}(b) = (1{-}a(E_b))\,C_{\rm ext}(E_b)/C_{\rm ext}
  (\lambda_{\rm ref})$ comes from a kext table covering the EUV (the
  WD01/D03 $R_V{=}3.1$ table, 1\,\AA--$10^4\,\mu$m:
  $C_{\rm abs}/{\rm H} = 1.83{\times}10^{-21}$ at 13.6~eV falling to
  $3.96{\times}10^{-22}\,{\rm cm^2}$ at 100~eV).  A trap worth
  recording: the D03 table is \emph{not monotonic} in wavelength (an
  out-of-order seam) --- without sorting, the linear interpolation
  search silently returned $s_{\rm abs}(100\,{\rm eV}) = 0.903$ instead
  of $0.832$ and shifted the dusty Str\"omgren radius by $+1.6\%$.
\item \textbf{Scattering.}  With \code{par\%ion\_dust\_scatter} the
  scattering coefficient joins the extinction
  ($s_{\rm sca}(b) = a\,C_{\rm ext}/C_{\rm ext,ref}$; in-band
  Henyey--Greenstein asymmetry $g(E_b) = 0.67$--$0.99$ from the table
  $\langle\cos\rangle$ column) and interactions are sampled: after the
  analytic zero-variance direct walk, a first interaction is forced
  along the ray with weight $w_1 = 1 - e^{-\tau_{\rm edge}}$, the
  packet keeps the albedo fraction
  $a_{\rm eff} = \kappa_{\rm sca,d}/\kappa_{\rm ext}$ at each
  interaction, scatters into an HG direction, and continues with
  unforced flights (whose segments tally into $J_b$) under Russian
  roulette ($w_{\min} = 10^{-4}$).  Stellar and diffuse packets share
  the walk; gas ionization, gas heating, and grain heating all derive
  from the same tally, so the energy split is bookkept by construction.
\item \textbf{Live ionization-dependent dust with a PAH split.}
  \code{dust\_model='laursen09\_live'} recomputes the dust density from
  the \emph{computed} $\xHI$ every iteration,
  $n_{\rm d} \propto (Z/Z_{\rm ref})\,n_{\rm H}\big[(1{-}f_{\rm PAH})
  (\xHI + f_{\rm ion}x_{\rm HII}) + f_{\rm PAH}(\xHI +
  f_{\rm ion,PAH}\,x_{\rm HII})\big]$, the PAH share carrying its own
  ionized-gas survival ($f_{\rm ion,PAH} = 0$ default: PAHs destroyed
  in ionized gas) --- the generalization including its
  PAH-specific option.
\item \textbf{Grain heating and the equilibrium temperature.}
  \code{Heat\_dust(leaf)} $= \sum_b [4\pi J\Delta\nu]_b\,
  \kappa_{\rm abs,d}$ [erg\,s$^{-1}$\,cm$^{-3}$] is consumed by
  \code{dust\_temp\_mod}: $f_{\rm emit}(T) = 4\pi\int s_{\rm abs}(\nu)
  B_\nu(T)\,d\nu$ is precomputed from the same kext table (log-$T$
  grid, 3.2--3200~K) and $T_{\rm d}({\rm leaf})$ solves
  ${\rm Heat\_dust} = (\rho\kappa/d_{\rm cm})f_{\rm emit}(T_{\rm d})$;
  outputs are the $T_{\rm d}$ leaf map and the grid-integrated
  modified-blackbody IR spectrum (1--1000\,$\mu$m).  FUV grain heating
  is transported by the \code{par\%add\_fuv} band extension (Section
  2.1): gas cross sections vanish below their thresholds automatically,
  so no new transport code is involved, and the two-segment grid keeps
  the ionizing bins at full resolution.  This is the
  single equilibrium \emph{mixture} temperature; the size- and
  material-resolved treatment is the SEDust coupling below.
\end{itemize}

\subsection{Stochastic dust emission (SEDust coupling)}
\label{sec:sedust}
With \code{par\%dust\_sed}, \code{sedust\_mod} (adapted from the MoCafe
\code{dustemis\_mod}) computes the grid-integrated dust SED with
size-resolved stochastic heating and PAH bands, using the in-tree
SEDust library (built from source under \code{SEDust/sed} by
\code{build\_lib.sh}; astrodust, DL07, and Zubko grain models).  The
coupling follows the MoCafe convention that makes the normalization
exact by construction:
\begin{itemize}
\item \textbf{Shape from the local field.}  For each leaf, the band
  tally $J_b$ is converted to $J_\lambda$ on the SEDust wavelength grid
  (1129 points, 0.0912--$4{\times}10^4\,\mu$m for astrodust).  The band
  reaches into the EUV, below the SEDust optics grid; that energy is
  deposited into the shortest SEDust bin, conserving the absorbed power
  (the hardness of the stochastic temperature spikes is slightly
  underestimated there --- a documented approximation).  SEDust then
  returns the emission spectrum of the grain population in this field:
  equilibrium large grains plus the full temperature-distribution solve
  for small grains and PAHs.
\item \textbf{Normalization from the transport.}  The SEDust spectrum
  is used only as a shape: each leaf's emission is normalized to its
  absorbed power ${\rm Heat\_dust}\cdot V$ from the Monte Carlo tally
  (radiative equilibrium).  H\,\textsc{ii} regions are optically thin
  to the re-emitted IR, so no Lucy re-iteration is needed and the
  grid-integrated $L_\lambda$ is accumulated directly (an
  \code{MPI\_ALLREDUCE} over leaves distributed round-robin).
  $\int L_\lambda\,d\lambda = \sum {\rm Heat\_dust}\,V$ holds by
  construction and is written into the \code{\_dustsed.txt} header.
\item \textbf{Cost.}  The exact stochastic solve costs
  $\sim$0.1--0.2~s for each leaf with nonzero grain heating; it runs
  once, at output time, after convergence (the dusty Str\"omgren
  smoke test: $\sim$1.9~hours on 8 ranks for a $64^3$ grid ---
  the solve, not the transport, dominates).
\item \textbf{Smoke result.}  Dusty Str\"omgren at $4{\times}10^4$~K:
  $\int L_\lambda d\lambda = 4.296{\times}10^{38}$~erg/s $= \sum
  {\rm Heat\_dust}\,V$ to 4 digits; the SED shows the PAH bands at
  3.3/6.2/7.7/8.6/11.3/12.7/17~$\mu$m over a warm 20--60~$\mu$m
  plateau with the FIR falloff --- the canonical stochastically
  heated spectrum for a hot compact field.
\item \textbf{Live PAH share} (\code{par\%sed\_pah\_live}):
  \code{dust\_emission}'s channel output (astrodust: 'AD' + 'PAH')
  lets each leaf's spectrum be assembled with the PAH channel weighted
  by the leaf's PAH survival $(x_{\rm HI} + f_{\rm ion,pah}
  x_{\rm HII})$, divided under \code{laursen09\_live} by the dust
  survival already carried by \code{rhokap}; the
  ${\rm Heat\_dust}\cdot V$ normalization keeps energy closure (the
  PAH energy shifts to the large grains).  Smoke: the PAH bands drop
  to 0.26 (7.7~$\mu$m) and 0.15 (11.3~$\mu$m) of the
  PAH-everywhere reference, the FIR rises by 1.06--1.36$\times$, and
  the 8-$\mu$m map turns from centrally peaked into a \emph{ring} at
  the ionization front --- the classic Spitzer/JWST bubble morphology.
\end{itemize}

\subsection{Peel-off imaging}
\label{sec:peel}
With \code{par\%ion\_peel}, one extra transport pass runs on the
converged state with observer peeling (the MoCafe peeling-off
convention, adapted to the band): every emitted packet --- stellar and
diffuse --- peels a direct contribution
$L_{\rm pkt}w\,e^{-\tau}/(4\pi r^2)$ toward each observer, and, when
dust scattering is sampled, every interaction peels a scattered
contribution with the Henyey--Greenstein phase function of the
packet's bin (after the albedo weight, before the new direction is
sampled).  The optical depth to the box edge along the observer
direction is integrated through \code{kap\_ion} at the packet's bin by
a tally-free walk (tallying the virtual peel ray would double count),
so the images carry exactly the transport's extinction: gas, dust
absorption, and dust scattering.  Images are TAN projections on the
MoCafe observer geometry (\code{obsx/y/z} or angle lists, rotation
matrix, auto-sized pixel scale), written as surface brightness
[erg\,s$^{-1}$\,cm$^{-2}$\,sr$^{-1}$] split into direct/scattered and
ionizing/FUV channels.  The band tally is rebuilt by the imaging pass,
so the written rates reflect the same field the images carry.  With \code{par\%peel\_bins} the
channel axis becomes the band bin --- \code{direct\_cube}/
\code{scatt\_cube} $(n_x, n_y, n_\nu)$ --- for hardness maps and
synthetic EUV/FUV photometry (the band-integrated channel images are
still written, derived from the cube; the gate shows the expected
color: the FUV scattered halo shrinks from a 15.5-px to a 12.0-px
half-light radius between 6.2 and 13.2~eV).

\paragraph{External-field peel.}  The direct peel of an external
packet (Section~\ref{sec:sources}) replaces the isotropic
$1/(4\pi r^2)$ with its Lambert emission density about the
entry-surface inward normal $\hat n_{\rm in}$,
$p(\Omega) = (\hat k \cdot \hat n_{\rm in})/\pi$, so the contribution
is $L_{\rm pkt}\,w\,\max(\hat k_{\rm obs}\cdot\hat n_{\rm in}, 0)\,
e^{-\tau}/(\pi r^2)$: only the far-side surface (inward normal toward
the observer) contributes, and the image is the background field seen
through the medium --- the cloud in silhouette.  In the optically thin
limit the image-integrated flux is $F = J A_{\rm proj}/d^2$, and with
\code{save\_direc0} the ratio \code{direc0/direct}$\,= e^{+\tau}$ is
again the line-of-sight optical-depth map.  The scattered peel is
origin-agnostic and unchanged.

For the \emph{dust} bands no transport is needed --- the nebula is
optically thin to its own IR --- so \code{par\%dust\_emis\_bands}
instead stores SEDust band emissivities of each leaf
(\code{\_dustemis}, the same block layout as the \code{\_emis} file)
and the Python map maker projects them into flux-conserving dust-band
images exactly like line maps.

\subsection{Cartesian (car) grids}
\label{sec:uniform}
\code{par\%grid\_type = 'car'} builds a single-level Cartesian grid and
stores its leaves in raster order: leaf $=$ cell $=$
$1 + i_x + n_x(i_y + n_y i_z)$.  No tree topology, no neighbor
table, and no geometry arrays are built.  (\code{'uniform'} is accepted
as a backward-compatible alias.)  The leaf list comes either from a file
(\code{par\%amr\_file}, single level) or --- with \code{amr\_file} empty
--- directly from the namelist: \code{par\%nx/ny/nz} cells over the box
$[-x_{\max},x_{\max}]{\times}\dots$ (cubic cells required).

\paragraph{Uniform density and sphere/shell geometry.}  A gas-density
model is applied on the leaf centers before either grid (\code{'car'} or
\code{'amr'}) is built: \code{par\%nH\_const} overrides every leaf's
$n_{\rm H}$ with one uniform value, and \code{par\%rmax}/\code{rmin} zero
$n_{\rm H}$ outside a sphere of radius $r_{\max}$ / inside $r_{\min}$ (a
shell).  Because the cut is on the leaf centers, a sharp sphere or shell
edge can be resolved by the octree (\code{'amr'} with a refined leaf
list) while the density stays exactly uniform --- the case that motivates
uniform density on \code{'amr'} as well as \code{'car'}.

\paragraph{Analytic geometry.}  All geometry queries go through
accessor functions in \code{octree\_mod} (\code{cell\_cx/cy/cz/ch},
\code{leaf\_cell}, \code{leaf\_cx/cy/cz/half}).  On the octree these
are plain array reads; on the car grid they decode the raster
index and evaluate
$c_x = x_{\min} + (i_x + \tfrac12)\,\Delta$ --- the same
floating-point expression that previously filled the arrays, so
nothing changes numerically.  The car build therefore allocates
\emph{none} of the tree arrays: not only \code{children}
($8{\times}$), \code{neighbor} ($6{\times}$), and \code{parent}, but
also the cell centers, half-widths, and the leaf${\leftrightarrow}$cell
maps.  At $512^3$ this removes the entire ${\sim}10$~GB of shared
tree memory; the grid costs only its physical fields.  All physics
modules (thermal, ionization, diffuse, lines, dust) read geometry
exclusively through the accessors, so they run unchanged on either
grid.

\paragraph{Traversal (\code{par\%car\_walk}).}  The car grid is not a
physics choice --- both walks discretize the same transport integral and
agree to rounding --- so \code{par\%car\_walk} selects only \emph{how}
the ray is stepped:
\begin{itemize}
\item \code{'dda'} (default): dedicated incremental Amanatides--Woo
  walks, the natural car traversal.  Each ray precomputes, once, the
  parametric distance to the next face crossing per axis
  ($t_{\max,x/y/z}$) and the crossing period
  ($t_{\Delta,x/y/z} = \Delta/|k_{x/y/z}|$); the inner loop advances
  by comparisons and additions alone --- no exit-face solve at each
  step, no leaf re-lookup.  The tally expressions (analytic
  first-flight $w\,(e^{-\tau_{\rm in}} - e^{-\tau_{\rm out}})/\kappa$,
  pathlength sums, $\tau$-target stop) are identical to the \code{'amr'}
  walks.  It is the fastest mode: $\times 1.64$ over the \code{'amr'} walk,
  $\times 2.0$ over the octree (total Str\"omgren run time
  $13.0 \to 6.5$ minutes).
\item \code{'amr'} (verification): the octree band walks
  (\code{raytrace\_ion\_to\_edge\_amr}, \code{\_to\_tau\_amr},
  \code{\_tau\_only\_amr}) run unchanged; only their two geometry
  primitives branch on the grid mode --- the leaf finder
  \code{amr\_find\_leaf} becomes three integer divisions
  (\code{car\_index}) and the face-neighbor step
  \code{amr\_next\_leaf} becomes pure integer arithmetic on the
  raster index (\code{car\_next}).  The exit face and distance are
  still recomputed each step from the cell center and half-width
  (\code{amr\_cell\_exit}), exactly as on the octree.  Because the
  floating-point order along each ray is then unchanged, the same-seed
  Str\"omgren run is \emph{bit-identical} to the single-level-octree
  run ($\max|\Delta x_{\rm HI}| = 0$) --- the cross-check the DDA walk
  cannot provide (it differs at rounding by design,
  $\max|\Delta x_{\rm HI}| = 2.7\times10^{-15}$, $V_{\rm ion}$
  identical).  This certifies the raster backend rather than serving
  production (it is $\times 1.25$ over the octree, slower than DDA).
\end{itemize}

Everything downstream is untouched (the gas modules only consume
leaf-indexed arrays) and \code{refine\_front} is octree-only.

\subsection{Plane-parallel slab illumination}
\label{sec:slab}
An \emph{xy-periodic slab} models a medium that is horizontally infinite: a
packet that leaves an $x$ or $y$ face re-enters at the opposite face, while the
$z$ faces escape.  The medium may vary only in $z$ (a 1D column,
\code{nx\,=\,ny\,=\,1}) or be an arbitrary 3D tile (\code{nx,ny\,$>$\,1})
repeated periodically.  The wrap is source-agnostic --- it is a boundary
condition on the transport, so an internal source under \code{xy\_periodic} is
a horizontally periodic array of sources, and the external slab illumination
below rides on the same wrap.  In the octree walks the wrap re-enters the
packet at the opposite face's exact boundary coordinate (not by
$x\!\pm\!x_{\rm range}$, which can land a rounding-epsilon outside the box and
have the leaf lookup drop the packet); the incremental DDA wraps the integer
index (\code{modulo}), the parametric $t$-coordinates being
translation-invariant on a uniform grid.

External illumination (\code{source\_geometry\,=\,'slab'}) enters the top
($+z$) and/or bottom ($-z$) face, each as a collimated beam at incidence angle
$\theta$ or an isotropic (Lambert) field.  The physical normalization is the
flux crossing the horizontal face,
\begin{equation}
  F_z \;=\; \int_{\rm hemisphere} I_{\rm in}\cos\theta \, {\rm d}\Omega \;=\;
  \begin{cases} \mu\,F_n, & \text{beam } (\mu=\cos\theta,\ F_n \text{ normal to the beam)} \\
                \pi\,I,   & \text{isotropic } (I \text{ the specific intensity)}, \end{cases}
\end{equation}
so the user inputs the source strength ($F_n$ or $I$) and the geometric factor
($\mu$ or $\pi$) gives $F_z$.  The entering luminosity of one periodic tile is
$L = F_z\,A_{\rm zface}$ (the packet weight is $L/N_{\rm phot}$); $A_{\rm zface}$
cancels in every per-volume result, which therefore scales linearly with $F_z$
and is independent of the tile area and lateral resolution.  The two faces
compose additively (both lit $\Rightarrow L=(F_{z,{\rm top}}+F_{z,{\rm bot}})A$).

The emergent observable is the boundary intensity: a packet escaping a $z$-face
carries its surviving luminosity into a $(\mu=|k_z|,\text{face})$ histogram,
written as $I(\mu)$ [erg\,s$^{-1}$\,cm$^{-2}$\,sr$^{-1}$] with
$F=2\pi\!\int\! I\mu\,{\rm d}\mu$, together with the reflected/transmitted/
absorbed energy budget.  Both the direct beam (the analytic edge walk) and,
with dust scattering, the scattered/diffuse flights contribute, so the budget
closes to unity.  Point-observer peel-off is disabled under \code{xy\_periodic}
(the periodic images admit no finite-distance observer).  The quasi-random
(\code{sobol}) launch covers the slab in both modes (a six-dimensional layout:
frequency, face, the two in-face coordinates, and the Lambert $\mu,\phi$).

\subsection{PDR-zone physics}
\label{sec:pdr}
Three switches (default off) extend the FUV zone beyond ion structure:
\begin{itemize}
\item \code{par\%metal\_ne}: electrons from the metal cascade join the
  $n_e$ closure of the cell solve,
  $n_e = n_{\rm H}[x_{\rm HII} + y_{\rm He}(\ldots)] +
  \sum_{\rm el} A_{\rm el} n_{\rm H} \sum_i (i{-}1) f_i$ --- a
  $\lesssim 10^{-3}$ correction inside the H\,\textsc{ii} region but
  the \emph{only} electrons beyond the front
  ($n_e \simeq n_{\rm CII} \approx A_{\rm C} n_{\rm H}$ in the PDR
  zone), where they control the [C\,\textsc{ii}]/Mg\,\textsc{ii}
  excitation and the photoelectric efficiency.
\item \code{par\%metal\_heat}: metal photoheating enters the thermal
  balance; the heating integrals
  $\mathcal{H}_i = \sum_b [4\pi J\Delta\nu]_b \sigma_i (1 - E_{{\rm th},i}/E_b)$
  are accumulated together with the metal $\Gamma$'s.
\item \code{par\%grain\_pe}: grain photoelectric heating and grain
  recombination cooling with the Bakes \& Tielens (1994) fits,
  $\Gamma_{\rm pe} = 1.3{\times}10^{-24}\,\epsilon\,G_0\,n_{\rm H}$
  with
  $\epsilon = 4.87{\times}10^{-2}/[1 + 4{\times}10^{-3}x^{0.73}]
  + 3.65{\times}10^{-2}(T/10^4)^{0.7}/[1 + 2{\times}10^{-4}x]$,
  $x = G_0\sqrt{T}/n_e$, and
  $\Lambda_{\rm rec} = 4.65{\times}10^{-30}T^{0.94}x^{\beta}n_e n_{\rm H}$,
  $\beta = 0.74/T^{0.068}$, both scaled by \code{pe\_scale} times the
  leaf's dust amount relative to Milky-Way dust, measured from the
  leaf opacity itself:
  $(\rho\kappa/d_{\rm cm})/(n_{\rm H}\,C_{\rm ext,MW}(V))$ with
  $C_{\rm ext,MW}(V) = 4.868{\times}10^{-22}\,{\rm cm^2/H}$ --- robust
  against the DGR-versus-\code{cext\_dust} input conventions (the
  first gate run had DGR $= 1$ with the MW extinction folded into
  \code{cext\_dust}; a naive DGR$/0.01$ scale made the heating
  $100\times$ too strong and drove the zone to a false
  Ly$\alpha$-capped 12~kK equilibrium).  $G_0$ is the local FUV
  field in Habing units summed from the FUV bins of the tally
  (\code{add\_fuv} required).  Three further guards, each found by the
  gate:
  (i) the BT94 fits are evaluated at $\min(T, 10^4\,{\rm K})$ ---
  outside their fit range the $T^{0.7}$ term of $\epsilon$ diverges;
  (ii) both grain terms carry an $x_{\rm HI}$ weight --- the BT94
  charging parameter knows only the FUV field, so in ionized gas
  (EUV-charged grains, grain energy already routed to the IR via
  \code{Heat\_dust}) it would overheat badly; (iii) the package
  includes the \emph{H-impact} fine-structure cooling of
  [C\,\textsc{ii}] 158\,$\mu$m and [O\,\textsc{i}] 63\,$\mu$m
  (two-level, low-density limit;
  $q_{ul}^{\rm H} = 7.6{\times}10^{-10}(T/100)^{0.14}$ and
  $9.2{\times}10^{-11}(T/100)^{0.67}\,{\rm cm^3\,s^{-1}}$, Barinovs et
  al.\ 2005) --- the cooling fits are electron-impact only, and with
  $n_{\rm HI}/n_e \sim 10^3$--$10^4$ in the PDR zone the H collisions
  ARE the coolant.  Without (iii) the zone has no coolant between the
  fine-structure and Ly$\alpha$ regimes, and the
  $\epsilon(x)$--$n_e$ feedback (collisional ionization raises $n_e$,
  neutralizing the grains and raising $\epsilon$ fifty-fold) drives it
  to a false $\sim$10$^4$~K equilibrium.
\end{itemize}
Scope note: this is the atomic PDR skin --- no H$_2$/CO chemistry, no
line transfer, no dust--gas collisional coupling --- so PDR
temperatures are first-order estimates, adequate for the ion structure
and rough [C\,\textsc{ii}] emissivities.

\subsection{$n$-level solver and line luminosities}
\label{sec:nlevel}
\code{nlevel\_mod} evaluates exactly the recipe of
Section~\ref{sec:nlev} (Clenshaw evaluation of the Chebyshev series)
and solves the statistical equilibrium
$\dot n = 0$, $\sum_i n_i = 1$ (partial-pivot elimination;
de-excitation $q_{ul} = 8.629{\times}10^{-6}\,\Upsilon/(g_u\sqrt T)$,
excitation by detailed balance).  \code{lines\_mod} integrates
$L_{\rm line} = \sum_{\rm leaf} j\,n_{\rm ion}V$ on the converged state
and normalizes to the SH95 H$\beta$.  The Fortran chain was verified
against the Python reference solver on an identical state: 26 line
luminosities agree to four significant digits.
With \code{par\%emis\_output} the same solve also writes the
emissivities leaf by leaf (\code{\_emis}: an
\code{emis\_<el>\_<stage>} block for each ion and \code{emis\_h} for
the SH95 lines, [erg\,s$^{-1}$\,cm$^{-3}$] with
$\sum {\rm emis}\cdot V = L_{\rm line}$, verified to $4{\times}10^{-4}$\%),
together with \code{n\_H}, $n_e$, $T_e$, the H/He fractions, and the
metal stage fractions --- a self-contained input for line maps and
position--velocity work.

\subsection{The two-photon continuum: emission, cooling, transport}
\label{sec:twophoton}
The metastable $2s$ ($2^1{\rm S}$ in He\,\textsc{i}) levels decay by
simultaneous emission of two photons sharing the transition energy.  That
one channel enters the code in three separate places, and the three
treatments are deliberately \emph{not} the same; they are collected here
because the differences are easy to mistake for inconsistencies.

\textbf{Emission} (\code{nebcont\_mod}, Section~\ref{sec:contdata}).  The
emission coefficient is written per $n_e n_{\rm ion}$ and is fed by the
\emph{recombination} cascade, one term for each of the three ions:
\begin{center}
\begin{tabular}{llll}
\toprule
ion & weight & $\nu_0$ [Ryd] & $\alpha_{\rm eff}(2s)$ \\
\midrule
H\,\textsc{i}   & $n_e n_{\rm HII}$   & 0.7496 &
  $0.8368{\times}10^{-13}(T_4)^{-0.723}$ (Pengelly 1964) \\
He\,\textsc{i}  & $n_e n_{\rm HeII}$  & 1.514  &
  $\alpha_{\rm eff}(2^1{\rm S}) + [1-p_{584}]\,\alpha_{\rm eff}(2^1{\rm P})$
  (BSS99) \\
He\,\textsc{ii} & $n_e n_{\rm HeIII}$ & 3.00   &
  SH95 table, $6.16 \rightarrow 2.04{\times}10^{-13}$ over 5--30~kK \\
\bottomrule
\end{tabular}
\end{center}
with the Nussbaumer \& Schmutz (1984) shape $A(y)$, $y = \nu/\nu_0$, for
H\,\textsc{i} and (scaled by $Z^6 = 64$ in both $A$ and $A_{2q}$) for
He\,\textsc{ii}, and the Almog \& Netzer (1989) tabulated shape for
He\,\textsc{i}.  The He\,\textsc{i} entry has a history worth recording.  The
port took MOCASSIN's $6.23{\times}10^{-14}(T_4)^{-0.827}$, which is the
recombination coefficient to \emph{all} singlets --- the exact sum
$2.06{\times}10^{-14}+4.17{\times}10^{-14}$ of the $2^1$S and $2^1$P
coefficients of Eq.~(\ref{eq:hei-branch}), refitted as a single power law, so
the exponent $-0.827$ is a fit to the sum and not a sum of exponents.  Using
it is MOCASSIN's stated assumption that every He\,\textsc{i} singlet ends in
$2^1$S, i.e.\ the $p_{584}\rightarrow0$ limit of Eq.~(\ref{eq:hei-584}), and
it disagreed with the $p_{584}\rightarrow1$ limit \code{diffuse\_mod} was
transporting.  That constant is no longer used.  \code{gamma\_2q\_HeI} now
builds the emitting $2^1$S population from the two coefficients separately,
\begin{equation}
 \alpha_{2q}({\rm He\,I}) = 2.06{\times}10^{-14}(T_4)^{-0.451}
   + \left[1-p_{584}\right]4.17{\times}10^{-14}(T_4)^{-0.695},
\end{equation}
with the same local $p_{584}$ that \code{diffuse\_mod} folds into its branch
weights, so the written continuum and the transported field now describe one
nebula (Section~\ref{sec:diffuse}).  The MOCASSIN constant is recovered
exactly as $p_{584}\rightarrow0$ and $\alpha_{\rm eff}(2^1{\rm S})$ alone as
$p_{584}\rightarrow1$.  Only H\,\textsc{i} carries the
collisional suppression
\begin{equation}
  \gamma_{2q} \rightarrow \frac{\gamma_{2q}}
       {1 + n_e q_{2s\text{-}2p}/A_{2q}},
  \qquad A_{2q} = 8.2249\ {\rm s^{-1}},
\end{equation}
with $q_{2s\text{-}2p} = 5.31{\times}10^{-4}$ (at $T \le 10^4$~K) falling to
$4.71{\times}10^{-4}\,{\rm cm^3\,s^{-1}}$ (at $T \ge 2{\times}10^4$~K), so
the continuum is quenched above
$n_{\rm crit} \simeq 1.6{\times}10^4\,{\rm cm^{-3}}$.  No factor is applied
to the helium terms: $A_{2q}$ is $51.3\,{\rm s^{-1}}$ for
He\,\textsc{i} $2^1{\rm S}$ and $526\,{\rm s^{-1}}$ for He\,\textsc{ii}
$2s$, putting their critical densities one to two orders of magnitude
higher.  The three terms are summed into the total $L_\nu$ of
\code{\_nebcont.txt} and also written separately as its last column.
Because the population is the recombination cascade alone, collisional
excitation of $2s$ contributes nothing to the emitted continuum, which is
therefore a lower bound in the partly neutral gas of an ionization front.

\textbf{Cooling} (\code{cooling\_mod}, Section~\ref{sec:cool}).  The
H\,\textsc{i} line-cooling fit is the $n_e \rightarrow 0$
statistical-equilibrium limit of the 25-level CHIANTI atom, and level 2 of
that atom \emph{is} $2s$ ($82258.956\,{\rm cm^{-1}}$, $A = 8.229\,{\rm
s^{-1}}$, the two-photon rate) with the $1s \rightarrow 2s$ collision
strength present.  Collisional excitation of $2s$ therefore cools with its
full 10.2~eV, and the leading term of the fit,
$A_1 = 4.4262{\times}10^{-17}$ at $T_1 = 1.1844{\times}10^{5}$~K, is exactly
that excitation temperature.

The cooling must \emph{not} carry the density quenching the continuum
does, and does not: proton and electron $l$-mixing converts $2s$ to $2p$
and the two-photon continuum into Ly$\alpha$, but the gas loses the same
10.2~eV on either route.  (CHIANTI \code{h\_1} tabulates no
$2s \leftrightarrow 2p$ collision strength, so the equilibrium solve leaves
$2s$ unquenched --- harmless for the same reason.  Nor does the level
produce a spurious line: \code{lines\_mod} extracts only the
H$\alpha$/H$\beta$/H$\gamma$ windows of the collisional component, so
nothing is double counted against the continuum.)  Helium has no
collisional-excitation line cooling at all --- \code{cooling\_total} carries
recombination, free-free, collisional ionization and the H\,\textsc{i} line
term only --- so collisional two-photon cooling of He\,\textsc{i}
$2^1{\rm S}$ and He\,\textsc{ii} $2s$ is absent, as it is in \textsc{Mocassin},
\textsc{Torus}, Wood et al., and \textsc{CMacIonize}.

\textbf{Transport} (\code{diffuse\_mod}, Section~\ref{sec:diffuse}).  Only
helium two-photon photons can ionize anything: the H\,\textsc{i} continuum
stops at 10.2~eV, so omitting it costs nothing.  The He\,\textsc{i}
$2^1{\rm S}$ branch is one of the four diffuse channels and carries its
H-ionizing fraction 0.55644 from the Drake table (Section~\ref{sec:diffuse};
the flat surrogate's 0.56 is no longer used).  The He\,\textsc{ii}
continuum, which extends to
40.8~eV and can ionize both H and He\,\textsc{i}, is \emph{not} transported,
and neither is He\,\textsc{ii} Ly$\alpha$ --- of no consequence wherever
He$^{2+}$ is absent (both Lexington benchmarks), an omission to close for
the hard planetary-nebula and AGN spectra of Section~\ref{sec:limits}.

%==========================================================================
\section{Validation gates and results}
\label{sec:validation}
\subsection{Rate integrals vs.\ analytic attenuation}
Uniform sphere ($n_{\rm H} = 1$, $R = 1$~pc, $\xHI = 0.1$,
$y_{\rm He} = 0.1$), central $10^5$~K source, 16 bins, $10^8$ packets.
Against the cell-volume-averaged analytic attenuation with identical
bins: $\Gamma_{\rm HI}$ mean bias $+0.004\%$, median $0.38\%$
(pure Monte Carlo noise), max $2.6\%$; $\Gamma_{\rm HeI}$ bias
$+0.011\%$ (Figure~\ref{fig:g0}).

\begin{figure}[p]
\centering
\includegraphics[width=0.8\textwidth]{figures/g0_gamma_check.png}
\caption{Rate-integral gate: $\Gamma(r)$ from the Monte Carlo tally (points)
against the analytic point-source attenuation (lines); the lower panel
shows the ratio (shaded regions excluded: cell-size effects at the
center, the voxelized sphere edge outside).}
\label{fig:g0}
\end{figure}

\subsection{Str\"omgren sphere}
Uniform $n_{\rm H} = 100$ sphere, $Q_{\rm H} = 10^{49}\,{\rm s^{-1}}$
($T_\star = 4{\times}10^4$~K, 32 bins), case B, $T_e = 10^4$~K.  The
reference is an exact 1D radial marching with identical bins and rates
--- necessarily \emph{Gauss--Seidel} (each shell solved with the
already-updated interior opacity): a Jacobi sweep advances the front
only one absorption length per pass and never converges.  Sharp-edge
analytic $R_S = 3.153$~pc; the 1D reference gives
$R_{\rm eff} = (3V_{\rm ion}/4\pi)^{1/3} = 3.0365$~pc (the $\sim$4\%
gap is physical: He consumes photons and the model re-emits none).

\begin{center}
\begin{tabular}{lrrr}
\toprule
run & leaves & $R_{\rm eff}$ [pc] & vs.\ 1D \\
\midrule
uniform level 6            & 262\,144 & 3.0534 & $+0.56\%$ \\
front-refined levels 4--6  & 74\,320  & 3.0534 & $+0.56\%$ \\
center-refined levels 4--6 & 7\,568   & 3.1078 & $+2.35\%$ \\
\bottomrule
\end{tabular}
\end{center}
The front-refined octree reproduces the uniform grid to four decimals
with 3.5$\times$ fewer cells; the center-refined control (coarsest
cells \emph{at} the front) shows why the I-front is the one place
needing resolution (Figure~\ref{fig:g1}).

\begin{figure}[p]
\centering
\includegraphics[width=\textwidth]{figures/g1_stromgren_check.png}
\caption{Str\"omgren gate: H and He ionization structure against the 1D exact
reference (black); uniform level-6 (blue), center-refined (orange),
front-refined (green, overlying blue).}
\label{fig:g1}
\end{figure}

\subsection{Thermal H/He nebula}
Same configuration with the thermal balance on
(Figure~\ref{fig:g2a}): $T_e(r)$ matches the 1D thermal reference to
median $0.136\%$, mean bias $+0.006\%$, max $0.95\%$ over
$0.3 < r < 2.8$~pc --- including the central He\,\textsc{iii}-zone
temperature structure and the front drop.  Without metals the nebula
runs at 19--24~kK, the known pure-H/He overshoot that motivates the
metal frame.

\begin{figure}[p]
\centering
\includegraphics[width=\textwidth]{figures/g2a_thermal_check.png}
\caption{Thermal gate: $T_e(r)$ and $\xHI(r)$ with the thermal balance,
against the 1D exact reference.}
\label{fig:g2a}
\end{figure}

\subsection{The MOCASSIN Lexington benchmarks}
\label{sec:hii}
HII20/HII40: blackbody 20/40~kK, $Q_{\rm H} = 1.006/4.26 \times
10^{49}\,{\rm s^{-1}}$, $n_{\rm H} = 100$, hollow shell
$R_{\rm in} = 3{\times}10^{18}$~cm, abundances He 0.1, C $2.2{\times}
10^{-4}$, N $4{\times}10^{-5}$, O $3.3{\times}10^{-4}$, Ne
$5{\times}10^{-5}$, S $9{\times}10^{-6}$.  The stored MOCASSIN 2.02.73.2 baselines are
\emph{reduced} runs ($10^5$ photons, 3 iterations) --- at iteration 3
only 27\% of cells have converged.  To close the comparison MOCASSIN
itself was re-run converged ($2{\times}10^6$ photons, 40 iterations,
72\% cells; a finer $\mathrm{nx}{=}25$ grid confirms resolution), and
MoCHII with the \emph{corrected} charge exchange of
Section~\ref{sec:cx}.

\textbf{Temperatures.}  $T_e$ weighted by $n_{\rm HII} n_e$, MoCHII with
its current defaults (threshold-aligned grid, corrected charge exchange
of Section~\ref{sec:cx}, and the density-suppressed cooling of
Section~\ref{sec:cool}); the benchmark inputs carry the ionizing-band
luminosity that reproduces the specified $Q_{\rm H}$, transport the
diffuse field, and use case-B recombination lines:

\begin{center}\small
\begin{tabular}{lccccc}
\toprule
 & MoCHII & MoCHII & \textsc{Cloudy} & Lexington-2000 & MOCASSIN \\
 & (\code{local\_ne}) & (\code{low\_density}) & c25.00 & (Ercolano+03) & 40-it \\
\midrule
HII20 & 6925 K & 6376 K & 6910 K & 6402--6749 K & --- \\
HII40 & 8169 K & 7894 K & 8210 K & 7720--8199 K & 8176 K \\
\bottomrule
\end{tabular}
\end{center}
With the density-suppressed cooling default (\code{par\%cooling\_model =
local\_ne}) MoCHII gives $T_e = 8169$~K (HII40) and 6925~K (HII20),
matching a current \textsc{Cloudy} c25.00 run (8210 and 6910~K) to
$0.50\%$ and $0.22\%$.  Both modern codes sit a few hundred K \emph{above} the
2000-era published ranges of Ercolano et al.\ (2003) ---
$\langle T[N_pN_e]\rangle = 7720$--$8199$~K (median 8026) at HII40 and
$6402$--$6749$~K (median 6662) at HII20 --- whose tabulated columns carry
older atomic data; the two modern codes now agree with each other far
more closely than either does with the historical spread.

The $\sim$9\% HII40 excess reported in earlier revisions (a superseded
$T_e = 8913$~K) is \emph{resolved}: it was two defects in the Tier-1
cooling fits (Section~\ref{sec:cool}), compounded by the then-missing
density suppression.  \textbf{(D1)} The ground-term fine-structure
populations were distributed by Boltzmann factors --- the
\emph{high}-density limit --- instead of the $n_e\rightarrow0$ limit,
which suppressed the far-infrared fine-structure cooling
([O\,\textsc{iii}] 52/88\,$\mu$m, [N\,\textsc{iii}] 57\,$\mu$m,
[C\,\textsc{ii}] 158\,$\mu$m) by factors $1.6$--$2.9$.
\textbf{(D2)} The \code{.scups} Burgess--Tully \emph{scaling} energy was
used as the physical transition energy in both the Boltzmann factor and
the radiated photon, in place of the observed \code{.elvlc} level
difference (e.g.\ [O\,\textsc{ii}] cooling low by $0.58\times$ at 8~kK).
Both were located by a fixed-state cooling-budget comparison against
\textsc{Cloudy} c25.00 (its \code{save cooling each} output): D1 was
23.9\% and D2 10.6\% of the HII40 cooling budget, while missing channels
([S\,\textsc{iv}] 10.5\,$\mu$m at 1.2\%, now added; all others $<0.2\%$)
and trace-species coupling ($n_e$ plus metal photoheating, together
$+0.2\%$) were excluded.  Regenerating the fits from the
$n_e\rightarrow0$ statistical-equilibrium limit with observed energies
gives HII40 7894~K and HII20 6376~K (the \code{low\_density} option);
restoring the density suppression at the local $n_e$ (the
\code{local\_ne} default) raises them to 8211 and 7025~K, and the
subsequent corrections to the diffuse field --- the Milne ground continua,
the He\,\textsc{i} cascade and the $\alpha_1$ correction of
\S\ref{sec:diffuse} --- settle them at 8169 and 6925~K.  The former
``cooling-side residual'' was thus the low-density over-cooling of the
low-$n_{\rm crit}$ fine-structure lines, and there is no open temperature
discrepancy with the current reference code.

\textbf{Ionization structure} ($n_e$-weighted means, HII40).  The He and
metal ionization is set correctly once the ionizing grid resolves the
He\,\textsc{i} edge.  With \code{par\%ion\_align\_edges} on (the default)
or $\mathrm{nnu\_ion} \geq 128$ the He-ionizing photon budget returns to
the exact blackbody value ($Q_{\rm HeI}/Q_{\rm HI} = 0.108$ against 0.127
for a 32-bin log grid), He\,\textsc{i} rises from 0.058 to 0.163
(MOCASSIN 0.207, closing most of the gap), and the metal 3+ fractions
fall correspondingly.  A separate in-code swap of MOCASSIN-style
recombination (Pequignot 1991 RR + Nussbaumer--Storey 1983 DR) showed the
recombination-coefficient vintage to be a minor lever: MoCHII's Badnell
2023 $\alpha(\mathrm{X\,\textsc{iii}}\to\mathrm{X\,\textsc{ii}})$ runs
$\sim$0.8$\times$ MOCASSIN's, but the swap moved the fractions by only
$\sim$0.026 and closed $\sim$17\% of the ionization gap.  The table below
is the old 32-bin log grid, the reference point for the alignment
discussion (corrected CX vs.\ the converged 40-iteration MOCASSIN):

\begin{center}\small
\begin{tabular}{lcc@{\qquad}lcc}
\toprule
ion & MoCHII (log) & MOCASSIN & ion & MoCHII (log) & MOCASSIN\\
\midrule
C\,\textsc{ii}   & 0.188 & 0.355 & O\,\textsc{ii}  & 0.344 & 0.491 \\
C\,\textsc{iii}  & 0.805 & 0.640 & O\,\textsc{iii} & 0.653 & 0.496 \\
N\,\textsc{ii}   & 0.180 & 0.386 & Ne\,\textsc{ii} & 0.524 & 0.587 \\
N\,\textsc{iii}  & 0.818 & 0.600 & Ne\,\textsc{iii}& 0.474 & 0.401 \\
S\,\textsc{ii}   & 0.194 & 0.066 & S\,\textsc{iii} & 0.750 & 0.786 \\
\bottomrule
\end{tabular}
\end{center}
Switching on the aligned-grid default brings the He and metal fractions
substantially closer to MOCASSIN --- He\,\textsc{i} $0.058 \to 0.163$,
O\,\textsc{iii} $0.653 \to 0.601$, C\,\textsc{iii} $0.805 \to 0.694$ ---
so most of the apparent over-ionization in the table is the log-grid
discretization at the He\,\textsc{i} edge, not a physics error.  At HII20
the CX correction is negligible (low excitation): the structure matches
everywhere except Ne\,\textsc{i}/\textsc{ii} (the 21.6-eV threshold at a
20-kK source, a known code-spread hotspot) and the
S\,\textsc{ii}/\textsc{iii} split.

\textbf{Line fluxes} (ratios to H$\beta$).  With the corrected cooling
the optical collisional lines --- formerly systematically high because
$T_e$ ran $\sim$9\% hot --- now fall inside the published spread:
[O\,\textsc{iii}] 5007 drops from $\sim$4.3 to $\sim$2 against the
published 1.64--3.27, and the [N\,\textsc{ii}], [S\,\textsc{ii}] and
[Ne\,\textsc{iii}] ratios move into range.  Absolute
$L({\rm H}\beta) \approx 1.9 \times 10^{37}$ (HII40) and
$4.7 \times 10^{36}\,{\rm erg\,s^{-1}}$ (HII20) sit $\sim$3\% below the
published $2.01$--$2.10 \times 10^{37}$ and $4.83$--$5.04 \times 10^{36}$,
consistent with helium absorption of the $>$24.6-eV photons and the weak
temperature dependence of $\alpha_{\rm B}$; with the corrected cooling the
HII40 front is captured within the prescribed sphere (outermost-shell
$x_{\rm HII} = 0.11$, recombinations 95\% of $Q_{\rm H}$).  The emissivity
chain itself is exact (Section~\ref{sec:nlevel}).

\subsection{Re-refinement gate and the TNG demonstration}
The re-refinement of Section~\ref{sec:refine} on the Str\"omgren
configuration (level-5 start, rebuild to level 7 at iteration 15,
32\,768 $\rightarrow$ 464\,248 leaves):

\begin{center}
\begin{tabular}{lrrr}
\toprule
run & leaves & $R_{\rm eff}$ [pc] & vs.\ 1D \\
\midrule
re-refined 5$\rightarrow$7 & 464\,248   & 3.0432 & $+0.22\%$ \\
native uniform level 7     & 2\,097\,152 & 3.0430 & $+0.21\%$ \\
\bottomrule
\end{tabular}
\end{center}
The two agree to $+0.008\%$ with $4.5\times$ fewer leaves, and both
improve on level 6 ($+0.56\%$) --- resolution-convergent, as the
resolution-convergence requirement expects.  The front profiles overlie
exactly.

The post-processing demonstration runs MoCHII on an Illustris-TNG
60-kpc cutout (the MoCafe converter output, $n_{\rm H}$ and $\xHI$
columns from the simulation) with a $Q_{\rm H} = 10^{53}\,{\rm s^{-1}}$,
40-kK source at the density peak, thermal balance + metals +
re-refinement (Figure~\ref{fig:tng}): dense clumps cast sharp neutral
shadows, the circumgalactic gas sits at $\xHI \sim 10^{-4}$ with
$n_e$-weighted mean $T_e = 7502$~K, and the refinement (32\,768
$\rightarrow$ 617\,135 leaves) tracks the \emph{distributed} ionization
topology of a turbulent medium --- clump skins and partially ionized
gas, not a single thin shell.  The run also demonstrates why the
max-cell convergence test stalls in such media (clump-surface jitter);
a volume-integrated criterion is queued.

\begin{figure}[p]
\centering
\includegraphics[width=0.98\textwidth]{figures/g4_tng_maps.png}
\caption{MoCHII on an Illustris-TNG cutout: gas density, converged
neutral fraction (note the clump shadows), electron temperature, and
the re-refined cell size, in the source plane.}
\label{fig:tng}
\end{figure}

\subsection{Dusty Str\"omgren gate}
With the MW D03 dust of Section~\ref{sec:dust} added to the Str\"omgren
configuration ($\tau_{\rm d,abs}(13.6\,{\rm eV}) = 1.7$ over the
dust-free radius), against the 1D reference carrying the identical
dust absorption in every bin:

\begin{center}
\begin{tabular}{lccc}
\toprule
case & $R_{\rm eff}$(1D) [pc] & $R_{\rm eff}$(MC) [pc] & dev \\
\midrule
dust-free            & 3.0365 & 3.0534 & $+0.56\%$ \\
full MW dust        & 2.1983 & 2.2058 & $+0.34\%$ \\
\code{laursen09\_live} & 3.0218 & 3.0425 & $+0.68\%$ \\
\bottomrule
\end{tabular}
\end{center}
The classic dusty reduction ($R/R_0 = 0.724$ in 1D, 0.722 between the two
MC runs) is reproduced within the gate tolerance; the live
ionization-dependent dust drives itself to near dust-free
($R/R_0 = 0.996$: interior dust vanishes with the computed
$x_{\rm HII}$), confirming the dust--ionization feedback loop converges
self-consistently.  Energy budget: dust absorbs 61\% of the band
luminosity in the full-dust case
($L_{\rm abs} = 1.933\times10^{38}$ of $3.178\times10^{38}$ erg\,s$^{-1}$).

\textbf{Scattering bracket.}  Turning the scattering on must land the
sphere between two analytically defined bounds: treating the dust
extinction entirely as absorption (scattered photons destroyed; 1D
$R_{\rm eff} = 2.059$~pc) and the absorption-only treatment (scattered
photons never interact; 2.198~pc, MC counterpart 2.206~pc).  The
full-scattering run gives $R_{\rm eff} = 2.197$~pc --- inside the
bracket and close to the absorption-only bound, exactly as the strong
forward scattering ($g = 0.67$--$0.99$) predicts: forward-scattered
photons continue nearly radially, and only the small diffused fraction
is lost from the radial budget.  The dust-absorbed luminosity rises by
$0.7\%$ (rescattered photons reabsorbed) --- energy bookkeeping
consistent.

\textbf{Dust temperature and IR.}  With the band widened to 6~eV
(the band luminosity rescaled by the exact blackbody ratio 2.0969 so
$Q_{\rm H}$ is preserved) and scattering on, the equilibrium mixture
temperature runs from $T_{\rm d} = 87$~K at 0.3~pc to 28~K at the
sphere edge --- the canonical HII-region gradient --- and the IR
spectrum integral closes on the absorbed power to
$L_{\rm IR}/L_{\rm abs} = 1.000$.  (This demo predates the two-segment
\code{add\_fuv} grid and rescaled the single-segment band by hand; the
$+1.5\%$ bin-coarsening shift it measured is what the two-segment grid
now removes.)

\subsection{FUV extension gate: PDR-side C\,\textsc{ii} and Mg\,\textsc{ii}}
The \code{add\_fuv} smoke test (dusty Str\"omgren, 32 ionizing + 16
FUV bins, thermal + metals with Mg/Fe, dust scattering on;
\code{tests/d\_dusty/d\_fuv\_opt.in}) verifies the band mechanics and
the physics it was built for.  Mechanics: the two-segment grid reports
$L_{\rm FUV}/L_{\rm ion} = 1.0969$ for the $4{\times}10^4$~K Planck
source at a 6~eV edge --- exactly the blackbody band ratio applied by
hand in the earlier demo --- with $Q_{\rm H}$ preserved by
construction.  Physics: inside the front (2.2~pc with this dust and
cooling) the metals sit at C\,\textsc{iii}/Mg\,\textsc{iii} as before;
\emph{beyond} it, where H is neutral and the ionizing bins are
exhausted, the FUV bins propagate under dust extinction alone and hold
the low-threshold stages fully singly ionized --- radially averaged
C\,\textsc{ii} $= 1.000$ and Mg\,\textsc{ii} $= 1.000$ from the front
out to the sphere edge, the PDR-side structure that the ionizing-only
band cannot produce (there the same zone recombines to
C\,\textsc{i}/Mg\,\textsc{i}).  Two documented limitations in that
zone: electrons from metal ionization are not fed back into $n_e$ (the
trace-species design), and with photoheating $\approx 0$ the thermal
bisection pins $T_e$ at \code{te\_min} --- PDR temperatures require
photoelectric heating, outside the present scope.  The front-zone
cells with vanishing heating also dominate the $\max|\Delta T_e|$
convergence test (another argument for the queued volume-integrated
criterion); the H\,\textsc{ii}-interior structure and $T_e$
(8.1--10.5~kK) converge normally.

\subsection{External field and multiple sources}
\label{sec:extval}
The source-component model of Section~\ref{sec:sources} is gated in the
\code{ext\_field}, \code{multi\_src}, and \code{peel\_ext} directories
under \code{tests/}:
\begin{itemize}\itemsep2pt
\item \textbf{External field.}  In an optically thin box
  ($\tau \sim 7{\times}10^{-3}$) the interior band-integrated mean
  intensity recovers $J_{\rm ext}$ to 0.9981 (\code{'rec'}) and 0.9991
  (\code{'sph'}), uniform and isotropic (axis asymmetry $< 0.1\%$); the
  entering power equals $\pi J A_{\rm surf}$ [Eq.~(\ref{eq:extlum})]
  exactly.
\item \textbf{Superposition.}  Two point sources against the analytic
  sum $\sum_i L_i/(16\pi^2 r_i^2)$: median deviation 0.76\%, mean ratio
  0.9989; a mixed point $+$ external run, median 0.75\%.
  Two-temperature additivity --- a two-source run against the sum of
  the two single-source runs --- agrees to a bin-wise median of 0.12\%
  and to 1.0001 in the band total.
\item \textbf{Spectra.}  A multi-column \code{src\_spectrum\_file}
  tabulating the same Planck curves as its \code{src\_tstar} twin
  reproduces it to 0.010\% in the worst bin; \code{ext\_spectrum}
  against its \code{ext\_tstar} twin to 0.019\% (this check also caught
  and fixed the single-component path ignoring
  \code{ext\_tstar}/\code{ext\_spectrum}).
\item \textbf{External peel.}  The thin-box image-integrated direct
  flux against $F = J A_{\rm proj}/d^2$: $-0.22\%$ (\code{'rec'}); the
  \code{direc0} flux $+0.15\%$ at chord $\tau \sim 2$; the measured
  central-pixel $\tau_b = -\ln(\mathrm{direct}/\mathrm{direc0})$
  reproduces $n_{\rm H}\,\sigma_{\rm HI}(E_b)\,\ell$ bin by bin to four
  digits with the exact hydrogenic cross section (deviating from the
  $E^{-3}$ power law by up to $\sim$10\%, as $\sigma_{\rm exact}$
  does).
\item \textbf{Physical-unit spectra.}  A \code{'per\_ev'} source against its
  \code{'shape'} twin agrees to 0.004\% in the worst bin (the derived
  \code{src\_lum} is exact); the \code{'per\_ang'} and \code{'per\_um'}
  wavelength conversions reproduce it to 0.013\%; rescaling a physical
  file to a set scale matches its shape twin to 0.0000\%.  A physical
  \code{ext\_spectrum} $J_E$ file recovers the interior
  $\langle J\rangle/J = 0.998$, and doubling \code{ext\_intensity}
  doubles the field exactly (2.0000).
\item \textbf{ISRF presets.}  The measured FUV energy density matches the
  analytic value to $u_{\rm FUV,meas}/u_{\rm FUV} = 1.0002$--$1.0003$:
  $8.94$ (Draine), $5.23$ (Habing), $6.01{\times}10^{-14}$~erg\,cm$^{-3}$
  (Mathis), with exactly zero in the ionizing bins.
\item \textbf{Backward compatibility.}  Single-source runs are
  bit-compatible with the pre-change binary (worst difference 1~ULP,
  $9.9{\times}10^{-16}$ relative).
\end{itemize}

\subsection{Plane-parallel slab}
\label{sec:slabval}
The slab of Section~\ref{sec:slab} is gated under \code{tests/slab/} against
the analytic plane-parallel transport, with the \code{xy\_periodic} wrap
exercised throughout:
\begin{itemize}\itemsep2pt
\item \textbf{xy-periodic wrap.}  A tall thin periodic column ionizes fully
  ($\langle x_{\rm HII}\rangle = 0.9999$) where the isolated (non-periodic)
  control reaches only $0.7072$ --- the wrap re-injects the lateral flux; the
  octree (\code{car\_walk\,=\,'amr'} / \code{grid\_type\,=\,'amr'}) and the
  DDA agree to $8.5{\times}10^{-17}$ for \code{nx}$\ge$2, and after the
  exact-face wrap fix (Section~\ref{sec:slab}) also at \code{nx\,=\,1}
  ($1.3{\times}10^{-18}$).  The AMR leaf-list and car density-cube file inputs
  both run correctly under \code{xy\_periodic}.
\item \textbf{Beam attenuation.}  A frozen neutral slab illuminated by a
  collimated beam gives a photoionization-rate profile $\Gamma_{\rm HI}(z)$
  that is a \emph{perfect} single exponential ($R^2 = 1.00000$); at
  $\theta = 60^{\circ}$ ($\mu = 0.5$) the decay slope is exactly
  slope$_0/\mu$ (ratio $1.0000$) --- the $\mathrm{e}^{-\tau/\mu}$ law and the
  $\cos\theta$ convention are correct.  The isotropic mode gives $F_z = \pi I$
  exactly and a shallower monotone profile.
\item \textbf{Emergent $I(\mu)$ and energy budget.}  A normal beam transmits
  $\mathrm{e}^{-\tau_\perp} = 0.8874$ exactly (budget closes,
  $0.88739 + 0.11261 = 1.0$), with $I(\mu)$ a single spike at $\mu = 1$ and
  zero reflection (no scattering); a $\theta = 60^{\circ}$ beam transmits
  $\mathrm{e}^{-\tau/0.5} = 0.7874$ exactly; isotropic illumination gives a
  broad $I(\mu)$ over all bins (mean $\mu = 0.583$, forward-shifted by the
  $\mathrm{e}^{-\tau/\mu}$ weighting).  Two-face illumination is a symmetric
  superposition.
\item \textbf{Photoionization slab.}  A full ionization solve (case~B OTS)
  builds a plane-parallel ionization front at the Str\"omgren depth
  $d = \Phi/(n_{\rm H}^2\alpha_B)$; the front position matches to within
  $\sim$5--7\% (recombination-coefficient and finite-front-width uncertainty)
  and scales with the incident flux ($d \propto F_z$, ratio $1.94$ for a
  factor-of-two change, the small deficit being the roughly constant front
  width).
\end{itemize}

\subsection{Argon as a data operation}
Adding argon touched no algorithm: two cooling fits (Ar\,\textsc{iii},
Ar\,\textsc{iv}), two $n$-level files, one element file (VFKY96 rows
$Z{=}18$, Voronov, Badnell RR + SVS82 DR2, KF96 CX), one abundance
parameter.  PyNeb cross-check: [Ar\,\textsc{iii}] 7136/7751 within
0.7\%; [Ar\,\textsc{iv}] 4740/4711 within 4.0\% (a density pair --- the
same database-spread class as [O\,\textsc{ii}] and [S\,\textsc{ii}];
PyNeb uses RB97/RGJ19 there).  End-to-end smoke test: [Ar\,\textsc{iii}]
8.99\,$\mu$m, 7136, 7751 appear in the line table with
7136/7751 $= 4.17$, the $A$-ratio, as expected.

\subsection{Magnesium and iron}
Mg and Fe repeat the recipe with three new data wrinkles, all absorbed
by the readers: the Fe\,\textsc{ii} power-law RR (\code{RR2}), the
SVS82-form DR (\code{DR2}), and the Mg/Fe charge-exchange forms and
stage attribution (Section~\ref{sec:cx}).  Products: cooling fits
Mg\,\textsc{ii} (0.05\% max error), Fe\,\textsc{ii} (0.11\%),
Fe\,\textsc{iii} (0.42\%); $n$-level files Mg\,\textsc{ii} (3 levels),
Fe\,\textsc{ii} (13 levels, the IR/optical forbidden lines),
Fe\,\textsc{iii} (14 levels); element files with VFKY96 $Z{=}12, 26$
rows.  PyNeb cross-check: the [Fe\,\textsc{iii}] 4658/4702 ratio agrees
to $-6$\% to $-16$\% --- the largest gate deviation so far, and a known
property of iron: PyNeb's default Fe\,\textsc{iii} collision strengths
(BB14) differ from CHIANTI v11's by tens of percent
(the iron collision-strength spread in the literature), compounded by
the 14-level truncation of a 34-level model.  The number is recorded as
a data-vintage bracket, not a code error --- the same chain reproduces
[O\,\textsc{iii}] to 0.6\%.  For iron-line studies
\code{par\%fe\_levels\_full} loads expanded 16/34-level
Fe\,\textsc{ii}/\textsc{iii} models (\code{nlevel\_fe\_\{2,3\}\_full.txt},
298 [Fe\,\textsc{iii}] lines vs 38, e.g.\ 4658/4702 $+3\%$) in place of
the compact defaults.  End-to-end smoke test (HII20 sphere with
$n_{\rm Mg}/n_{\rm H} = 3{\times}10^{-6}$,
$n_{\rm Fe}/n_{\rm H} = 3{\times}10^{-7}$, depleted gas-phase values):
Mg loads 3 stages, Fe 4 stages; the line table gains the Mg\,\textsc{ii}
2796/2803 doublet (ratio 1.98, the $A$-ratio limit) and the
[Fe\,\textsc{ii}] 1.257\,$\mu$m, 5.34\,$\mu$m, and 26\,$\mu$m lines.

\subsection{Silicon, chlorine, calcium (five stages)}
Si, Cl, and Ca bring the registry to eleven elements, each tracked
through five stages.  The one framework extension was chlorine's cross
section (the VY95 path of Section~3.1 --- Cl is absent from the VFKY96
outer-shell table); Si carries the KF96 charge exchange from the
MOCASSIN \code{chex} rows (recombination direction), Cl and Ca have no
transcription available.  $n$-level files (all $\leq 0.65\%$ worst
transition): Si\,\textsc{ii}/\textsc{iii}/\textsc{iv},
Cl\,\textsc{ii}/\textsc{iii}/\textsc{iv}, Ca\,\textsc{ii},
Ca\,\textsc{v}; Cl\,\textsc{v} and Ca\,\textsc{iii}/\textsc{iv} have
no CHIANTI v11 level data and are tracked without lines (the
Ar\,\textsc{ii} pattern).  Cooling fits: Si\,\textsc{ii} 1.1\%,
Si\,\textsc{iv} 0.2\%, Ca\,\textsc{ii} 0.4\%; Si\,\textsc{iii} 12\%
and the Cl ions 8--9\% resist the multi-exponential form --- accepted,
since at $\sim$$3{\times}10^{-7}$ abundances their share of the
thermal balance is $\sim$$10^{-4}$ and the diagnostics run through the
exact $n$-level solve.  PyNeb gates: [Cl\,\textsc{ii}] 8579/9124 to
1.5\%; Si\,\textsc{iii}] 1892/1883 to 6.9\% and [Cl\,\textsc{iii}]
5518/5538 to 6.6\% (PyNeb's DK94/BZ89 compilations against CHIANTI
v11 --- the established database-spread class).  End-to-end smoke:
[Si\,\textsc{ii}] 34.8\,$\mu$m enters the line table at 4.4\% of
H$\beta$ (the strongest metal line after [C\,\textsc{ii}] ---
the ionization-front/PDR coolant it is), with
Si\,\textsc{iii}] 1883/1892, Si\,\textsc{iv} 1394/1403,
[Cl\,\textsc{ii}] 8581/9126, and the [Cl\,\textsc{iii}] pair at
$5518/5538 = 1.402$, the low-density limit.

\subsection{Line and continuum outputs}
The converged-state outputs now include: the SH95 H\,\textsc{i}
recombination lines (smoke-state check: H$\alpha$/H$\beta = 2.94$ at
$\sim$7~kK); the $n$-level collisional lines for the 24 registry ions
(including [C\,\textsc{ii}] 158\,$\mu$m, C\,\textsc{iii}] 1909,
[N\,\textsc{iii}] 57\,$\mu$m, Mg\,\textsc{ii} 2796/2803, and the
[Fe\,\textsc{ii}] IR lines); and the grid-integrated nebular
continuum $L_\nu$ (free-bound + free-free from the verified tables and
the Milne extension above the edges, plus the three two-photon
continua), written on a 400-point log-$\nu$ grid from 11.4\,$\mu$m to
217\,\AA.

\subsection{He\,\textsc{i} 2$^3$S metastable and 10833~\AA\ absorption}
The 2$^3$S diagnostic (Section~\ref{sec:heimeta}) is gated in
\code{tests/hei\_meta/} on three runs (low/high density and an
\code{add\_fuv} run).  \textbf{Stage 1}: the data-file analytic limits
reproduce the low-density ratio $\alpha_3 n_e/A_{31} =
1.6509{\times}10^{-7}$ (spec $1.65{\times}10^{-7}$ at $n_e=100$, $10^4$~K),
the plateau $6.513{\times}10^{-6}$, and $n_{\rm crit} = 3945$~cm$^{-3}$ (the
$n_e$-scan half-plateau crossing lands at 3926); the Fortran module $n_3$
matches the independent Python closed form on each run's own state to worst
relative $4{\times}10^{-15}$ (machine precision); the low-density core sits
within median 2.8\%\,/\,max 3.9\% of the radiative limit and the
high-density run ($n_e = 1.1{\times}10^4 > n_{\rm crit}$) at 0.676 of the
plateau; $\Phi_3$ from the band tally matches an independent $J_\nu$
integration to $1.5{\times}10^{-9}$ (max $9.07{\times}10^{-5}$~s$^{-1}$); and
$\alpha_3(10^4\,{\rm K}) = 0.998\times0.75\,\alpha_B({\rm He\,II})$, the
diffuse \code{HEI\_F3S} consistency.  \textbf{Stage 2 (10833 absorption)}:
the LINE10833 constants carry the $f$ ratio 1:3:5 and $A =
1.0216{\times}10^7$~s$^{-1}$ each (NIST/p-winds, not the CHIANTI
$1.15{\times}10^7$); a synthetic uniform cube reproduces the column
conservation $\sum N A_{\rm pix} = \sum n V$ and the interior column
$n_3\,\ell$ exactly, and the line-center $\tau_0 = \int\!\kappa_0\,{\rm d}l$
per component and for the doublet matches the hand integral
$\kappa_0^{\rm hand}\,\ell$ to $<2{\times}10^{-16}$; the optically-thin case
raises no warning while a $\times10^5$ scaled case pushes $\tau_0 > 1$ and
fires it.  On the high-density run the column conserves to
$9{\times}10^{-15}$ and the doublet/$J'{=}0$ $\tau_0$ ratio is 8.003, equal
to the $(f_1{+}f_2)/f_0 = 8.002$ oscillator-strength ratio of the blended
lines (that run's metastable column is optically thick, $\tau_0 \sim
10^4$ --- flagged by the warning, as designed).

%==========================================================================
\section{Known limitations and next steps}
\label{sec:limits}
\begin{itemize}
\item By default He recombination emits only the three ground continua;
  the He\,\textsc{i} excited-channel ionizing photons ($\sim$60\% of He
  recombinations) are re-emitted only when \code{par\%hei\_diffuse} is
  set --- their effect on the He\,\textsc{i}/\textsc{ii} split against
  MOCASSIN is small (the split is set by the $>$24.6-eV budget).
  The He\,\textsc{ii} excited cascade has no such channel: its
  Ly$\alpha$ (40.8~eV) and two-photon continuum can ionize H and
  He\,\textsc{i} but are not transported (Section~\ref{sec:twophoton}).
  Immaterial wherever He$^{2+}$ is absent, as in both Lexington
  benchmarks; it belongs to the hard-spectrum track.
\item The EUV part of the leaf field, below the SEDust optics grid, is
  folded into the shortest SEDust bin (energy-conserving; spike
  hardness slightly underestimated).
\item The thermal-loop line cooling is suppressed to the local $n_e$ by
  default (\code{cooling\_model = local\_ne}) through the $(T,n_e)$ table
  of Section~\ref{sec:cool}, and the H-impact [C\,\textsc{ii}]
  158\,$\mu$m / [O\,\textsc{i}] 63\,$\mu$m term (\code{metal\_cooling\_H},
  weighted by $x_{\rm HI}$) carries its own two-level saturation.  What
  remains is the \emph{cold} limit: the cooling fits hold over
  $10^3$--$10^5$~K and collapse below that floor (their lowest excitation
  temperature is 2054~K for C\,\textsc{i}, far above the 23.6~K of
  [C\,\textsc{i}] 609.1\,$\mu$m), and no H-collider term exists for the
  [C\,\textsc{i}] fine-structure lines --- so gas colder than
  $\sim$$10^3$~K, where hydrogen is the collider, is outside the cooling
  model.  Ar\,\textsc{ii} carries no cooling or lines (no CHIANTI v11
  level data).
\item PDR zones: with the Section~\ref{sec:pdr} switches off (the
  default), the FUV extension gives the PDR-side \emph{ion structure}
  only; with them on, $n_e$, metal photoheating, and photoelectric
  heating make the PDR temperatures first-order estimates ---
  H$_2$/CO chemistry and line-transfer cooling remain out of scope.
\item The cell-max convergence test stalls at the front-cell Monte
  Carlo noise floor (and at clump-surface jitter in turbulent media);
  the volume-integrated criterion (\code{conv\_crit = 'vol'}) resolves
  this and is the recommended production setting.
\item The MOCASSIN Lexington benchmark temperatures are resolved
  (Section~\ref{sec:hii}): with the density-suppressed cooling default
  MoCHII (HII40 8169~K, HII20 6925~K) matches a current \textsc{Cloudy}
  c25.00 run to $0.50\%$ and $0.22\%$, the earlier excess having been traced
  to two Tier-1 cooling defects (D1/D2) and the missing density suppression,
  all fixed.  The one remaining note is that both modern codes sit a few
  hundred K above the 2000-era published Lexington temperatures --- a
  data-vintage difference, not a code defect.
\item The HII40 He\,\textsc{i} front sits $0.178$~pc inside the
  \textsc{Cloudy} c25.00 one and that difference is unexplained
  (Section~\ref{sec:diffuse}).  Every single-parameter candidate is now
  excluded by measurement, including the two that were still open --- the
  helium ionization network (Section~\ref{sec:henetwork}) and the metal
  opacity (Section~\ref{sec:metalop}).  The untested hypothesis is a
  combination of several few-percent differences compounding through the
  He$^0$ optical depth, which needs a joint rather than a one-at-a-time
  experiment.
\item The registry-wide recombination audit (Section~\ref{sec:recaudit})
  found no defect --- no transcription error, and no form a better source
  could replace --- but it leaves one question open: the rates that form
  \emph{neutral} atoms differ between Badnell and Verner \& Ferland (1996)
  by $10$--$16\%$, and for the metals no third determination is available
  to say which is right.  Adjudicating it needs an independent
  determination for each ion (R-matrix C\,\textsc{i}/O\,\textsc{i}/%
  Ne\,\textsc{i} recombination, or CHIANTI totals not derived from
  Badnell).  For He\,\textsc{i}, where such a determination does exist,
  MoCHII is within $0.3\%$ of it.
\item Further metal ions as science needs them.
\end{itemize}

\end{document}
